\documentclass[11pt]{article}
\usepackage[a4paper,margin=27mm]{geometry}
\usepackage[T1]{fontenc}
\usepackage[utf8]{inputenc}
\usepackage{lmodern}
\usepackage{microtype}
\usepackage{amsmath,amssymb,amsthm,amscd,mathtools,mathrsfs}
\usepackage{booktabs,array}
\usepackage{enumitem}
\usepackage{natbib}
\usepackage{xurl}
\usepackage{xcolor}
\usepackage[colorlinks=true,linkcolor=blue!45!black,citecolor=blue!45!black,urlcolor=blue!45!black]{hyperref}
\usepackage[nameinlink,noabbrev]{cleveref}

\setlength{\parindent}{0pt}
\setlength{\parskip}{0.52em}
\numberwithin{equation}{section}
\newtheorem{theorem}{Theorem}[section]
\newtheorem{proposition}[theorem]{Proposition}
\newtheorem{lemma}[theorem]{Lemma}
\newtheorem{corollary}[theorem]{Corollary}
\theoremstyle{definition}
\newtheorem{definition}[theorem]{Definition}
\theoremstyle{remark}
\newtheorem{remark}[theorem]{Remark}

\newcommand{\R}{\mathbb R}
\newcommand{\C}{\mathbb C}
\newcommand{\Z}{\mathbb Z}
\newcommand{\T}{\mathbb T}
\newcommand{\Spec}{\operatorname{Spec}}
\newcommand{\dd}{\mathrm d}
\newcommand{\id}{\mathrm{id}}
\newcommand{\opnorm}[1]{\left\lVert #1\right\rVert_{\mathrm{op}}}
\newcommand{\Verified}{\textbf{Verified.}\ }

\title{Cusp Spectral Degeneration and Local Toric Kummer Extension\\
\large Conditional results from the reconstructed candidate $S^6$
complex-torus fibration}
\author{A specialist-verification draft\\
\small based on the candidate $S^6$ construction of Levent Alp\"oge}
\date{Revised conditional specialist-verification draft, 4 September 2026}

\begin{document}
\maketitle

\begin{abstract}
Use the period, integral-monodromy, and toric-cusp data from a recent
manuscript proposing a complex structure on $S^6$.  Those inputs have been
reconstructed and separately verified in the accompanying workbench, but
they are imported, rather than reproved, in this standalone note.  The
regular fibres are marked complex two-tori
\[
 F_z=\C^2/\mathsf P_z\Z^4
\]
with an explicit real period matrix $\mathsf P_z$.  We equip the covering space $\R^4\cong\C^2$ with its coordinate Euclidean metric and study the resulting scalar fibre Laplacians.  Every fixed fibre has a positive first eigenvalue.  Nevertheless, on the actual manuscript cusp, where $T=\operatorname{Im}\tau\to\infty$, $\mu$ and $\beta+\tau$ remain bounded, one has the sharp asymptotics
\[
 \lambda_1(F_z)=\frac{4\pi^2}{T^2}
 \left(1+\frac{2\min\{b_{\mathrm I,0},0\}}{T}+O(T^{-2})\right),
 \qquad
 \operatorname{Vol}(F_z)=T^2-b_{\mathrm I,0}T+O(1),
\]
where $b_{\mathrm I,0}=\operatorname{Im}(\beta+\tau)(0)$.  The two possible shortest
signed Fourier pairs, their exact comparison, and their first corrections are
determined below.  After rescaling each four-torus to unit volume, the first
eigenvalue still tends to zero, now as $4\pi^2/T$.  The decomposable vertical
operator obtained by integrating the mean-zero fibre Laplacians along a cusp
ray is injective but has zero at the bottom of its essential spectrum.

The degeneration is not represented by only one soft Fourier mode.  If
$\mathcal D_T:=|\det\mathsf P_z|=-TD$ and $N_T(E)$ counts scalar
eigenvalues with Fourier multiplicity, then throughout the first transverse
band $0<E<4\pi^2$ one has the fixed-energy law
\[
 N_T(E)=\frac{E\mathcal D_T}{4\pi}+O_E(T).
\]
Thus a two-dimensional density of low modes grows on the order of the fibre
covolume.  We also give the unit-volume counting law and the exact finite
count at the collapsing $T^{-2}$ scale.

The native marked four-tori do not converge to $\R^4$: their pointed smooth
limit is $\T^2\times\R^2$.  Replacing the two fixed unit periods by $M$ times
those periods gives degree-$M^2$ local fibre covers.  They converge pointed
  smoothly to $\R^4$, along sequences with $T\to\infty$, exactly when
  $M\to\infty$ for the unscaled metric.  For
the pullback of the base-fibre unit-volume metric the condition is
$M/T^{1/2}\to\infty$; if every cover is instead normalized to have its own
volume one, no choice of $M$ has pointed smooth limit $\R^4$.  After the
cyclic base change $t_c=v^M$ the cover monodromy is integral, and the complete
cover spectrum separates its deck-invariant pullback sector from its new
momentum sectors.  These finite-cover limits are different from both the
native $\T^2\times\R^2$ limit and the quotient constructions below.

The full monodromy has, up to sign, a unique primitive invariant integral character.  It gives a real, non-holomorphic circle quotient of the regular group family.  In the convention above, with $m=\operatorname{Im}\mu$, $D=\operatorname{Im}\beta-6m^2/T<0$, and $b_{\mathrm I,0}=\operatorname{Im}(\beta+\tau)(0)$, its exact quotient circumference and radial asymptotic are
\[
 \ell_\gamma=\frac{-TD}{\sqrt{T^2+m^2}}
 =T-b_{\mathrm I,0}+O(T^{-1}).
\]
Four synchronized copies, with no base change or fibre cover, give a real four-torus quotient over the full regular base whose radial fibres converge pointed locally isometrically to $\R^4$.  Four copies are minimal within the class of fixed primitive full-monodromy-invariant kernels equipped with the induced product quotient metric.  The literal map to this fixed product $\mathbb T^4$ target does not extend continuously through the asserted toric cusp and does not factor through the fourfold smash collapse; these statements concern that target and that regular-fibre restriction only.

The included local toric theorem gives the exact Kummer extension after the
base change $t=v^N$.  If $Z(s)=sB_0+C(t)$, the source action is not the
naive target action with $v$ substituted for $t$: it is
\[
 \Psi^{(N)}_{\bar\lambda}(\xi,v)
 =\left(\exp\!\left(\frac{2\pi i}{N}C(v^N)\bar\lambda\right)
 v^{B_0\bar\lambda}\xi,v\right).
\]
With this action, $(\xi_1,\xi_2,v)\mapsto
(\xi_1^N,\xi_2^N,v^N)$ descends to a finite flat map of degree $N^3$ to
the original completed cusp family.  Its factor to the base-changed target
has generic degree $N^2$; for $N>1$ that relative map is quasi-etale but not
flat along the preimages of the double and triple central-fibre strata.
This is a completed local cusp morphism, not a cover over the full base.

For the classical Yang--Mills functional at the product connection, Coulomb
reduction gives an exact Hessian gap
\[
 \gamma_T=\frac{4\pi^2}{g_{\mathrm{YM}}^2T^2}
 \left(1+\frac{2\min\{b_{\mathrm I,0},0\}}{T}+O(T^{-2})\right).
\]
The corresponding finite cubical cochain model has exact gap
$\rho_N/(4\pi^2)$ times this value, where
$\rho_N=4N^2\sin^2(\pi/N)$; it closes as $T^{-2}$ uniformly for every
$N\ge2$, although resolving the varying continuum metric requires
$N/T\to\infty$.  These are linear classical fluctuation operators, not a
nonlinear lattice gauge theory or a quantum Hamiltonian.

An independent compact-link and Gaussian construction keeps the cover index
$M_{\mathrm{cov}}$ separate from the cubical subdivision count
$N_{\mathrm{mesh}}$.  It proves the exact holonomy and plaquette formulas for
an explicit periodic $\mathrm{SU}(2)$ flux and the complete finite cochain
spectrum.  The associated finite-dimensional Gaussian oscillator has a
unique ground state and gap $\sqrt{\lambda_{\min}}$.  Along the explicit
sequence
\[
 R_j=j^2,\qquad M_{\mathrm{cov},j}=j,\qquad
 N_{\mathrm{mesh},j}=j^3,
\]
the covering tori converge pointed locally isometrically to $\R^4$, the mesh
resolves the varying metric, and both the Gaussian gap and the displayed
logarithmic flux action tend to zero.  This Gaussian Hilbert space is an
$L^2$-space on a linear Coulomb slice; it is not the compact-link Haar
Hilbert space and does not define an interacting Hamiltonian.

For every fixed nondegenerate integral alternating form $\mathsf E$ on the
marked four-torus, the associated constant-curvature magnetic Laplacian has
oscillator frequencies $b_1,b_2>0$ satisfying
\[
 b_1b_2=\frac{4\pi^2|\operatorname{Pf}(\mathsf E)|}{\mathcal D_T},
 \qquad
 \operatorname{gap}\le
 4\pi\sqrt{\frac{|\operatorname{Pf}(\mathsf E)|}{\mathcal D_T}}
 =O(T^{-1}).
\]
This is a fibrewise statement in the marked cusp family and does not by
itself globalize the flux through monodromy.  The actual lattice nevertheless
has the explicit invariant full-flux class
\[
 \eta=u\wedge w+6\gamma\wedge\delta,
 \qquad \operatorname{Pf}(\eta)=6.
\]
Writing $R$ for logarithmic cusp height, its two charge-one frequencies obey
$\mathcal B_{R,+}\mathcal B_{R,-}=24\pi^2/D_R$ and have leading terms
$12\pi/R$ and $2\pi/R$.  Every joint Landau level has complex multiplicity
six, and the first distinct spacing is $4\pi/R+O(R^{-2})$.  The corresponding
maximal multiplication operators on both the fixed and varying measurable
radial Hilbert fields are self-adjoint and have purely continuous spectrum
$[0,\infty)$, empty point spectrum, and zero kernel.  The radial variable in
these operators is only a direct-integral parameter: there is no radial
derivative or boundary condition.

The split connection on
$\mathscr L_{\eta,R}\oplus\mathscr L_{\eta,R}^{-1}$ is a non-flat reducible
$\mathrm{SU}(2)$ Yang--Mills critical point.  Its closed Coulomb Jacobi
operator has exactly three negative bands on a sufficiently deep cusp tail,
each of real multiplicity $48$, and hence real Morse index $144$.  For every
nonzero trace-zero integral weight vector
$\mathbf m\in\mathbb Z^r$, the split-Cartan $\mathrm{SU}(r)$ construction has
exact rootwise negative bands and deep-tail index
$36r\sum_jm_j^2$.  These are spectra of classical fibre Hessians at explicit
reducible backgrounds, not the spectrum of interacting four-dimensional
quantum Yang--Mills theory.

By contrast, the specific degree-six magnetic construction pulled back from
the holomorphic first-coordinate elliptic quotient requires
$6\mu=a\tau+b$ for integers $a,b$.  We prove, from the stated transformation
and cusp laws, that this exceptional locus has no cusp-accumulating sequence
and is finite on the completed orbifold base.  Consequently that
first-coordinate Landau spectrum cannot supply the cusp limit.  Finally, all
fibre and finite Gaussian degenerations are distinguished from the scalar
Laplacian of any smooth metric on the compact total space, which necessarily
has compact resolvent and a positive first eigenvalue.
\end{abstract}

\paragraph{Status and logical scope.}
All statements involving the global $S^6$ fibration use the unpublished
construction in \citet{AlpogeS6_2026}.  Its topology, singular-fibre maps,
and direct-image argument have been reconstructed and separately verified in
the accompanying proof-verification workbench; this standalone note imports
that geometric package as stated input instead of reproducing those proofs.
A contemporaneous downstream paper treats
that construction as an assumption rather than supplying an independent
proof \citep{LiangLiuZhang_2026}.  The linear algebra, Fourier spectra, cusp
estimates, and direct-integral argument below are independent deductions
from the displayed period and cusp formulae.  The regular-locus quotient
theorem additionally uses the asserted integral monodromy matrices, which
are restated and checked below; its nonextension statement additionally uses
the asserted toric cusp chart, and is only a nonextension result for the
literal fixed product-circle target with the stated regular-fibre map.  The
quotient map is real.  Its radial
circumference and pointed-limit statements use the separately chosen
Euclidean metrics on marked cusp fibres; no deck-invariant global fibre
metric is asserted.  It is neither the holomorphic first-coordinate quotient
nor the separate degree-$M^2$ cover construction.
The counting theorem and the arbitrary-flux magnetic theorem are likewise
statements along the marked cusp.  A general alternating form $\mathsf E$ is
fixed only in the marked lattice; a global magnetic line bundle over the
regular family would require additional descent data.  Separately, the
displayed class $\eta=u\wedge w+6\gamma\wedge\delta$ is proved invariant
under the asserted monodromy generators and is realized by an explicit factor
of automorphy on every marked fibre.  The radial direct-integral theorem uses
Lebesgue measure in logarithmic height and a multiplication operator with the
exact maximal domain; it includes no radial derivative.
The local toric Kummer theorem included below is proved for a completed cusp
disc.  It uses the modified source factor
$\exp(2\pi iC(v^N)\bar\lambda/N)$ and distinguishes the degree-$N^3$ map
to the original cusp family from the generic degree-$N^2$ relative map to
the base-changed target.  It supplies no full-base cover.  The separate cone
targets, simultaneous-root construction, and Seifert filling proofs in the
canonical workbench are not imported or claimed here.  This paper does prove
statements about scalar and magnetic fibre Laplacians, their specified
decomposable radial multiplication operators, classical Yang--Mills Hessians
at the product and explicit reducible split backgrounds, an honest
compact-link flux configuration, and explicitly defined finite cochain and
Gaussian models.  It constructs neither nonlinear quantum Yang--Mills theory
nor a four-dimensional physical Hamiltonian, and makes no physical mass-gap
claim.

\section{Conditional period data}

Let $\mathfrak h_z$ be the marked parameter upper half-plane used in the proposed construction.  The asserted holomorphic period functions are
\[
 \tau:\mathfrak h_z\longrightarrow\mathfrak h,
 \qquad \mu,\beta:\mathfrak h_z\longrightarrow\C,
\]
and the complex period matrix and its real form are
\begin{equation}\label{eq:Z-P}
 Z(z)=\begin{pmatrix}6\mu(z)&\tau(z)\\ \beta(z)&\mu(z)\end{pmatrix},
 \qquad
 \mathsf P_z=
 \begin{pmatrix}
 6\Re\mu&\Re\tau&1&0\\
 6\Im\mu&\Im\tau&0&0\\
 \Re\beta&\Re\mu&0&1\\
 \Im\beta&\Im\mu&0&0
 \end{pmatrix}.
\end{equation}
Write
\begin{equation}\label{eq:D}
 T=\Im\tau>0,
 \qquad m=\Im\mu,
 \qquad p=\Im\beta,
 \qquad D=p-\frac{6m^2}{T}.
\end{equation}
The manuscript asserts that an additive normalization of $\beta$ can be chosen so that $D<0$ everywhere and proves the determinant identity
\begin{equation}\label{eq:detP}
 \det\mathsf P_z=T D.
\end{equation}
Thus $\mathsf P_z\Z^4$ is a full lattice in $\R^4$ at each regular parameter.

At the distinguished cusp, let $t_c$ be a local coordinate with $t_c=0$ at the added point.  The manuscript gives
\begin{equation}\label{eq:cusp-input}
 q=e^{2\pi i\tau}=t_cu_1(t_c),\qquad u_1(0)\ne0,
\end{equation}
while $\mu$ and $b:=\beta+\tau$ extend holomorphically across $t_c=0$.
On a sufficiently small cusp disc choose a holomorphic logarithm of $u_1$ and
put
\begin{equation}\label{eq:cusp-s-coordinate}
 h=\frac{1}{2\pi i}\log u_1,\qquad
 s=\tau-h,\qquad e^{2\pi i s}=t_c.
\end{equation}
The positive radial lift is $s=iR$, where $R\to\infty$.  We reserve
$T=\operatorname{Im}\tau=R+\operatorname{Im}h$ for the convention in
\eqref{eq:D}; the radial height $R$ is therefore not denoted by $T$ in this
paper.  Along that lift put
\begin{equation}\label{eq:radial-b-data}
 B_R:=\operatorname{Im}b(e^{-2\pi R}),\qquad
 b_{\mathrm I,0}:=\operatorname{Im}b(0),\qquad
 m_0:=\operatorname{Im}\mu(0).
\end{equation}
Holomorphy at the cusp gives
\begin{equation}\label{eq:radial-holomorphic-errors}
 B_R=b_{\mathrm I,0}+O(e^{-2\pi R}),\qquad
 m=m_0+O(e^{-2\pi R}),
\end{equation}
and, because $b=\beta+\tau$,
\begin{equation}\label{eq:radial-covolume-exact}
 p=B_R-T,\qquad
 \mathcal D_T:=|\det\mathsf P_z|=-TD
 =T(T-B_R)+6m^2>0.
\end{equation}
In particular, along every path in a fixed cusp sector on which $T\to\infty$,
one has
\begin{equation}\label{eq:cusp-asymptotic-input}
 m=O(1),\qquad p=-T+O(1),\qquad D=-T+O(1).
\end{equation}
The coarse asymptotics suffice for the leading spectral term; the exact
radial identities above determine its first correction.

The generator laws needed later are
\begin{equation}\label{eq:generator-laws}
 \tau\circ g_0=\tau-1,\qquad \mu\circ g_0=\mu,
 \qquad
 \tau\circ g_2=-\frac1\tau,\qquad
 \mu\circ g_2=1+\frac\mu\tau.
\end{equation}

For the regular-locus quotient, we also use the asserted integral lattice
\begin{equation}\label{eq:monodromy-lattice}
 \Lambda=\Z\langle\widehat\gamma,\widehat u,
                    \widehat w,\widehat\delta\rangle
\end{equation}
in the column order of \eqref{eq:Z-P}.  Let
$\Delta=\langle g_1,g_2\mid g_1^3=g_2^4=1\rangle$ denote the manuscript's
triangle-group deck action.  It must not be identified with the ordinary
fundamental group of
$B^\circ=B\setminus\{p_0,p_1,p_2\}$.  Put $g_0=(g_1g_2)^{-1}$.  Its asserted
action on $\Lambda$ is
generated by
\begin{equation}\label{eq:monodromy-matrices}
 A_1=
 \begin{pmatrix}
 1&0&0&0\\6&0&1&0\\-6&-1&-1&0\\-2&1&0&1
 \end{pmatrix},\qquad
 A_2=
 \begin{pmatrix}
 1&0&0&0\\0&0&-1&0\\-6&1&0&0\\3&0&1&1
 \end{pmatrix},
\end{equation}
and the cusp monodromy is
\begin{equation}\label{eq:cusp-monodromy-matrix}
 M_0=(A_1A_2)^{-1}=
 \begin{pmatrix}
 1&0&0&0\\0&1&0&0\\0&1&1&0\\-1&0&0&1
 \end{pmatrix}.
\end{equation}
Let $\mathfrak h_z^\circ$ be the inverse image of $B^\circ$ in the marked
parameter cover.  We use here the manuscript's assertion that the elliptic
fixed-point orbits have been removed, that the deck action on
$\mathfrak h_z^\circ$ is free, and that its lift to the complex torus family
is holomorphic and complex-linear.  The fibre over the marked cover is
$\C^2/\mathsf P_z\Lambda$, and in flat angular coordinates the lifted action
is
\begin{equation}\label{eq:flat-monodromy-action}
 (z,x+\Lambda)\longmapsto(gz,A_gx+\Lambda).
\end{equation}
It is linear and fixes the zero section.  Consequently its quotient
$J=X^\circ=f^{-1}(B^\circ)$ is, conditionally on this input, a holomorphic
family of complex Lie groups rather than merely an affine torus bundle.

\section{The fixed-fibre flat spectrum}

Equip $\R^4$ with the Euclidean metric in the coordinates used in \eqref{eq:Z-P}, and put
\begin{equation}\label{eq:flat-fibre}
 F_z^{\mathrm{flat}}=\R^4/\mathsf P_z\Z^4.
\end{equation}
This metric is a chosen auxiliary metric determined by the displayed coordinate realization.  It is not asserted to be canonical for the complex torus or to extend across a singular fibre of the completed family.

\begin{proposition}[Fourier spectrum on a marked flat fibre]\label{prop:fixed-spectrum}
If $\det\mathsf P_z\ne0$, let $\Delta_z$ be the nonnegative self-adjoint
operator associated with the closed Dirichlet form
\[
 q_z[f]=\int_{F_z^{\mathrm{flat}}}|\dd f|^2\,
 \dd\operatorname{vol}_{F_z},\qquad \operatorname{Dom}(q_z)=H^1(F_z^{\mathrm{flat}}).
\]
Then $\operatorname{Dom}(\Delta_z)=H^2(F_z^{\mathrm{flat}})$ and
$\Delta_z$ has spectrum
\begin{equation}\label{eq:fixed-spectrum}
 \Spec(\Delta_z)=
 \left\{4\pi^2\left\lVert\mathsf P_z^{-\mathsf T}k\right\rVert^2:
 k\in\Z^4\right\},
\end{equation}
with multiplicities.  Its kernel is the constants and
\begin{equation}\label{eq:fixed-gap}
 \lambda_1(z)=4\pi^2
 \min_{k\in\Z^4\setminus\{0\}}
 \left\lVert\mathsf P_z^{-\mathsf T}k\right\rVert^2>0.
\end{equation}
For an arbitrary $k=(r,s,n_1,n_2)^{\mathsf T}$, set
\[
 \alpha=\Re\mu,\qquad \eta=\Re\tau,\qquad c=\Re\beta,
\]
\begin{equation}\label{eq:shifted-Fourier-indices}
 R_z(k)=r-6\alpha n_1-cn_2,
 \qquad S_z(k)=s-\eta n_1-\alpha n_2.
\end{equation}
The complete dual-vector formula is
\begin{equation}\label{eq:complete-dual-vector}
 \mathsf P_z^{-\mathsf T}k
 =\left(n_1,
 \frac{mR_z(k)-pS_z(k)}{-TD},n_2,
 \frac{-TR_z(k)+6mS_z(k)}{-TD}\right).
\end{equation}
\end{proposition}

\begin{proof}
The characters
\[
 e_k([x])=\exp\bigl(2\pi i\langle\mathsf P_z^{-\mathsf T}k,x\rangle\bigr),
 \qquad k\in\Z^4,
\]
form the Fourier basis for the compact abelian group $F_z^{\mathrm{flat}}$.
Applying $-\sum_j\partial_{x_j}^2$ gives the eigenvalue in
\eqref{eq:fixed-spectrum}.  The standard closed-form construction and
elliptic regularity give the stated form and operator domains.  The dual
lattice is discrete, so its shortest nonzero vector has positive length.
Standard flat-torus spectral theory gives completeness \citep{Chavel_1984}.
Finally, solving $\mathsf P_z^{\mathsf T}x=k$ gives $x_1=n_1$ and
$x_3=n_2$; the remaining system is
\[
 \begin{pmatrix}6m&p\\T&m\end{pmatrix}
 \binom{x_2}{x_4}=\binom{R_z(k)}{S_z(k)}.
\]
Its determinant is $-TD$, and direct inversion proves
\eqref{eq:complete-dual-vector}.
\end{proof}

\section{Sharp degeneration on the actual cusp}

The following result uses the actual cusp behavior \eqref{eq:cusp-asymptotic-input}; it is not based on an arbitrary family in the ambient period-matrix space.

\begin{theorem}[Shortest cusp modes and their first corrections]
\label{thm:cusp-gap}
Assume the conditional period data in Section~1 and restrict to the positive
radial lift $s=iR$.  Put
\[
 e_\gamma=(1,0,0,0)^{\mathsf T},\qquad
 e_u=(0,1,0,0)^{\mathsf T},\qquad
 \ell_{\gamma,T}=\mathsf P_z^{-\mathsf T}e_\gamma,\qquad
 \ell_{u,T}=\mathsf P_z^{-\mathsf T}e_u.
\]
For all sufficiently large $T=\operatorname{Im}\tau$, every shortest
nonzero vector of the dual lattice belongs to
\[
 \{\pm\ell_{\gamma,T},\ \pm\ell_{u,T}\}.
\]
Their exact squared lengths and exact comparison are
\begin{align}
 \|\ell_{\gamma,T}\|^2
 &=\frac{T^2+m^2}{\mathcal D_T^2},&
 \|\ell_{u,T}\|^2
 &=\frac{(T-B_R)^2+36m^2}{\mathcal D_T^2},
 \label{eq:unit-mode-exact-lengths}\\
 \mathcal D_T^2
 \bigl(\|\ell_{\gamma,T}\|^2-\|\ell_{u,T}\|^2\bigr)
 &=B_R(2T-B_R)-35m^2.
 \label{eq:unit-mode-comparison}
\end{align}
Thus the shortest indices are $\pm e_u$ when the right-hand side of
\eqref{eq:unit-mode-comparison} is positive, $\pm e_\gamma$ when it is
negative, and both signed pairs when it vanishes.

The two candidate eigenvalues have the expansions
\begin{align}
 4\pi^2\|\ell_{\gamma,T}\|^2
 &=\frac{4\pi^2}{T^2}
   \left(1+\frac{2b_{\mathrm I,0}}{T}+O(T^{-2})\right),
 \label{eq:gamma-eigenvalue-correction}\\
 4\pi^2\|\ell_{u,T}\|^2
 &=\frac{4\pi^2}{T^2}
   \left(1+O(T^{-2})\right).
 \label{eq:u-eigenvalue-correction}
\end{align}
More generally, for each fixed $(r,s)\in\Z^2$,
\begin{equation}\label{eq:general-plane-mode-correction}
 \frac{T^2}{4\pi^2}\lambda_T(r,s,0,0)
 =r^2+s^2+\frac{2b_{\mathrm I,0}r^2-10m_0rs}{T}+O_{r,s}(T^{-2}).
\end{equation}
Consequently
\begin{equation}\label{eq:cusp-gap-result}
 \boxed{\lambda_1(z)=\frac{4\pi^2}{T^2}
 \left(1+\frac{2\min\{b_{\mathrm I,0},0\}}{T}+O(T^{-2})\right).}
\end{equation}
If $b_{\mathrm I,0}>0$, the bottom indices are eventually $\pm e_u$; if
$b_{\mathrm I,0}<0$, they are eventually $\pm e_\gamma$.  If
$b_{\mathrm I,0}=0$, the exact comparison
\eqref{eq:unit-mode-comparison} decides; in particular, if $m_0\ne0$ then
the bottom indices are eventually $\pm e_\gamma$.  The first positive
eigenvalue has complex multiplicity two when exactly one signed pair
minimizes and multiplicity four when both pairs minimize.  All statements
are independent of the real parts of $\tau,\mu,\beta$.
\end{theorem}

\begin{proof}
Write $k=(r,s,n_1,n_2)\in\Z^4$ and solve
$\mathsf P_z^{\mathsf T}x=k$.  The third and fourth equations are
\begin{equation}\label{eq:x1x3}
 x_1=n_1,\qquad x_3=n_2.
\end{equation}
Hence every mode with $(n_1,n_2)\ne(0,0)$ has dual squared norm at least
one.  For $k=(r,s,0,0)$, the remaining system is
\[
 \begin{pmatrix}6m&p\\T&m\end{pmatrix}
 \begin{pmatrix}x_2\\x_4\end{pmatrix}
 =\begin{pmatrix}r\\s\end{pmatrix}.
\]
Its determinant is $-TD=\mathcal D_T$, and $p=B_R-T$, so
\begin{equation}\label{eq:dual-exact}
 \mathsf P_z^{-\mathsf T}(r,s,0,0)
 =\left(
 0,\frac{mr+(T-B_R)s}{\mathcal D_T},
 0,\frac{-Tr+6ms}{\mathcal D_T}
 \right).
\end{equation}
Define
\begin{equation}\label{eq:MT}
 M_T:=\frac{T}{\mathcal D_T}
 \begin{pmatrix}m&T-B_R\\-T&6m\end{pmatrix}.
\end{equation}
Equations \eqref{eq:radial-holomorphic-errors} and
\eqref{eq:radial-covolume-exact} give
\begin{equation}\label{eq:MT-J}
 M_T=J+\frac1T
 \begin{pmatrix}m_0&0\\-b_{\mathrm I,0}&6m_0\end{pmatrix}
 +O_{\mathrm{op}}(T^{-2}),
 \qquad J=\begin{pmatrix}0&1\\-1&0\end{pmatrix}.
\end{equation}
In particular, because $J$ is orthogonal, a constant $C$ exists such that
for every $v\in\R^2$ and every sufficiently large $T$,
\begin{equation}\label{eq:singular-value-bound}
 \left(1-\frac CT\right)\|v\|
 \le\|M_Tv\|
 \le\left(1+\frac CT\right)\|v\|.
\end{equation}
Thus every plane vector with $r^2+s^2\ge2$ is eventually longer than all
four plane vectors with $r^2+s^2=1$, while the vectors excluded by
\eqref{eq:x1x3} have norm at least one and the four candidates have norm
$T^{-1}(1+O(T^{-1}))$.  This proves the asserted exhaustive minimizing list.

Putting $(r,s)=(1,0)$ and $(0,1)$ in \eqref{eq:dual-exact} proves
\eqref{eq:unit-mode-exact-lengths}.  Subtracting their numerators gives
\[
 T^2+m^2-(T-B_R)^2-36m^2
 =B_R(2T-B_R)-35m^2,
\]
which proves \eqref{eq:unit-mode-comparison}.  Finally, squaring
\eqref{eq:MT-J} on a fixed vector $(r,s)$ gives
\[
 \|M_T(r,s)^{\mathsf T}\|^2
 =r^2+s^2+\frac{2b_{\mathrm I,0}r^2-10m_0rs}{T}+O_{r,s}(T^{-2}).
\]
Since $T^2\|\mathsf P_z^{-\mathsf T}(r,s,0,0)\|^2
=\|M_T(r,s)^{\mathsf T}\|^2$, this proves
\eqref{eq:general-plane-mode-correction} and the two special expansions.
The smaller of the relative corrections $2b_{\mathrm I,0}$ and $0$ is
$2\min\{b_{\mathrm I,0},0\}$, proving \eqref{eq:cusp-gap-result}; the exact comparison
and Fourier signed-pair counting prove the remaining assertions.
\end{proof}

\begin{corollary}[Volume and unit-volume degeneration]\label{cor:unit-volume}
Let $V_z$ be the Euclidean volume of $F_z^{\mathrm{flat}}$.  Then
\begin{equation}\label{eq:volume-asymptotic}
 V_z=|\det\mathsf P_z|=\mathcal D_T
 =T(T-B_R)+6m^2
 =T^2-b_{\mathrm I,0}T+6m_0^2+O(Te^{-2\pi R}).
\end{equation}
If $g_z$ is the flat metric and
\begin{equation}\label{eq:unit-volume-metric}
 \widehat g_z=V_z^{-1/2}g_z
\end{equation}
so that $\operatorname{Vol}(F_z,\widehat g_z)=1$, then
\begin{equation}\label{eq:unit-volume-gap}
 \lambda_1(F_z,\widehat g_z)
 =\frac{4\pi^2}{T}
 \left(1+\frac{-b_{\mathrm I,0}/2+2\min\{b_{\mathrm I,0},0\}}{T}
 +O(T^{-2})\right)\longrightarrow0.
\end{equation}
In the normalized metric the two fixed periods have length
$\mathcal D_T^{-1/4}=T^{-1/2}(1+O(T^{-1}))$, whereas the two long periods
have length $T^{1/2}(1+O(T^{-1}))$.
\end{corollary}

\begin{proof}
Equation \eqref{eq:radial-covolume-exact} gives the exact volume and
\eqref{eq:radial-holomorphic-errors} gives its expansion.  In real dimension
four, replacing $g$ by
$\alpha g$ multiplies volume by $\alpha^2$ and Laplace eigenvalues by
$\alpha^{-1}$.  Therefore \eqref{eq:unit-volume-metric} multiplies eigenvalues by $V_z^{1/2}$.  Combine this with Theorem~\ref{thm:cusp-gap}.
The period-length assertions follow directly from the columns of
$\mathsf P_z$: the last two have length one, while the first two have length
$T(1+O(T^{-1}))$; all lengths are multiplied by $\mathcal D_T^{-1/4}$.
\end{proof}

\begin{proposition}[Fixed-energy counting law]\label{prop:fixed-energy-count}
Along the radial cusp of Theorem~\ref{thm:cusp-gap}, retain the positive
covolume
\begin{equation}\label{eq:covolume-DT}
 \mathcal D_T=|\det\mathsf P_z|=-TD
 =\operatorname{Vol}(F_z^{\mathrm{flat}})
\end{equation}
from \eqref{eq:radial-covolume-exact}.
For $E\ge0$, count eigenvalues with their Fourier multiplicities by
\begin{equation}\label{eq:counting-definition}
 \lambda_T(k):=4\pi^2\|\mathsf P_z^{-\mathsf T}k\|^2,
 \qquad
 N_T(E):=\#\{k\in\Z^4:\lambda_T(k)\le E\}.
\end{equation}
Then, for every fixed $E$,
\begin{equation}\label{eq:fixed-energy-all-bands}
 N_T(E)=\pi\mathcal D_T
 \sum_{n\in\Z^2}
 \left(\frac{E}{4\pi^2}-\|n\|^2\right)_+
 +O_E(T),
 \qquad (x)_+:=\max\{x,0\}.
\end{equation}
The sum is finite.  In particular, throughout the first transverse band
$0<E<4\pi^2$,
\begin{equation}\label{eq:fixed-energy-first-band}
 \boxed{N_T(E)=\frac{E}{4\pi}\mathcal D_T+O_E(T).}
\end{equation}
Since $\mathcal D_T=T^2(1+O(T^{-1}))$, the number of modes below any fixed
energy in this band grows as $ET^2/(4\pi)+O_E(T)$.

For the unit-volume metric $\widehat g_z$ of
\eqref{eq:unit-volume-metric}, let
\begin{equation}\label{eq:unit-volume-count-definition}
 \widehat N_T(E):=
 \#\{k\in\Z^4:\lambda_T(k;\widehat g_z)\le E\}.
\end{equation}
For every fixed $E>0$,
\begin{equation}\label{eq:unit-volume-count}
 \widehat N_T(E)
 =\frac{E}{4\pi}\mathcal D_T^{1/2}
 +O_E(\mathcal D_T^{1/4}+1),
 \qquad
 \lim_{T\to\infty}
 \frac{\widehat N_T(E)}{\mathcal D_T^{1/2}}
 =\frac{E}{4\pi}.
\end{equation}

At the collapsing spectral scale define
\begin{equation}\label{eq:collapsing-count-definition}
 N_T^{\mathrm{coll}}(E):=
 \#\left\{k\in\Z^4:
 \lambda_T(k)\le\frac{4\pi^2E}{T^2}\right\}.
\end{equation}
If $E$ is not a sum of two integer squares, then for all sufficiently large
$T$,
\begin{equation}\label{eq:collapsing-count}
 N_T^{\mathrm{coll}}(E)
 =\#\{(r,s)\in\Z^2:r^2+s^2\le E\}.
\end{equation}
If $E$ is a sum of two integer squares, then
\begin{equation}\label{eq:threshold-count}
 \begin{split}
 \#\{(r,s):r^2+s^2<E\}
 &\le\liminf_{T\to\infty}N_T^{\mathrm{coll}}(E)\\
 &\le\limsup_{T\to\infty}N_T^{\mathrm{coll}}(E)
 \le\#\{(r,s):r^2+s^2\le E\}.
 \end{split}
\end{equation}
For a threshold vector $(r,s)$, the coefficient
$2b_{\mathrm I,0}r^2-10m_0rs$ in \eqref{eq:general-plane-mode-correction} decides
membership whenever it is nonzero; the exact formula
\eqref{eq:counting-exact-eigenvalue} and higher Taylor coefficients decide
the remaining cases.
\end{proposition}

\begin{proof}
Write
\[
 \alpha=\Re\mu,\qquad \eta=\Re\tau,\qquad c=\Re\beta,
\]
and index $k\in\Z^4$ as $k=(r,s,n_1,n_2)$.  Solving
$\mathsf P_z^{\mathsf T}x=k$ first gives $x_1=n_1$ and $x_3=n_2$.
The remaining two equations are
\begin{equation}\label{eq:counting-two-by-two-system}
 \mathcal A_T\binom{x_2}{x_4}
 =\binom{r}{s}-\theta_T(n),
 \quad
 \mathcal A_T:=\begin{pmatrix}6m&p\\T&m\end{pmatrix},
 \quad
 \theta_T(n):=\binom{6\alpha n_1+cn_2}{\eta n_1+\alpha n_2}.
\end{equation}
Here
\[
 \det\mathcal A_T=6m^2-pT=-TD=\mathcal D_T>0.
\]
Consequently the exact eigenvalue formula is
\begin{equation}\label{eq:counting-exact-eigenvalue}
 \frac{\lambda_T(r,s,n_1,n_2)}{4\pi^2}
 =\|n\|^2+
 \left\|\mathcal A_T^{-1}
 \left(\binom{r}{s}-\theta_T(n)\right)\right\|^2,
 \qquad n=(n_1,n_2).
\end{equation}

Put $R_E^2=E/(4\pi^2)$.  A fixed $n$ contributes nothing when
$\|n\|^2>R_E^2$.  Otherwise, with
$\rho_n=(R_E^2-\|n\|^2)^{1/2}$, its admissible pairs $(r,s)$ are precisely
the integer points of the translated ellipse
\begin{equation}\label{eq:counting-ellipse}
 \theta_T(n)+\mathcal A_T\overline B_{\R^2}(0,\rho_n).
\end{equation}
When $\rho_n>0$, this ellipse has area
$\pi\mathcal D_T\rho_n^2$ and perimeter $O_E(T)$, because
$m=O(1)$, $p=-T+O(1)$, and hence
$\opnorm{\mathcal A_T}=O(T)$.  When $\rho_n=0$, its contribution is zero or
one and is absorbed by $O_E(1)$.

For completeness, every translated compact convex set $K\subset\R^2$
satisfies
\begin{equation}\label{eq:area-perimeter-lattice-bound}
 \#(K\cap\Z^2)=\operatorname{area}(K)
 +O\bigl(\operatorname{length}(\partial K)+1\bigr),
\end{equation}
with an absolute implied constant.  Indeed, center a unit square at every
point of $K\cap\Z^2$ and let $U$ be their union.  The squares have disjoint
interiors and $\operatorname{area}(U)=\#(K\cap\Z^2)$.  Every point of
$U\mathbin\triangle K$ lies within distance $\sqrt2$ of $\partial K$:
this is immediate for $U\setminus K$, while a point of $K\setminus U$ farther
than $\sqrt2$ from the boundary would have its nearest lattice point in $K$
and hence lie in the corresponding square.  The inner and outer boundary
strips of a planar convex set of fixed width have area bounded by an absolute
constant times its perimeter plus one (first for polygons, by summing edge
rectangles and vertex sectors, and then by polygonal approximation).  This
proves \eqref{eq:area-perimeter-lattice-bound}.

Only finitely many $n$ occur for fixed $E$.  Applying
\eqref{eq:area-perimeter-lattice-bound} to
\eqref{eq:counting-ellipse}, and summing over those $n$, proves
\eqref{eq:fixed-energy-all-bands}.  If $0<E<4\pi^2$, only $n=0$ occurs, and
its area is
$\pi\mathcal D_T E/(4\pi^2)=E\mathcal D_T/(4\pi)$, proving
 \eqref{eq:fixed-energy-first-band}.

Under $\widehat g_z=\mathcal D_T^{-1/2}g_z$, every eigenvalue is multiplied
by $\mathcal D_T^{1/2}$, so
$\widehat N_T(E)=N_T(E/\mathcal D_T^{1/2})$.  For large $T$ only $n=0$
can occur in \eqref{eq:counting-exact-eigenvalue}.  Its ellipse has radius
parameter
\[
 \rho=\left(\frac{E}{4\pi^2\mathcal D_T^{1/2}}\right)^{1/2},
\]
area $E\mathcal D_T^{1/2}/(4\pi)$, and perimeter
$O_E(T\mathcal D_T^{-1/4}+1)=O_E(\mathcal D_T^{1/4}+1)$.
The same lattice estimate proves \eqref{eq:unit-volume-count}.

At the collapsing scale, every index with $(n_1,n_2)\ne(0,0)$ contributes
at least $4\pi^2$ and hence contributes nothing for large $T$.  For
$n_1=n_2=0$, equations \eqref{eq:MT}--\eqref{eq:MT-J} give uniformly
\[
 \frac{T^2\lambda_T(r,s,0,0)}{4\pi^2}
 =\|M_T(r,s)^{\mathsf T}\|^2\longrightarrow r^2+s^2.
\]
The lower singular-value bound \eqref{eq:singular-value-bound} restricts a
fixed rescaled window to a fixed finite set of pairs $(r,s)$.  Termwise
comparison proves \eqref{eq:collapsing-count} away from thresholds and
\eqref{eq:threshold-count} at thresholds.  The final refinement follows
from \eqref{eq:general-plane-mode-correction} and the exact eigenvalue
formula.
\end{proof}

\begin{remark}[What the determinant condition does and does not imply]
The pointwise condition $\det\mathsf P_z\ne0$ proves that every individual quotient is a compact torus and hence individually gapped.  It supplies no uniform lower bound over a noncompact parameter family.  The cusp gives an explicit failure of the implication
\[
 \text{pointwise nondegeneracy}\quad\Longrightarrow\quad
 \inf_z\lambda_1(z)>0.
\]
The failure persists after volume normalization, so it is anisotropic rather than a mere global dilation.
\end{remark}

\section{The native pointed limit and finite-cover unspooling}

Let $e_1,e_2,e_3,e_4$ be the standard basis in the real coordinate order
of \eqref{eq:Z-P}, and on the positive radial lift put
\begin{equation}\label{eq:primal-period-decomposition}
 v_{\gamma,T}:=\mathsf P_ze_\gamma,\qquad
 v_{u,T}:=\mathsf P_ze_u,\qquad
 \mathsf P_z\Z^4
 =\Z v_{\gamma,T}\oplus\Z v_{u,T}
  \oplus\Z e_1\oplus\Z e_3.
\end{equation}
Define the fixed real two-torus
\begin{equation}\label{eq:native-fixed-torus}
 \T^2_{\mathrm{re}}
 :=(\R e_1\oplus\R e_3)/(\Z e_1\oplus\Z e_3).
\end{equation}

\begin{theorem}[Native limit and expanding marked covers]
\label{thm:native-and-cover-limits}
Let $o_T$ be the origin of $F_z^{\mathrm{flat}}$.  With the unscaled
Euclidean quotient metric,
\begin{equation}\label{eq:native-pointed-limit}
 (F_z^{\mathrm{flat}},o_T)
 \longrightarrow
 (\T^2_{\mathrm{re}}\times(\R e_2\oplus\R e_4),(0,0))
 \cong(\T^2\times\R^2,(0,0))
\end{equation}
in pointed smooth and pointed Gromov--Hausdorff topology.  For every fixed
$r<\infty$, the radius-$r$ balls are exactly isometric for all sufficiently
large $T$.  Moreover, eventually
\begin{equation}\label{eq:native-systole}
 \operatorname{sys}(\mathsf P_z\Z^4)=1,
 \qquad \operatorname{inj}(F_z^{\mathrm{flat}})=\frac12.
\end{equation}

For $M\in\Z_{\ge1}$ define
\begin{equation}\label{eq:cover-sublattice}
 \Lambda_{T,M}
 :=\Z v_{\gamma,T}\oplus\Z v_{u,T}
   \oplus\Z(Me_1)\oplus\Z(Me_3),
 \qquad
 \widetilde F_{T,M}:=\R^4/\Lambda_{T,M}.
\end{equation}
The identity on $\R^4$ induces a regular holomorphic local-isometric cover
\begin{equation}\label{eq:degree-cover}
 c_{T,M}:\widetilde F_{T,M}\longrightarrow F_z^{\mathrm{flat}}
\end{equation}
of degree $M^2$, with deck group $(\Z/M\Z)^2$.  There are constants
$C,C'>0$, independent of large $T$ and $M$, for which
\begin{equation}\label{eq:cover-systole-two-sided}
 \min\{T-C,M\}\le\operatorname{sys}(\Lambda_{T,M})
 \le\min\{T+C',M\}.
\end{equation}
Consequently, for an integer-valued $M=M(T)$,
\begin{equation}\label{eq:cover-R4-limit}
 (\widetilde F_{T,M(T)},\widetilde o_{T,M(T)})
 \longrightarrow(\R^4,0)
\end{equation}
in pointed smooth and pointed Gromov--Hausdorff topology if and only if
$M(T)\to\infty$.

If the cover instead carries the pullback of the base-fibre unit-volume
metric $\widehat g_z=\mathcal D_T^{-1/2}g_z$, then the same two pointed
limits are $\R^4$ if and only if
\begin{equation}\label{eq:unit-volume-cover-condition}
 \frac{M(T)}{\mathcal D_T^{1/4}}\longrightarrow\infty,
 \qquad\text{equivalently}\qquad
 \frac{M(T)}{\sqrt T}\longrightarrow\infty.
\end{equation}
By contrast, if each covering torus is rescaled by
\begin{equation}\label{eq:cover-own-unit-volume-metric}
 g^{\mathrm{cov},1}_{T,M}:=(M^2\mathcal D_T)^{-1/2}g_z
\end{equation}
so that the cover itself has volume one, no choice of $M(T)$ gives pointed
smooth limit $\R^4$.
\end{theorem}

\begin{proof}
Project the two long periods in \eqref{eq:primal-period-decomposition} to
the $(e_2,e_4)$-plane.  Their projected column matrix is
\begin{equation}\label{eq:long-plane-matrix}
 Y_T=\begin{pmatrix}6m&T\\B_R-T&m\end{pmatrix}
 =T\begin{pmatrix}0&1\\-1&0\end{pmatrix}
  +\begin{pmatrix}6m&0\\B_R&m\end{pmatrix}.
\end{equation}
The second matrix is uniformly bounded.  Hence, after increasing $C$,
\begin{equation}\label{eq:long-plane-lower-bound}
 \left\|Y_T\binom ab\right\|
 \ge(T-C)\sqrt{a^2+b^2}
 \qquad(a,b\in\R).
\end{equation}
First quotient $\R^4$ by $\Z e_1\oplus\Z e_3$.  The result is the product
$\T^2_{\mathrm{re}}\times(\R e_2\oplus\R e_4)$.  Every nonidentity
remaining deck translation has displacement at least $T-C$ by
\eqref{eq:long-plane-lower-bound}.  If $T-C>4r$, no such translate can
identify two points of the radius-$r$ ball or shorten their distance, while
every minimizing quotient geodesic of length at most $r$ lifts from the
origin.  The two radius-$r$ balls are therefore exactly isometric.  This
proves \eqref{eq:native-pointed-limit}.  A nonzero period with a nonzero long
coefficient has length at least $T-C$, whereas the shortest nonzero vector
of $\Z e_1\oplus\Z e_3$ has length one; this proves
\eqref{eq:native-systole}.

The four vectors in \eqref{eq:primal-period-decomposition} are a lattice
basis, so \eqref{eq:cover-sublattice} has index $M^2$ and quotient
$(\Z/M\Z)^2$.  Finite-index inclusion of complex lattices makes
\eqref{eq:degree-cover} holomorphic and local-isometric.  A nonzero vector
\[
 av_{\gamma,T}+bv_{u,T}+M(pe_1+qe_3)
\]
has length at least $T-C$ if $(a,b)\ne(0,0)$, by its plane projection, and
at least $M$ otherwise.  Conversely, $Me_1$ has length $M$, while either
long generator has length at most $T+C'$ for a uniform $C'$.  This proves
\eqref{eq:cover-systole-two-sided}.  A flat quotient has radius-$r$ balls
exactly Euclidean once its lattice systole exceeds $4r$; hence $M(T)\to
\infty$ proves \eqref{eq:cover-R4-limit}.  If $M(T)$ does not tend to
infinity, integrality gives a subsequence on which $M(T)=M_*$ is constant.
Quotienting first by $M_*\Z e_1\oplus M_*\Z e_3$ and repeating the preceding
exact-ball argument shows that this subsequence converges to
$(\R/M_*\Z)^2\times\R^2$, not to $\R^4$.  This excludes both asserted
$\R^4$ topologies and proves necessity.

Under the pullback of $\widehat g_z$, all lengths are multiplied by
$\mathcal D_T^{-1/4}$.  The long lower bound becomes
$(T-C)\mathcal D_T^{-1/4}\sim\sqrt T$ and diverges automatically, while
$Me_1$ has length $M\mathcal D_T^{-1/4}$.  This proves the equivalence
\eqref{eq:unit-volume-cover-condition} in the sufficient direction.  For
necessity, if $\ell_T=M(T)\mathcal D_T^{-1/4}$ does not tend to infinity,
pass to a subsequence with $\ell_T\to\ell<\infty$.  Once the remaining long
translations leave every fixed ball, the balls are those of
$(\R/\ell_T\Z)^2\times\R^2$.  If $\ell>0$ the pointed limit is
$(\R/\ell\Z)^2\times\R^2$; if $\ell=0$ the pointed
Gromov--Hausdorff limit is $\R^2$ and there is collapse, so there is no
pointed smooth four-dimensional limit.  In neither case is the limit $\R^4$.
Finally, under
\eqref{eq:cover-own-unit-volume-metric}, the explicit periods $Me_1$ and
$v_{\gamma,T}$ have respective lengths
\begin{equation}\label{eq:cover-own-unit-volume-lengths}
 \sqrt{\frac MT}\,(1+o(1)),\qquad
 \sqrt{\frac TM}\,(1+o(1)).
\end{equation}
Their leading factors have product one, so at least one stays bounded.  The
injectivity radius cannot tend to infinity, which excludes a pointed smooth
$\R^4$ limit.
\end{proof}

\begin{theorem}[Base-changed monodromy and exact cover momenta]
\label{thm:cover-momentum}
Put
\begin{equation}\label{eq:cover-index-matrix}
 \mathsf S_M=\operatorname{diag}(1,1,M,M),
 \qquad \Lambda^{(M)}=\mathsf S_M\Z^4.
\end{equation}
Then $\mathsf P_z\Lambda^{(M)}=\Lambda_{T,M}$.  Writing
$M_0=I+R_0$ for the cusp monodromy in
\eqref{eq:cusp-monodromy-matrix}, one has
\begin{equation}\label{eq:cover-monodromy-algebra}
 R_0^2=0,\qquad
 \ker R_0=\operatorname{im}R_0
 =\Z\langle\widehat w,\widehat\delta\rangle,
 \qquad
 \mathsf S_M^{-1}M_0^M\mathsf S_M=M_0.
\end{equation}
Thus, after the cyclic base change $t_c=v^M$, the covers form a regular
relative degree-$M^2$ cover over $v\ne0$.  This is punctured descent; the
completed local extension and its different absolute and relative degrees
are treated in Section~\ref{cusp-cover-section}.

With normalized Haar measure, the complete orthonormal character basis of
$\widetilde F_{T,M}$ is
\begin{equation}\label{eq:cover-characters}
 \widetilde\chi_{T,M,k}([x])
 =\exp\!\left(2\pi i
 \left\langle\mathsf P_z^{-\mathsf T}\mathsf S_M^{-1}k,x\right\rangle
 \right),\qquad k\in\Z^4,
\end{equation}
and its scalar spectrum, with multiplicity, is
\begin{equation}\label{eq:cover-scalar-spectrum}
 \Spec(\widetilde\Delta_{T,M})
 =\left\{4\pi^2
 \left\|\mathsf P_z^{-\mathsf T}\mathsf S_M^{-1}k\right\|^2:
 k\in\Z^4\right\}.
\end{equation}
The pullback, equivalently deck-invariant, sector consists exactly of
\begin{equation}\label{eq:cover-invariant-sector}
 k\in\mathsf S_M\Z^4
 \quad\Longleftrightarrow\quad
 k_3\equiv k_4\equiv0\pmod M.
\end{equation}
There is a constant $C_{\mathrm{mom}}$, independent of large $T$ and
$M\ge1$, such that
\begin{equation}\label{eq:cover-momentum-mesh}
 \sup_{p\in\R^4}\operatorname{dist}\!\left(
 p,\mathsf P_z^{-\mathsf T}\mathsf S_M^{-1}\Z^4\right)
 \le C_{\mathrm{mom}}\left(\frac1T+\frac1M\right).
\end{equation}
Hence if $T_j,M_j\to\infty$, the cover momentum lattices become locally
dense in $\R^4$, and every $E\ge0$ is a limit of eigenvalues in
\eqref{eq:cover-scalar-spectrum}.
\end{theorem}

\begin{proof}
From \eqref{eq:cusp-monodromy-matrix},
$R_0(a,b,c,d)^{\mathsf T}=(0,0,b,-a)^{\mathsf T}$.  This proves the kernel,
image, and square-zero assertions.  The binomial theorem gives
$M_0^M=I+MR_0$, and direct evaluation on $(a,b,c,d)^{\mathsf T}$ gives
\[
 \mathsf S_M^{-1}M_0^M\mathsf S_M(a,b,c,d)^{\mathsf T}
 =(a,b,c+b,d-a)^{\mathsf T}=M_0(a,b,c,d)^{\mathsf T}.
\]
Thus $M_0^M$ preserves $\Lambda^{(M)}$, which is exactly descent after
$t_c=v^M$.

The dual of $\mathsf P_z\mathsf S_M\Z^4$ is
$\mathsf P_z^{-\mathsf T}\mathsf S_M^{-1}\Z^4$.  Fourier theory proves
\eqref{eq:cover-characters}--\eqref{eq:cover-scalar-spectrum}.  Moreover,
$c_{T,M}^*\chi_{T,\ell}=\widetilde\chi_{T,M,\mathsf S_M\ell}$, so a cover
character descends exactly under the congruence in
\eqref{eq:cover-invariant-sector}.

Let $b_{j,T,M}$ be the columns of
$\mathsf P_z^{-\mathsf T}\mathsf S_M^{-1}$.  The exact system
\eqref{eq:counting-two-by-two-system}, boundedness of all holomorphic cusp
data, and $\mathcal D_T\asymp T^2$ give a uniform constant $C_0$ with
\[
 \|b_{1,T,M}\|+\|b_{2,T,M}\|\le\frac{C_0}{T},
 \qquad
 \|b_{3,T,M}\|+\|b_{4,T,M}\|\le\frac{C_0}{M}.
\]
Express any $p\in\R^4$ uniquely as $p=\sum_jt_jb_{j,T,M}$ and choose
$n_j\in\Z$ with $|t_j-n_j|\le1/2$.  Then
\[
 \left\|p-\sum_jn_jb_{j,T,M}\right\|
 \le\frac12\sum_j\|b_{j,T,M}\|
 \le\frac{C_0}{2}\left(\frac1T+\frac1M\right),
\]
which proves \eqref{eq:cover-momentum-mesh}.  For $E>0$, approximate a
vector $p$ with $4\pi^2\|p\|^2=E$; continuity gives eigenvalues tending to
$E$.  For $E=0$, take $k=e_\gamma$, whose momentum is $O(T^{-1})$.
\end{proof}

\section{The invariant real quotient and its four-copy limit}

Define the short-plane lattice, the primitive integral character, and its
kernel lattice by
\begin{equation}\label{eq:invariant-character}
 \begin{aligned}
 K_\Lambda&=\Z\langle\widehat w,\widehat\delta\rangle,\\
 \varepsilon_\gamma
 (a\widehat\gamma+b\widehat u+c\widehat w+d\widehat\delta)&=a,
 &L_\Lambda=\ker\varepsilon_\gamma
 &=\Z\langle\widehat u,\widehat w,\widehat\delta\rangle.
 \end{aligned}
\end{equation}

\begin{theorem}[Regular-locus quotient and exact radial metric]
\label{thm:global-real-quotient}
Assume the conditional global period family and the integral monodromy data
in Section~1.  Then the following statements hold.

\begin{enumerate}[label=(\roman*),leftmargin=2.3em]
\item The formula
\begin{equation}\label{eq:global-circle-map}
 \chi_\gamma([z,\mathsf P_zx])
 =\bigl([z],\varepsilon_\gamma(x)\bmod\Z\bigr)
\end{equation}
defines a smooth real-analytic fibrewise group epimorphism
\begin{equation}\label{eq:global-circle-quotient}
 \chi_\gamma:J\longrightarrow B^\circ\times(\R/\Z).
\end{equation}
Its kernel is a connected real-three-torus family with fibre
$(L_\Lambda\otimes\R)/L_\Lambda$.  The kernel family need not be a product.
The quotient is real; it is not a holomorphic quotient of complex tori.
Moreover,
\begin{equation}\label{eq:invariant-character-lattice}
 \operatorname{Hom}(\Lambda,\Z)^\Delta=\Z\varepsilon_\gamma,
\end{equation}
so this primitive invariant character is unique up to sign.

\item Give each marked fibre the Euclidean quotient metric already used in
Section~2 and give the target circle its Riemannian quotient metric.  On a
radial cusp fibre,
\begin{equation}\label{eq:radial-circle-character}
 \chi_{\gamma,z}([y])
 =\left\langle\mathsf P_z^{-\mathsf T}e_\gamma,y\right\rangle\bmod\Z,
 \qquad e_\gamma=(1,0,0,0)^{\mathsf T},
\end{equation}
and its circumference is exactly
\begin{equation}\label{eq:circle-circumference}
 \ell_\gamma(z)
 =\frac{1}{\lVert\mathsf P_z^{-\mathsf T}e_\gamma\rVert}
 =\frac{-TD}{\sqrt{T^2+m^2}}.
\end{equation}
If $b_{\mathrm I,0}=\operatorname{Im}b(0)$, then on the positive radial lift
\begin{equation}\label{eq:circle-circumference-asymptotic}
 \ell_\gamma(z)=T-b_{\mathrm I,0}+O(T^{-1})\longrightarrow\infty.
\end{equation}
The pullback of the first circle Fourier mode is the fibre character with
index $e_\gamma$, whose eigenvalue is $4\pi^2/\ell_\gamma(z)^2$.  The
character is global, but the chosen Euclidean fibre metrics need not be
invariant under the deck transitions.  On the positive radial marked lift,
this mode gap-closes with the same leading rate as
Theorem~\ref{thm:cusp-gap}; it is not asserted to be the exact shortest mode
at every radial parameter.

\item Let
\begin{equation}\label{eq:fourfold-fibre-product}
 J^{[4]}=J\times_{B^\circ}J\times_{B^\circ}J\times_{B^\circ}J.
\end{equation}
The fourfold product of \eqref{eq:global-circle-map} is a quotient
\begin{equation}\label{eq:fourfold-circle-map}
 \chi_\gamma^{[4]}:J^{[4]}\longrightarrow
 B^\circ\times(\R/\Z)^4
\end{equation}
with connected twelve-torus kernel.  Along the synchronized positive radial
cusp, its quotient fibre with the induced product metric is
\begin{equation}\label{eq:four-circle-fibre}
 G_z=(\R/\ell_\gamma(z)\Z)^4.
\end{equation}
Its systole is $\ell_\gamma(z)$.
Consequently
\begin{equation}\label{eq:four-circle-limit}
 (G_z,0)\longrightarrow(\R^4,0)
\end{equation}
in pointed local-isometric, pointed smooth, and pointed
Gromov--Hausdorff topology.  For every fixed $r>0$, the radius-$r$ ball is
exactly Euclidean once $\ell_\gamma(z)>4r$.

Equivalently, the external Cartesian product defines
\begin{equation}\label{eq:external-four-circle-map}
 J^4\longrightarrow(B^\circ)^4\times(\R/\Z)^4,
\end{equation}
and $J^{[4]}$ is its pullback along the diagonal
$B^\circ\to(B^\circ)^4$.  At four independent positive radial cusp
parameters $z_1,\ldots,z_4$, its quotient lattice is
\[
 \operatorname{diag}\bigl(\ell_\gamma(z_1),\ldots,
 \ell_\gamma(z_4)\bigr)\Z^4,
\]
so the same three pointed limits hold whenever
$\min_i\ell_\gamma(z_i)\to\infty$.

Among connected torus quotients of $n$ synchronized copies defined by a
fixed primitive full-monodromy-invariant kernel lattice and equipped with the
Riemannian quotient of the product Euclidean metric, any quotient whose
systole tends to infinity has rank at most $n$.  Hence four copies are
minimal for a rank-four all-expanding quotient in this class, and
\eqref{eq:fourfold-circle-map} attains the bound.
\end{enumerate}
\end{theorem}

\begin{proof}
For a row covector $(a,b,c,d)$, invariance under the first matrix in
\eqref{eq:monodromy-matrices} gives $b=2c$ and $d=3c$, whereas invariance
under the second gives $b=c$ and $d=2c$.  Thus $b=c=d=0$, which proves
\eqref{eq:invariant-character-lattice}.  In particular
$\varepsilon_\gamma A_g=\varepsilon_\gamma$ for every $g\in\Delta$.
Equation \eqref{eq:flat-monodromy-action} therefore proves that
\eqref{eq:global-circle-map} is well defined.  Primitivity of
$\varepsilon_\gamma$ proves surjectivity and connectedness of its kernel,
and the induced action on $\Lambda/L_\Lambda\simeq\Z$ is trivial.  Hence the
target circle family is the displayed product.  Since the kernel has odd real
dimension, it is not a complex subtorus, proving (i).

Putting $(r,s)=(1,0)$ in \eqref{eq:dual-exact} gives
\begin{equation}\label{eq:circle-dual-vector}
 \mathsf P_z^{-\mathsf T}e_\gamma
 =\left(0,\frac{m}{-TD},0,\frac{-T}{-TD}\right)^{\mathsf T}.
\end{equation}
The length of a Riemannian quotient circle defined by a primitive dual
character is the reciprocal norm of that character, which proves
\eqref{eq:circle-circumference}.  Write
$b_{\mathrm I}(t_c)=\operatorname{Im}b(t_c)$.  Since
$p=b_{\mathrm I}(t_c)-T$,
\begin{equation}\label{eq:circumference-numerator}
 -TD=T^2-Tb_{\mathrm I}(t_c)+6m^2.
\end{equation}
Holomorphy at $t_c=0$ gives
$b_{\mathrm I}(t_c)=b_{\mathrm I,0}+O(e^{-2\pi R})$ and
$m=\operatorname{Im}\mu(0)+O(e^{-2\pi R})$ on the positive radial lift.
Because $T=R+O(1)$, expansion of
\eqref{eq:circle-circumference} proves
\eqref{eq:circle-circumference-asymptotic}.  Formula
\eqref{eq:radial-circle-character} is the standard character formula from
Proposition~\ref{prop:fixed-spectrum}, and its exponential is precisely the
index-$e_\gamma$ eigenfunction.  This proves (ii).

Taking four products proves the quotient and kernel statements in (iii).
After multiplying each angular target coordinate by $\ell_\gamma(z)$, its
fibre is $\R^4/\ell_\gamma(z)\Z^4$.  If $x,y\in B_r(0)$ and
$0\ne\lambda\in\ell_\gamma(z)\Z^4$, then
\[
 \lVert x-y+\lambda\rVert
 \ge \lVert\lambda\rVert-\lVert x-y\rVert
 >\ell_\gamma(z)-2r>2r\ge\lVert x-y\rVert
\]
when $\ell_\gamma(z)>4r$.  The covering map is therefore an exact isometry
on the radius-$r$ ball, proving all three pointed convergences.  Taking the
four quotient maps over independent base parameters gives
\eqref{eq:external-four-circle-map}; the identical estimate with
$\ell_\gamma(z)$ replaced by $\min_i\ell_\gamma(z_i)$ proves its pointed
limit and pulling back along the diagonal recovers $J^{[4]}$.

For minimality, equation \eqref{eq:Z-P} gives
$\mathsf P_z\widehat w=(1,0,0,0)^{\mathsf T}$ and
$\mathsf P_z\widehat\delta=(0,0,1,0)^{\mathsf T}$.  Thus each fixed element
of $K_\Lambda^{\oplus n}$ has length bounded independently of the radial
parameter.  Let $N\subset\Lambda^{\oplus n}$ be the fixed primitive invariant
kernel lattice of an all-expanding quotient.  If a fixed
$k\in K_\Lambda^{\oplus n}$ did not lie in $N$, primitivity would give
$N_\R\cap\Lambda^{\oplus n}=N$, so $k$ would determine a nonzero quotient
period of bounded length, contradicting divergence of the quotient systole.
Therefore $K_\Lambda^{\oplus n}\subset N$.

The first matrix in \eqref{eq:monodromy-matrices} satisfies
$A_1\widehat w=\widehat u-\widehat w$.  Invariance of $N$ consequently gives
$L_\Lambda^{\oplus n}\subset N$, and hence
\[
 \operatorname{rank}(\Lambda^{\oplus n}/N)
 \le 4n-3n=n.
\]
Rank four requires $n\ge4$, while
$N=L_\Lambda^{\oplus4}$ gives the quotient already constructed.
\end{proof}

\begin{proposition}[The expanding quotient of the native cusp fibre]
\label{prop:native-cusp-quotient}
On a radial cusp fibre define
\[
 \rho:\R^4\longrightarrow\R^2,
 \qquad \rho(x_1,x_2,x_3,x_4)=(x_2,x_4),
\]
and
\begin{equation}\label{eq:native-long-matrix}
 Y_z=\begin{pmatrix}6m&T\\p&m\end{pmatrix}.
\end{equation}
Then $\rho$ induces an exact sequence
with
\[
 \T^2_{\mathrm{re}}
 :=\frac{\R\langle e_1,e_3\rangle}{\Z\langle e_1,e_3\rangle},
\]
namely
\begin{equation}\label{eq:native-cusp-exact-sequence}
 0\longrightarrow\T^2_{\mathrm{re}}\longrightarrow
 F_z^{\mathrm{flat}}\longrightarrow
 Q_z:=\R^2/Y_z\Z^2\longrightarrow0,
\end{equation}
and $\operatorname{sys}(Q_z)\ge T-C$ for large $T$.  Thus $(Q_z,0)$ has
pointed local-isometric limit $(\R^2,0)$, and two synchronized copies satisfy
\begin{equation}\label{eq:native-quotient-two-copy-limit}
 (Q_z\times Q_z,0)\longrightarrow(\R^4,0).
\end{equation}
The native fibre itself still has the different pointed limit
$\T^2_{\mathrm{re}}\times\R^2$ from
Theorem~\ref{thm:native-and-cover-limits}.
The one-copy quotient has kernel lattice $K_\Lambda$, and the displayed
two-copy quotient has kernel lattice $K_\Lambda^{\oplus2}$.  This is not the
global quotient of four copies in Theorem~\ref{thm:global-real-quotient}.
No finite cover of $B^\circ$ makes either of these particular kernels
full-monodromy invariant.
\end{proposition}

\begin{proof}
The last two columns of \eqref{eq:Z-P} are killed by $\rho$, and the
projections of the first two columns are the columns of $Y_z$.  Moreover,
\[
 \det Y_z=6m^2-Tp=-TD>0.
\]
If $\rho(x)=Y_z(a,b)^{\mathsf T}$, then
$x-\mathsf P_z(a,b,0,0)^{\mathsf T}$ lies in
$\ker\rho=\R\langle e_1,e_3\rangle$.  Conversely, if
$\mathsf P_zk\in\ker\rho$ for $k\in\Z^4$, then
$Y_z(k_1,k_2)^{\mathsf T}=0$; invertibility of $Y_z$ gives $k_1=k_2=0$.
Thus
\[
 \ker\bigl(F_z^{\mathrm{flat}}\longrightarrow Q_z\bigr)
 \simeq
 \frac{\R\langle e_1,e_3\rangle}
      {\Z\langle e_1,e_3\rangle}
 =\T^2_{\mathrm{re}},
\]
which proves exactness without identifying this kernel with the invariant
kernel $L_\Lambda$.  By
\eqref{eq:cusp-asymptotic-input},
\[
 T^{-1}Y_z=
 \begin{pmatrix}0&1\\-1&0\end{pmatrix}+O_{\mathrm{op}}(T^{-1}).
\]
The least singular value of $Y_z$ is therefore at least $T-C$, and the same
is true of the shortest nonzero vector of $Y_z\Z^2$.  The pointed-limit
argument in the proof of Theorem~\ref{thm:native-and-cover-limits} now applies
first to $Q_z$ and then to $Q_z^2$.

For the finite-cover assertion, direct multiplication gives
\begin{equation}\label{eq:no-finite-cover-matrix}
 H_*=A_2A_1=
 \begin{pmatrix}
 1&0&0&0\\6&1&1&0\\0&0&1&0\\-5&0&-1&1
 \end{pmatrix}=I+N_*,\qquad N_*^2=0.
\end{equation}
Consequently, for every nonzero integer $n$,
\begin{equation}\label{eq:no-finite-cover-orbit}
 H_*^n\widehat w
 =\widehat w+n(\widehat u-\widehat\delta),
 \qquad H_*^n\widehat\delta=\widehat\delta,
\end{equation}
so $H_*^nK_\Lambda\ne K_\Lambda$.  Choose a loop in $B^\circ$ whose
monodromy is $H_*$.  Every finite-index subgroup of the ordinary
$\pi_1(B^\circ)$ contains a positive power of this loop, and that power still
fails to preserve $K_\Lambda$.  Hence no finite base cover globalizes this
particular one-copy marked quotient.  Under diagonal monodromy,
$(\widehat w,0)$ proves the same assertion for $K_\Lambda^{\oplus2}$ and thus
for its two-copy product.
\end{proof}

\begin{remark}[Copies, quotients, and covers]
The native fibre limit in Theorem~\ref{thm:native-and-cover-limits} is
$\T^2\times\R^2$.  The same theorem obtains $\R^4$ only after passing to
degree-$M(T)^2$ fibre covers with $M(T)\to\infty$.  Neither statement is a
quotient construction.  By contrast, the four-circle limit in
Theorem~\ref{thm:global-real-quotient} changes
neither the base nor the fibre lattice: it takes a quotient of four
synchronized copies.  Proposition~\ref{prop:native-cusp-quotient} instead
uses a rank-two marked-fibre quotient whose all-expanding behavior is radial
cusp behavior and whose kernel does not descend over the full regular base.
These quotient constructions are distinct from the degree-$M^2$ fibre covers,
whose punctured descent requires the base change $t_c=v^M$; no such
finite-cover construction is used in this paper's globalization theorem.
\end{remark}

\begin{proposition}[Literal fixed-target nonextension through the toric cusp]
\label{prop:circle-cusp-nonextension}
Assume the toric cusp chart asserted for the proposed completion.  Then
neither $\chi_\gamma$ to the fixed product-circle target nor
$\chi_\gamma^{[4]}$ to the fixed product $\mathbb T^4$ target extends
continuously through the toric cusp fibre with its stated regular-fibre
restriction.
\end{proposition}

\begin{proof}
On the positive radial lift $s=iR$, equations \eqref{eq:Z-P} and
\eqref{eq:cusp-s-coordinate} give
\[
 Z_R=iR B_0+C(t_c),\qquad
 B_0=\begin{pmatrix}0&1\\-1&0\end{pmatrix},\qquad
 C(t_c)=\begin{pmatrix}6\mu&h\\b-h&\mu\end{pmatrix},
 \qquad t_c=e^{-2\pi R}.
\]
For $a\in\R^2$, let
$x_R(a)=\exp(2\pi iZ_Ra)\in(\C^*)^2$ in the dense torus of the cusp model.
Its normalized logarithmic position is
\begin{equation}\label{eq:cusp-log-position}
 \frac{(\log|x_{R,1}(a)|,\log|x_{R,2}(a)|)}{\log|t_c|}
 =B_0a+\frac{\operatorname{Im}C(t_c)a}{R}.
\end{equation}
In the affine toric chart belonging to
$\sigma=\operatorname{conv}(0,e_1,e_2)$, the coordinates are
\[
 z_0=\frac{t_c}{x_1x_2},\qquad z_1=x_1,\qquad z_2=x_2.
\]
Writing the left side of \eqref{eq:cusp-log-position} as
$y_R=(y_{R,1},y_{R,2})$, one has exactly
\[
 |z_0|=|t_c|^{1-y_{R,1}-y_{R,2}},\qquad
 |z_1|=|t_c|^{y_{R,1}},\qquad
 |z_2|=|t_c|^{y_{R,2}}.
\]
If the limit of $y_R$ lies in the interior of $\sigma$, all three exponents
are eventually positive, so all three coordinates tend to zero and the path
approaches the same toric fixed point.

Take
\[
 a=\left(-\frac14,\frac14\right),\qquad
 a'=\left(-\frac12,\frac14\right).
\]
Then
\[
 B_0a=\left(\frac14,\frac14\right),\qquad
 B_0a'=\left(\frac14,\frac12\right)
\]
are both interior to $\sigma$, whereas
\[
 Z_Ra=\mathsf P_z(a_1,a_2,0,0)^{\mathsf T},\qquad
 Z_Ra'=\mathsf P_z(a'_1,a'_2,0,0)^{\mathsf T},
\]
in the real coordinate identification, and hence
\[
 \chi_{\gamma,R}(x_R(a))=\frac34,
 \qquad
 \chi_{\gamma,R}(x_R(a'))=\frac12
 \quad\text{in }\R/\Z.
\]
Thus two paths converge to the same cusp point while their character values
have different limits.  This rules out continuous extension of the literal
map to the fixed circle target.  Holding the other three factors on identical
radial paths proves the assertion for the fixed product $\mathbb T^4$ target.
\end{proof}

\begin{corollary}[The fixed $\mathbb T^4$ target and the smash collapse]
\label{cor:product-smash-obstruction}
Assume in addition the claimed recognition $X\cong_{\mathrm{diff}}S^6$.
Then
\[
 X^4\cong_{\mathrm{diff}}(S^6)^4\not\simeq S^{24},
 \qquad H_6(X^4;\Z)\cong\Z^4.
\]
After choices of base point and of the recognition, one has
$X^{\wedge4}\cong S^{24}$, but the literal fourfold quotient to the fixed
$\mathbb T^4$ target has no extension that factors through that smash
sphere.  More precisely, no map
$\Phi:X^4\to(\R/\Z)^4$ restricts on a regular fourfold fibre to
$(\chi_{\gamma,z})^4$, and no map
$\Psi:X^{\wedge4}\to(\R/\Z)^4$ has that restriction after precomposition
with the smash-collapse map $X^4\to X^{\wedge4}$.
These statements obstruct only the displayed character quotient with the
fixed target $(\R/\Z)^4$ and the stated regular-fibre restriction; they do
not rule out codomain-changing cone targets or any other real or singular
globalization mechanism.
\end{corollary}

\begin{proof}
The product homology statement follows from the integral Kunneth theorem, so
$X^4$ is not a homotopy sphere.  The smash identification follows from a
based recognition $X\cong S^6$.  On
$F_z^{\mathrm{flat}}\times\{0\}^3$, the fourfold quotient is
\[
 x\longmapsto(\chi_{\gamma,z}(x),0,0,0),
\]
which is nonconstant.  It is therefore not constant on the fat wedge and
does not factor through the fibre smash product.

For the stronger extension statement, $H^1(X^4;\Z)=0$.  If a map
\[
 \Phi:X^4\longrightarrow(\R/\Z)^4
\]
restricted on one regular fourfold fibre to
$(\chi_{\gamma,z})^4$, the pullback of the
first target $H^1$-generator would restrict to the nonzero class
$\varepsilon_\gamma$ on the first fibre factor, while every class pulled
back to $X^4$ is zero.  This is impossible.  A map from $X^{\wedge4}$ with
the corresponding restriction would give such a $\Phi$ after precomposition
with the smash-collapse map $X^4\to X^{\wedge4}$, so it is impossible as
well.
\end{proof}

\section{The completed local toric Kummer morphism}

The spectral degeneration and the invariant-character quotient do not by
themselves extend a fibre cover across the cusp.  The following local theorem
performs that separate task.  Its source quotient uses the modified action
with $C_N(v)=C(v^N)/N$; consequently the toric power map is equivariant.
The degree to the original cusp family is $N^3$, whereas $N^2$ is the generic
degree only after factoring through the base-changed target.  The construction
is local at the cusp and makes no full-base claim.

\subsection{The Kummer cover extends across the toric cusp}
\label{cusp-cover-section}

This section determines exactly what is true in the assertion that the
degree-$N^2$ cover of the regular fibres extends by a toric power map.  The
answer is positive for the completed \emph{local} cusp family, but the source
action has to be changed, and two different degrees must not be conflated.

Retain the cusp data
\[
 B_0=\begin{pmatrix}0&1\\-1&0\end{pmatrix},\qquad
 Z(s)=sB_0+C(t),\qquad t=e^{2\pi i s},
\]
where $C$ is holomorphic on a disc $\Delta_t$ about zero.  Put
\[
 c_{\bar\lambda}(t)
 :=\exp\!\bigl(2\pi iC(t)\bar\lambda\bigr)
 \in(\mathbb C^*)^2,
 \qquad \bar\lambda\in\bar\Lambda\simeq\mathbb Z^2.
\]
Let $\Sigma$ be the fan in
$N'=\mathbb Z^2\oplus\mathbb Z$ consisting of the cones over the cells of
the translation-invariant $A_2$ triangulation at height one, and write
$Y=Y_\Sigma$.  On the dense torus, with coordinates $(x_1,x_2,t)$, the
target action is
\begin{equation}
 \Psi_{\bar\lambda}(x,t)
 =\bigl(c_{\bar\lambda}(t)t^{B_0\bar\lambda}x,t\bigr).
 \label{cusp-cover-target-action}
\end{equation}
Its quotient is the completed cusp family
$f_0\colon N_0=Y/\bar\Lambda\to\Delta_t$.

Fix an integer $N\geq1$.  Let $v$ be a coordinate on a second disc and
make the base change
\begin{equation}
 \rho_N\colon\Delta_v\longrightarrow\Delta_t,\qquad
 \rho_N(v)=v^N,\qquad t=v^N.
 \label{cusp-cover-base-change}
\end{equation}
On a copy $Y^{(N)}$ of the same toric variety, with dense-torus coordinates
$(\xi_1,\xi_2,v)$, define
\begin{equation}
 C_N(v):=\frac1N C(v^N),\qquad
 c^{(N)}_{\bar\lambda}(v)
 :=\exp\!\bigl(2\pi iC_N(v)\bar\lambda\bigr)
 \label{cusp-cover-source-factor}
\end{equation}
and
\begin{equation}
 \Psi^{(N)}_{\bar\lambda}(\xi,v)
 :=\bigl(c^{(N)}_{\bar\lambda}(v)
          v^{B_0\bar\lambda}\xi,v\bigr).
 \label{cusp-cover-source-action}
\end{equation}

\begin{theorem}[Exact Kummer extension through the cusp]
\label{cusp-cover-extension-theorem}
\mbox{\Verified}\ 
After shrinking the $v$-disc, the following statements hold.

\begin{enumerate}
\item The maps in \eqref{cusp-cover-source-action} extend to a free,
properly discontinuous action on $Y^{(N)}$.  The quotient
\[
 f_N\colon N_0^{\langle N\rangle}
 :=Y^{(N)}/\bar\Lambda\longrightarrow\Delta_v
\]
is a proper holomorphic map from a smooth threefold.  Its central fibre
$W_N$ is reduced, irreducible, and nonnormal, and its normalization is a
toric del Pezzo surface of degree six.

\item Multiplication by $N$ on the cocharacter lattice defines a finite
flat toric morphism
\begin{equation}
 P_N\colon Y^{(N)}\longrightarrow Y,\qquad
 P_N(\xi_1,\xi_2,v)=(\xi_1^N,\xi_2^N,v^N),
 \label{cusp-cover-toric-power-map}
\end{equation}
of degree $N^3$.  It is equivariant for
\eqref{cusp-cover-source-action} and
\eqref{cusp-cover-target-action}, hence descends to a finite flat map
\begin{equation}
 \kappa_N\colon N_0^{\langle N\rangle}\longrightarrow N_0,
 \qquad f_0\circ\kappa_N=\rho_N\circ f_N,
 \label{cusp-cover-completed-map}
\end{equation}
of degree $N^3$.

\item The map $\kappa_N$ is etale over $v\ne0$, and
\begin{equation}
 \kappa_N^*W=NW_N,
 \qquad
 \operatorname{Ram}(\kappa_N)=(N-1)W_N.
 \label{cusp-cover-direct-ramification}
\end{equation}
In particular, the branch support of the map to the original cusp family
is the whole central fibre.

\item Let
\begin{equation}
 N_0[N]:=N_0\times_{\Delta_t}\Delta_v,
 \qquad t=v^N.
 \label{cusp-cover-base-changed-target}
\end{equation}
Then $N_0[N]$ is normal and $\kappa_N$ factors uniquely as
\begin{equation}
 N_0^{\langle N\rangle}
 \xrightarrow{\ Q_N\ }N_0[N]
 \longrightarrow N_0.
 \label{cusp-cover-relative-factorization}
\end{equation}
The map $Q_N$ has generic degree $N^2$.  For $N>1$ it is quasi-etale but
not flat: it is etale over the smooth locus of $N_0[N]$, and its non-etale
and nonflat loci are exactly the inverse images of the double and triple
strata of the base-changed central fibre.  Thus $Q_N$ is not a
normalization map.

\item On the reduced central fibres, the induced map
$W_N\to W$ has generic degree $N^2$.  It lifts to a commutative diagram
\begin{equation}
\begin{CD}
 E_N\simeq\mathrm{dP}_6 @>{[N]}>> E\simeq\mathrm{dP}_6\\
 @V{\nu_N}VV @VV{\nu}V\\
 W_N @>>> W,
\end{CD}
 \label{cusp-cover-normalization-diagram}
\end{equation}
where $[N]$ is toric multiplication by $N$, of degree $N^2$, ramified
with index $N$ along each of the six toric boundary curves.  The diagram
respects the paired conductor curves and the two triple-point classes.

\item Over $v\ne0$, $Q_N$ is exactly the cover defined by
\[
 \mathbb Z\langle\widehat\gamma,\widehat u,
 N\widehat w,N\widehat\delta\rangle
 \subset
 \mathbb Z\langle\widehat\gamma,\widehat u,
 \widehat w,\widehat\delta\rangle.
\]
Its deck group is $(\mathbb Z/N\mathbb Z)^2$.  The pulled-back cusp
monodromy is $M_0^N$, and in the displayed cover basis it is $M_0$:
\begin{equation}
 S_N^{-1}M_0^NS_N=M_0,\qquad
 S_N=\operatorname{diag}(1,1,N,N).
 \label{cusp-cover-monodromy-conjugacy}
\end{equation}
\end{enumerate}
\end{theorem}

\begin{proof}
We prove the assertions in the order in which the constructions occur.

\medskip\noindent
\emph{The source action.}
For $\bar\lambda\in\bar\Lambda$, the integral linear map
\[
 \varphi_{\bar\lambda}(y,y_3)
 =(y+y_3B_0\bar\lambda,y_3)
\]
translates the height-one slice by the integral vector
$B_0\bar\lambda$.  It preserves the $A_2$ triangulation and hence induces
a toric automorphism $\Phi_{\bar\lambda}$ of $Y^{(N)}$.  On the dense
torus this automorphism is
\[
 \Phi_{\bar\lambda}(\xi,v)
 =(v^{B_0\bar\lambda}\xi,v).
\]
Multiplication by the nowhere-zero holomorphic function
$c^{(N)}_{\bar\lambda}(v)$ extends across every toric chart.  Since
$C_N(v)\bar\lambda$ is additive in $\bar\lambda$, the resulting maps
satisfy
$\Psi^{(N)}_{\bar\lambda}\Psi^{(N)}_{\bar\mu}
=\Psi^{(N)}_{\bar\lambda+\bar\mu}$.

For completeness, freeness and proper discontinuity can be checked
directly.  Put
\[
 R_N(v):=-2\pi\operatorname{Im}C_N(v).
\]
It is bounded at $v=0$.  At a dense-torus point fixed by
$\bar\lambda$, logarithmic absolute values give
\[
 \log|v|B_0\bar\lambda+R_N(v)\bar\lambda=0.
\]
Choose the disc so small that
$|\log|v||>2\|R_N\|_\infty$.  Because $B_0$ is orthogonal, the last
equality is impossible for $\bar\lambda\ne0$.  On the central fibre,
$\Psi^{(N)}_{\bar\lambda}$ translates every bounded cell by the nonzero
vector $B_0\bar\lambda$, so it fixes no toric orbit.

For a dense-torus point set
\[
 y(\xi,v)=\frac{(\log|\xi_1|,\log|\xi_2|)}{\log|v|},
 \qquad
 B_{N,v}=B_0+\frac{R_N(v)}{\log|v|}.
\]
Then
\[
 y(\Psi^{(N)}_{\bar\lambda}(\xi,v))
 =y(\xi,v)+B_{N,v}\bar\lambda,
 \qquad
 |B_{N,v}\bar\lambda|\ge\tfrac12|\bar\lambda|.
\]
The logarithmic-position regions associated with any two fixed toric
charts are bounded.  Hence only finitely many $\bar\lambda$ can translate
one of them to meet the other.  Local finiteness of the fan reduces the
same assertion for compact sets to finitely many chart pairs, proving
proper discontinuity.  Finally, reducing
$B_{N,v}^{-1}y(\xi,v)$ modulo $\mathbb Z^2$ places every orbit meeting a
closed subdisc in finitely many closed unit toric charts.  Their union is
compact.  This proves properness of $f_N$.  The central-fibre and
normalization assertions follow exactly from the height-one fan: locally
$v$ is the product of the coordinates belonging to the rays through a
cell, every component occurs with multiplicity one, the lattice action is
transitive on vertices, and the star of a vertex is the complete smooth
hexagonal fan.

\medskip\noindent
\emph{The toric map and equivariance.}
Multiplication by $N$ on $N'$ sends every cone $\sigma\in\Sigma$ onto
itself as a real cone.  It therefore defines the toric morphism $P_N$.
For a maximal cone, let $z_0,z_1,z_2$ be the target characters dual to its
three primitive height-one ray generators, and let $w_0,w_1,w_2$ be the
corresponding source characters.  The map on convergent local rings is
\begin{equation}
 \mathbb C\{z_0,z_1,z_2\}\longrightarrow
\mathbb C\{w_0,w_1,w_2\},\qquad z_i\longmapsto w_i^N.
 \label{cusp-cover-maximal-chart-map}
\end{equation}
The target-chart preimage is the corresponding source chart, and the
right-hand ring is free of rank $N^3$ over the left-hand ring, with basis
\[
 \{w_0^{a_0}w_1^{a_1}w_2^{a_2}:0\le a_i<N\}.
\]
Thus $P_N$ is finite flat of degree $N^3$ on every maximal chart, hence
globally.

The choice \eqref{cusp-cover-source-factor} is forced by equivariance.
Indeed, writing powers coordinatewise,
\begin{align}
 P_N\Psi^{(N)}_{\bar\lambda}(\xi,v)
 &=\bigl((c^{(N)}_{\bar\lambda}(v))^N
          v^{NB_0\bar\lambda}\xi^N,v^N\bigr)\\
 &=\bigl(c_{\bar\lambda}(v^N)
          (v^N)^{B_0\bar\lambda}\xi^N,v^N\bigr)\\
 &=\Psi_{\bar\lambda}P_N(\xi,v).
 \label{cusp-cover-equivariance-computation}
\end{align}
Consequently $P_N$ descends to \eqref{cusp-cover-completed-map}.  To see
finiteness and flatness after the infinite-lattice quotient without any
implicit descent theorem, take an evenly covered open set in $N_0$ and
one of its lifts $V\subset Y$.  Equivariance gives a unique representative
of every source orbit above that open set whose image lies in $V$; two
such representatives cannot differ by a nontrivial lattice translate.
The descended map is therefore locally identified with
$P_N^{-1}(V)\to V$.  It is finite flat of degree $N^3$ by
\eqref{cusp-cover-maximal-chart-map}.

In the same chart the base characters are
\[
 t=z_0z_1z_2,\qquad v=w_0w_1w_2,\qquad P_N^*t=v^N.
\]
The Jacobian determinant of \eqref{cusp-cover-maximal-chart-map} is
\[
 N^3(w_0w_1w_2)^{N-1}=N^3v^{N-1}.
\]
This proves \eqref{cusp-cover-direct-ramification}; the local quotient
maps are biholomorphisms, so the calculation is unchanged downstairs.

\medskip\noindent
\emph{The relative degree-$N^2$ factor.}
The equality
\((f_0\circ\kappa_N)([\xi,v])=v^N=(\rho_N\circ f_N)([\xi,v])\)
gives the factorization
\eqref{cusp-cover-relative-factorization}.  At a triple point the exact
local-ring map is
\begin{equation}
 \begin{gathered}
 A_N:=
 \frac{\mathbb C\{z_0,z_1,z_2,v\}}
      {(z_0z_1z_2-v^N)},
 \qquad
 B:=\mathbb C\{w_0,w_1,w_2\},\\
 A_N\longrightarrow B,\qquad
 z_i\longmapsto w_i^N,\qquad
 v\longmapsto w_0w_1w_2.
 \end{gathered}
 \label{cusp-cover-basechange-local-ring}
\end{equation}
Let
\[
 G_N=\{(\varepsilon_0,\varepsilon_1,\varepsilon_2)\in\mu_N^3:
       \varepsilon_0\varepsilon_1\varepsilon_2=1\}
 \simeq\mu_N^2
\]
act diagonally on $B$.  A monomial
$w_0^{a_0}w_1^{a_1}w_2^{a_2}$ is $G_N$-invariant if and only if
\[
 a_0\equiv a_1\equiv a_2\pmod N.
\]
Writing their common residue as $r$, with $0\le r<N$, expresses that
monomial uniquely in the form
\[
 (w_0w_1w_2)^r
 (w_0^N)^{b_0}(w_1^N)^{b_1}(w_2^N)^{b_2}.
\]
The same coefficientwise argument applies to convergent power series.
Hence
\begin{equation}
 A_N=B^{G_N}.
 \label{cusp-cover-invariant-ring}
\end{equation}
It follows that $A_N$ is normal, that the map has generic degree
$|G_N|=N^2$, and that it is a finite quotient rather than a normalization.
The corresponding face-chart calculations prove normality along the
double and smooth strata as well, so $N_0[N]$ is normal everywhere.

The stabilizer in $G_N$ is trivial when none or exactly one of the $w_i$
vanishes, is isomorphic to $\mu_N$ when exactly two vanish, and is all of
$G_N$ when all three vanish.  Thus $Q_N$ is etale in codimension one and
has no ramification divisor, but it is not etale above the double and
triple strata.  It is also nonflat at every such point when $N>1$.  At a
triple point the closed fibre of \eqref{cusp-cover-basechange-local-ring}
has length
\begin{align}
 \dim_{\mathbb C}
 \frac{\mathbb C\{w_0,w_1,w_2\}}
      {(w_0^N,w_1^N,w_2^N,w_0w_1w_2)}
 &=N^3-(N-1)^3\\
 &=3N^2-3N+1,
 \label{cusp-cover-triple-fibre-length}
\end{align}
which exceeds the generic rank $N^2$.  At a point of a double stratum,
localizing at the nonzero third coordinate and choosing its holomorphic
$N$th root reduces the ramified factor to
\[
 \frac{\mathbb C\{z_0,z_1,v\}}{(z_0z_1-v^N)}
\longrightarrow\mathbb C\{w_0,w_1\},\qquad
 z_i\longmapsto w_i^N,\quad v\longmapsto w_0w_1.
\]
Its local generic rank is $N$, whereas its closed fibre has length
\begin{equation}
 \dim_{\mathbb C}
 \frac{\mathbb C\{w_0,w_1\}}
      {(w_0^N,w_1^N,w_0w_1)}=2N-1>N.
 \label{cusp-cover-double-fibre-length}
\end{equation}
The remaining transverse nonzero torus coordinate is etale locally.
Equations \eqref{cusp-cover-triple-fibre-length} and
\eqref{cusp-cover-double-fibre-length}, together with etaleness elsewhere,
identify the nonflat locus exactly.

\medskip\noindent
\emph{The regular fibres and monodromy.}
Choose a logarithm $r$ with $v=e^{2\pi ir}$.  The target period matrix at
$t=v^N$ is
\[
 Z(Nr)=NrB_0+C(v^N),
\]
whereas the source residual action is determined by
\[
 Z_N(r)=rB_0+\frac1NC(v^N)=\frac1N Z(Nr).
\]
The source fibre is therefore
\[
 \frac{\mathbb C^2}{\mathbb Z^2+Z_N(r)\mathbb Z^2},
\]
and multiplication by $N$ on $\mathbb C^2$ maps it to the target fibre.
Equivalently, in the unscaled additive coordinate it is the identity map
\begin{equation}
 \frac{\mathbb C^2}
      {N\mathbb Z^2+Z(Nr)\mathbb Z^2}
 \longrightarrow
 \frac{\mathbb C^2}
      {\mathbb Z^2+Z(Nr)\mathbb Z^2}.
 \label{cusp-cover-additive-fibre-map}
\end{equation}
The imaginary part of $Z(Nr)$ is invertible on a regular fibre, so its
period lattice meets $\mathbb Z^2$ trivially in the required real sense;
the index in \eqref{cusp-cover-additive-fibre-map} is $N^2$.  Its deck
group is $(\mathbb Z/N\mathbb Z)^2$.  In the ordered lattice basis
$(\widehat\gamma,\widehat u,\widehat w,\widehat\delta)$, this is exactly
the sublattice obtained by $S_N$.  Since $M_0=I+R$ with $R^2=0$ and
$R(a,b,c,d)=(0,0,b,-a)$, one has
\[
 M_0^N=I+NR,\qquad S_N^{-1}M_0^NS_N=M_0,
\]
proving \eqref{cusp-cover-monodromy-conjugacy}.

\medskip\noindent
\emph{The central normalization and conductor.}
Multiplication by $N$ fixes every ray of $\Sigma$.  It therefore maps a
chosen source component of the toric central fibre to the corresponding
target component.  The star of that ray is the hexagonal fan, so the
restriction is the toric power map
$[N]\colon\mathrm{dP}_6\to\mathrm{dP}_6$.  On its dense two-torus it has
degree $N^2$; in each toric boundary chart it is an $N$th-power map in the
normal coordinate, hence has ramification index $N$ along all six boundary
curves.  Equation \eqref{cusp-cover-equivariance-computation} says that
the source and target lattice identifications commute with this map.
Consequently opposite boundary curves, their conductor identifications,
and the two translation-orbits of triangle vertices descend compatibly.
This proves \eqref{cusp-cover-normalization-diagram} and completes the
proof.
\end{proof}

\begin{remark}[Why the unmodified action does not work]
\label{cusp-cover-naive-action-failure}
If one reuses the target action with $v$ in place of $t$ at the source,
then the two sides of equivariance have multiplicative factors
\[
 c_{\bar\lambda}(v)^N v^{NB_0\bar\lambda}
 \quad\hbox{and}\quad
 c_{\bar\lambda}(v^N)v^{NB_0\bar\lambda},
\]
respectively.  In general
$c_{\bar\lambda}(v)^N\ne c_{\bar\lambda}(v^N)$.  Thus the bare formula
$(\xi_1,\xi_2,v)\mapsto(\xi_1^N,\xi_2^N,v^N)$ extends on the ambient toric
variety but does not descend through the manuscript's quotient unless the
source factor \eqref{cusp-cover-source-factor} is inserted.  The explicit
logarithm $2\pi iC(v^N)\bar\lambda/N$ makes this root holomorphic; no
unjustified choice of an $N$th root is involved.
\end{remark}

\begin{remark}[Exact scope]
\label{cusp-cover-global-scope}
The theorem constructs a finite morphism of completed families over the
local cusp discs.  It is a degree-$N^3$ map to the original family and a
degree-$N^2$ map only after the target is base changed.  It is neither a
degree-$N^2$ self-map over the unchanged disc nor the normalization of the
base change.

Nothing here supplies a cover over the entire base $B$.  In fact the
displayed sublattice is not preserved by the original full monodromy when
$N>1$: for example, the printed matrix $A_1$ sends
$\widehat u$ to a vector whose $\widehat w$- and
$\widehat\delta$-coefficients are $-1$ and $1$.  A further global base
cover and compatible extensions through both finite fillings would have
to be constructed separately.  The local Kummer morphism also supplies
no map to a product sphere, smash sphere, $S^{12}$, or $S^{24}$.
\end{remark}


\paragraph{Context from the separate cone-target construction.}
The canonical workbench contains a different, codomain-changing construction
whose proof is not reproduced or claimed in this standalone paper.  In that
construction, collapsing the appropriate cusp boundary pieces gives based
maps
\[
 q_2^{\wedge2}:X^{\wedge2}\cong S^{12}
 \longrightarrow\mathcal A^{\wedge2},
 \qquad
 q_1^{\wedge4}:X^{\wedge4}\cong S^{24}
 \longrightarrow\mathcal C^{\wedge4}.
\]
Their synchronized radial target slices have intrinsic varying flat metrics
with pointed local-isometric $\R^4$ limits.  In every one fixed compatible
ambient metric on the compact cone targets, however, those same slices
collapse to the relevant cone point.  These sphere-source maps have different
codomains from the fixed $(\R/\Z)^4$ target in
Corollary~\ref{cor:product-smash-obstruction}, so they are outside that
literal nonextension statement.

\section{A gapless vertical direct-integral operator}

Choose a measurable cusp ray $T\mapsto z(T)$ for $T\ge T_0$ and normalize Haar measure on every fibre.  The marking identifies the fibre characters measurably.  Let
\begin{equation}\label{eq:direct-integral-space}
 \mathcal H_0=
 \int_{T_0}^{\infty}{}^{\oplus}
 L_0^2(F_{z(T)}^{\mathrm{flat}})\,w(T)\,\dd T,
\end{equation}
where $L_0^2$ denotes the orthogonal complement of the constants and $w$ is a positive locally integrable weight whose measure is nonzero on arbitrarily distant bounded intervals.  Define the decomposable operator
\begin{equation}\label{eq:direct-integral-operator}
 H_{\mathrm v}=
 \int_{T_0}^{\infty}{}^{\oplus}
 \Delta_{z(T)}\big|_{L_0^2}\,w(T)\,\dd T.
\end{equation}
Fourier expansion in the marked angular coordinates gives the concrete
unitary identification
\begin{equation}\label{eq:direct-integral-Fourier-sum}
 \mathcal H_0\cong
 \bigoplus_{k\in\Z^4\setminus\{0\}}
 L^2([T_0,\infty),w(T)\,\dd T),
\end{equation}
under which
\begin{equation}\label{eq:direct-integral-domain}
 \begin{split}
 (H_{\mathrm v}u)_k(T)&=\lambda_T(k)u_k(T),\\
 \operatorname{Dom}(H_{\mathrm v})
 &=\left\{u\in\mathcal H_0:
 \sum_{k\ne0}\int_{T_0}^{\infty}
 \lambda_T(k)^2|u_k(T)|^2w(T)\,\dd T<\infty\right\}.
 \end{split}
\end{equation}

\begin{theorem}[Continuous spectral threshold at zero]\label{thm:direct-integral}
The operator $H_{\mathrm v}$ is nonnegative and self-adjoint,
\begin{equation}\label{eq:direct-integral-conclusion}
 \inf\Spec(H_{\mathrm v})=0,
 \qquad
 \ker H_{\mathrm v}=\{0\}.
\end{equation}
Moreover,
\begin{equation}\label{eq:direct-integral-essential}
 0\in\sigma_{\mathrm{ess}}(H_{\mathrm v})
 \cap\sigma_{\mathrm{cont}}(H_{\mathrm v}),
\end{equation}
so there is no positive spectral gap.  If the fibrewise constants are
retained instead, the kernel contains the entire parameter-space factor and
is infinite-dimensional.
\end{theorem}

\begin{proof}
For each $k\ne0$, multiplication by the positive measurable real function
$\lambda_T(k)$ on $L^2([T_0,\infty),w(T)\,\dd T)$ is self-adjoint on its
maximal multiplication domain.  The orthogonal sum of these operators has
exactly the domain \eqref{eq:direct-integral-domain}, proving
self-adjointness and nonnegativity; this is also the standard direct-integral
construction \citep[Chapter XIII]{ReedSimonIV_1978}.  Let $e_{\gamma,T}$ be
the normalized fixed-fibre character corresponding to $k=e_\gamma$.
Proposition~\ref{prop:fixed-spectrum}, formula \eqref{eq:dual-exact}, and the
cusp estimates imply
\begin{equation}\label{eq:specific-mode-energy}
 \langle e_{\gamma,T},\Delta_{z(T)}e_{\gamma,T}\rangle
 \le \frac{C}{T^2}
\end{equation}
for large $T$.  Theorem~\ref{thm:global-real-quotient} additionally identifies
this already-defined character as the pullback of the first Fourier mode on
the global real circle quotient, but the direct-integral argument does not
depend on that identification and does not require it to be the exact shortest
fibre mode.  Choose pairwise disjoint intervals $I_j=[S_j,S_j+1]$ with $S_j\to\infty$ and positive $w(T)\dd T$ measure.  The sections
\[
 \Psi_j(T)=
 \frac{\mathbf 1_{I_j}(T)e_{\gamma,T}}
 {\left(\int_{I_j}w(T)\,\dd T\right)^{1/2}}
\]
belong to $\operatorname{Dom}(H_{\mathrm v})$, have unit norm, are mutually
orthogonal, hence converge weakly to zero, and satisfy
\[
 \|H_{\mathrm v}\Psi_j\|^2
 \le C^2/S_j^4\longrightarrow0.
\]
Weyl's criterion proves $0\in\sigma_{\mathrm{ess}}(H_{\mathrm v})$, and in
particular the spectral bottom is zero.

If $H_{\mathrm v}\Psi=0$, then nonnegativity and decomposability give
\[
 0=\int\langle\Psi(T),\Delta_{z(T)}\Psi(T)\rangle w(T)\,\dd T.
\]
The integrand is nonnegative, hence vanishes almost everywhere.  Every fibre vector then lies in the constants, but those were removed in \eqref{eq:direct-integral-space}; therefore $\Psi=0$.  A self-adjoint operator has no residual spectrum, so an injective operator with $0\in\Spec$ has $0$ in its continuous spectrum.  Retaining constants instead gives all sections $T\mapsto c(T)\mathbf 1_T$ as zero modes.
\end{proof}

\begin{remark}[Vertical, not total-space]
The operator \eqref{eq:direct-integral-operator} differentiates only along the torus fibres.  It has no derivative in the cusp parameter, no boundary coupling, and no claim to be the Laplacian of a smooth metric on the completed six-dimensional total space.  Its gaplessness is exact for the operator defined here, but it must not be silently transferred to a different operator.
\end{remark}

\section{Classical gauge fluctuation operators on the radial cusp}

Let $G$ be a compact Lie group with real Lie algebra $\mathfrak g$, equipped
with a positive-definite $\operatorname{Ad}(G)$-invariant inner product, and
put
\[
 d_G=\dim_{\R}\mathfrak g,\qquad
 \mathfrak g_{\C}=\mathfrak g\otimes_{\R}\C.
\]
Fix a coupling $g_{\mathrm{YM}}>0$.  Every operator in this section acts on
the trivial bundle over one of the auxiliary flat tori
$F_z^{\mathrm{flat}}$ along the positive radial cusp.  No physical spacetime
identification is made.

\subsection{Continuum Hessian and Coulomb spectrum}

On the trivial principal bundle
$\mathcal P_T=F_z^{\mathrm{flat}}\times G$, let $A_0$ be the product
connection and normalize
\begin{equation}\label{eq:YM-functional}
 \operatorname{YM}_T(A)=\frac1{2g_{\mathrm{YM}}^2}
 \int_{F_z^{\mathrm{flat}}}|F_A|^2\,\dd\operatorname{vol}_{F_z}.
\end{equation}

\begin{proposition}[Hessian at the product connection]
\label{prop:YM-Hessian}
For smooth $a,b\in\Omega^1(F_z^{\mathrm{flat}};\mathfrak g)$,
\begin{equation}\label{eq:YM-Hessian-bilinear}
 \operatorname{Hess}_{A_0}\operatorname{YM}_T(a,b)
 =g_{\mathrm{YM}}^{-2}\langle\dd a,\dd b\rangle_{L^2}.
\end{equation}
Let $\dd_1$ be the maximal closed exterior derivative on
$L^2\Omega^1(F_z^{\mathrm{flat}};\mathfrak g)$,
\[
 \operatorname{Dom}(\dd_1)=\{a\in L^2\Omega^1:
 \dd a\in L^2\Omega^2\text{ distributionally}\}.
\]
The closed Hessian form and its nonnegative self-adjoint operator are
\begin{align}
 \mathfrak l_T[a]&=g_{\mathrm{YM}}^{-2}\|\dd_1a\|_{L^2}^2,
 &\operatorname{Dom}(\mathfrak l_T)&=\operatorname{Dom}(\dd_1),
 \label{eq:YM-closed-form}\\
 \mathcal L_T&=g_{\mathrm{YM}}^{-2}\dd_1^*\dd_1,
 &\operatorname{Dom}(\mathcal L_T)
 &=\{a\in\operatorname{Dom}(\dd_1):
 \dd_1a\in\operatorname{Dom}(\dd_1^*)\}.
 \label{eq:YM-Hessian-operator}
\end{align}
Its differential expression on smooth one-forms is
$g_{\mathrm{YM}}^{-2}\dd^*\dd$.
\end{proposition}

\begin{proof}
With $F_{A_0+a}=\dd a+\tfrac12[a\wedge a]$,
\[
 F_{A_0+sa+tb}=s\,\dd a+t\,\dd b+O((|s|+|t|)^2).
\]
The mixed derivative of \eqref{eq:YM-functional} is therefore
\eqref{eq:YM-Hessian-bilinear}.  The maximal distributional exterior
derivative is closed and densely defined, so the first representation theorem
for closed nonnegative forms gives \eqref{eq:YM-closed-form} and
\eqref{eq:YM-Hessian-operator}, including the displayed operator domain.
\end{proof}

Gauge-longitudinal one-forms and harmonic one-forms are Hessian-zero
directions.  Remove both by defining
\begin{equation}\label{eq:continuum-Coulomb-space}
 \mathscr C_T=\{a\in L^2\Omega^1(F_z^{\mathrm{flat}};\mathfrak g):
 \dd^*a=0\text{ distributionally},\ 
 a\perp\mathcal H^1(F_z^{\mathrm{flat}};\mathfrak g)\}.
\end{equation}
This closed subspace reduces $\mathcal L_T$.  On it the form domain is
$H^1\Omega^1\cap\mathscr C_T$ and the operator domain is
$H^2\Omega^1\cap\mathscr C_T$: on the coclosed slice the graph norm of
$\dd_1$, together with the $L^2$ norm, is the Hodge graph norm, so the
claims follow from the Gaffney estimate and elliptic regularity on the compact
flat torus.

For $k\in\Z^4$ put
\[
 \xi_{T,k}=\mathsf P_z^{-\mathsf T}k,
 \qquad
 \chi_{T,k}([x])=e^{2\pi i\langle\xi_{T,k},x\rangle}.
\]
If $\varepsilon(\xi)$ denotes exterior multiplication by a covector and
$\iota_{\xi^\sharp}$ contraction by its Euclidean dual vector, direct
differentiation gives
\begin{align}
 \dd(\chi_{T,k}\alpha)
 &=2\pi i\chi_{T,k}\varepsilon(\xi_{T,k})\alpha,
 \label{eq:continuum-d-symbol}\\
 \dd^*(\chi_{T,k}\alpha)
 &=-2\pi i\chi_{T,k}\iota_{\xi_{T,k}^\sharp}\alpha,
 \label{eq:continuum-dstar-symbol}\\
 \dd^*\dd(\chi_{T,k}\alpha)
 &=4\pi^2\chi_{T,k}
 \bigl(|\xi_{T,k}|^2\alpha-
 \varepsilon(\xi_{T,k})\iota_{\xi_{T,k}^\sharp}\alpha\bigr).
 \label{eq:continuum-Hessian-symbol}
\end{align}

\begin{theorem}[Exact reduced continuum Hessian gap]
\label{thm:continuum-YM-gap}
For $k\ne0$, the Coulomb polarization space
\begin{equation}\label{eq:continuum-transverse-space}
 \mathcal T_{T,k}=\ker\iota_{\xi_{T,k}^\sharp}
 \subset(\R^4)^*_{\C}
\end{equation}
has complex dimension three.  On
$\mathcal T_{T,k}\otimes\mathfrak g_{\C}$ the Hessian eigenvalue is
\begin{equation}\label{eq:continuum-YM-mode}
 \lambda_T^{\mathrm{YM}}(k)
 =\frac{4\pi^2}{g_{\mathrm{YM}}^2}
 |\mathsf P_z^{-\mathsf T}k|^2
\end{equation}
with complex multiplicity $3d_G$ on that signed Fourier fibre.  If
$K_T(\lambda)=\{k\ne0:\lambda_T^{\mathrm{YM}}(k)=\lambda\}$, then the
total complex multiplicity, and the dimension of the real eigenspace, is
$3d_G\#K_T(\lambda)$.

For all sufficiently large $T$, the first eigenvalue on $\mathscr C_T$ is
\begin{equation}\label{eq:continuum-YM-gap-exact}
 \gamma_T=\frac{4\pi^2}{g_{\mathrm{YM}}^2}
 \min\left\{
 \frac{T^2+m^2}{\mathcal D_T^2},
 \frac{(T-B_R)^2+36m^2}{\mathcal D_T^2}
 \right\},
\end{equation}
and
\begin{equation}\label{eq:continuum-YM-gap-expansion}
 \boxed{\gamma_T=
 \frac{4\pi^2}{g_{\mathrm{YM}}^2T^2}
 \left(1+\frac{2\min\{b_{\mathrm I,0},0\}}T+O(T^{-2})\right).}
\end{equation}
If $I_T\subset\{\gamma,u\}$ records which of the two exact quantities in
\eqref{eq:continuum-YM-gap-exact} attain the minimum, the real multiplicity
of $\gamma_T$ is exactly $6d_G\#I_T$.
\end{theorem}

\begin{proof}
For $k\ne0$, invertibility of $\mathsf P_z$ gives $\xi_{T,k}\ne0$.
Contraction by $\xi_{T,k}^\sharp$ has rank one on complexified one-forms,
so \eqref{eq:continuum-transverse-space} has dimension three.  On its kernel
the second term of \eqref{eq:continuum-Hessian-symbol} vanishes, proving
\eqref{eq:continuum-YM-mode}.  Fourier completeness gives the whole reduced
spectrum.  Complex conjugation exchanges $k$ and $-k$; the reality condition
on their amplitudes shows that each conjugate pair contributes $6d_G$ real
dimensions, proving the multiplicity statement.  The zero Fourier mode is
precisely the harmonic sector already removed in
\eqref{eq:continuum-Coulomb-space}.

Minimizing \eqref{eq:continuum-YM-mode} is exactly the shortest-dual-vector
problem solved by Theorem~\ref{thm:cusp-gap}.  Equations
\eqref{eq:unit-mode-exact-lengths} and \eqref{eq:cusp-gap-result} prove the
exact gap, its expansion, and the bottom multiplicity.
\end{proof}

In lattice-coordinate coefficients $a_y$, the exact Coulomb projection at
mode $k$ is
\begin{equation}\label{eq:continuum-Coulomb-projection}
 k^{\mathsf T}(\mathsf P_z^{\mathsf T}\mathsf P_z)^{-1}a_y=0,
 \qquad
 \Pi_{T,k}a_y=a_y-k
 \frac{k^{\mathsf T}(\mathsf P_z^{\mathsf T}\mathsf P_z)^{-1}a_y}
 {k^{\mathsf T}(\mathsf P_z^{\mathsf T}\mathsf P_z)^{-1}k}.
\end{equation}
Thus Coulomb reduction uses the full physical covector metric, including all
cross terms; it is not obtained by deleting a coordinate component.

\subsection{Finite cubical cochains and simultaneous cusp--mesh limits}

Fix $N\ge2$ and write $\Lambda_N=(\Z/N\Z)^4$.  Its vertex $n$ represents
the point $\mathsf P_z(n/N)$ in the flat torus.  Let
$E=(\R^4)^*$ with its Euclidean inner product, let $\mathbf e_j$ be the
standard coordinate columns in lattice-coordinate space, and define the
physical coframe
\begin{equation}\label{eq:finite-physical-coframe}
 \vartheta_T^j=\mathsf P_z^{-\mathsf T}\mathbf e_j\in E,
 \qquad 1\le j\le4.
\end{equation}
For $0\le p\le4$, set
\begin{equation}\label{eq:finite-cochain-spaces}
 C_N^p(T;\mathfrak g_{\C})
 =\operatorname{Map}(\Lambda_N,\bigwedge^pE_{\C}\otimes\mathfrak g_{\C})
\end{equation}
and equip it with
\begin{equation}\label{eq:finite-inner-product}
 \langle a,b\rangle_{N,p,T}
 =\frac{\mathcal D_T}{N^4}\sum_{n\in\Lambda_N}
 \langle a(n),b(n)\rangle_{\wedge^pE\otimes\mathfrak g_{\C}}.
\end{equation}
This includes the exact Riemann-sum weight.  If
$G_T=\mathsf P_z^{\mathsf T}\mathsf P_z$, then
\begin{equation}\label{eq:finite-exterior-metric}
 \left\langle
 \vartheta_T^{i_1}\wedge\cdots\wedge\vartheta_T^{i_p},
 \vartheta_T^{j_1}\wedge\cdots\wedge\vartheta_T^{j_p}
 \right\rangle
 =\det\bigl((G_T^{-1})_{i_aj_b}\bigr)_{a,b=1}^p.
\end{equation}
Indeed,
$\langle\vartheta_T^i,\vartheta_T^j\rangle
=\mathbf e_i^{\mathsf T}\mathsf P_z^{-1}\mathsf P_z^{-\mathsf T}
\mathbf e_j=(G_T^{-1})_{ij}$, and the inner product induced on an exterior
power is the determinant of the matrix of pairings.

Let $(\mathsf T_ja)(n)=a(n+\mathbf e_j)$ and put
\begin{equation}\label{eq:finite-differences}
 D_{N,j}^+=N(\mathsf T_j-\id),\qquad
 (D_{N,j}^-a)(n)=N(a(n)-a(n-\mathbf e_j)).
\end{equation}
Define
\begin{equation}\label{eq:finite-exterior-derivative}
 \dd_Na=\sum_{j=1}^4\varepsilon(\vartheta_T^j)D_{N,j}^+a.
\end{equation}
This convention stores differential-form coefficients.  If $c_I(n)$ instead
stores the integral over a positively oriented $p$-cell, then
$a_I(n)=N^pc_I(n)$ carries the unscaled cellular coboundary
$\mathsf T_j-\id$, together with the inner product
\begin{equation}\label{eq:integrated-cochain-normalization}
 \langle c,c'\rangle^{\mathrm{int}}_{N,p,T}
 =\mathcal D_TN^{2p-4}\sum_{n\in\Lambda_N}
 \langle c(n),c'(n)\rangle,
\end{equation}
isometrically to \eqref{eq:finite-exterior-derivative} and
\eqref{eq:finite-inner-product}.  Thus both conventions give exactly the
same Hessian symbols and eigenvalues.  Explicitly,
$\mathcal D_TN^{2p-4}\sum|c|^2
=\mathcal D_TN^{-4}\sum|N^pc|^2$, and multiplying the unscaled coboundary
of $c$ by $N^{p+1}$ gives exactly $\dd_N(N^pc)$.
The differences commute while exterior multiplications anticommute, so
$\dd_N^2=0$.  Translation is unitary for
\eqref{eq:finite-inner-product}, and the adjoint of exterior multiplication
by $\vartheta_T^j$ is contraction by $(\vartheta_T^j)^\sharp$.  Hence
\begin{equation}\label{eq:finite-adjoint}
 \dd_N^*=-\sum_{j=1}^4
 \iota_{(\vartheta_T^j)^\sharp}D_{N,j}^-.
\end{equation}

For $k\in\Lambda_N$, put
\begin{align}
 \chi_{N,k}(n)&=e^{2\pi i k\cdot n/N},&
 \mathfrak d_N(k)_j&=N(e^{2\pi ik_j/N}-1),
 \label{eq:finite-coordinate-symbol}\\
 \zeta_{N,T}(k)&=\sum_j\mathfrak d_N(k)_j\vartheta_T^j
 =\mathsf P_z^{-\mathsf T}\mathfrak d_N(k).
 \label{eq:finite-physical-symbol}
\end{align}
Then, for every constant exterior amplitude $\alpha$,
\begin{align}
 \dd_N(\chi_{N,k}\alpha)
 &=\chi_{N,k}\varepsilon(\zeta_{N,T}(k))\alpha,\\
 \dd_N^*(\chi_{N,k}\alpha)
 &=\chi_{N,k}\iota_{\overline{\zeta_{N,T}(k)}^\sharp}\alpha,\\
 (\dd_N\dd_N^*+\dd_N^*\dd_N)(\chi_{N,k}\alpha)
 &=|\zeta_{N,T}(k)|^2\chi_{N,k}\alpha.
 \label{eq:finite-Hodge-symbol}
\end{align}
The last identity is the Hermitian exterior-algebra relation
$\iota_{\overline\zeta^\sharp}\varepsilon(\zeta)+
\varepsilon(\zeta)\iota_{\overline\zeta^\sharp}=|\zeta|^2\id$.
It also proves that the harmonic $p$-cochains are exactly the constant
cochains, of complex dimension $\binom4p d_G$.
Indeed, $\mathfrak d_N(k)=0$ precisely for $k=0$ in $\Lambda_N$, and
$\mathsf P_z$ is invertible, so the Hodge symbol vanishes precisely on the
zero Fourier mode.

Define the finite reduced Coulomb slice and Hessian by
\begin{align}
 \mathscr C_{N,T}
 &=\ker(\dd_N^*:C_N^1\to C_N^0)
 \cap(\ker\Delta_{N,T}^{(1)})^\perp,
 &\Delta_{N,T}^{(1)}&=\dd_N\dd_N^*+\dd_N^*\dd_N,
 \label{eq:finite-Coulomb-space}\\
 \mathcal L_{N,T}&=g_{\mathrm{YM}}^{-2}\dd_N^*\dd_N.
 \label{eq:finite-Hessian}
\end{align}
Equivalently, remove $k=0$ and at each $k\ne0$ restrict the one-form
amplitude to
$\ker\iota_{\overline{\zeta_{N,T}(k)}^\sharp}$, a complex
three-dimensional space.  The exact projection is
\begin{equation}\label{eq:finite-Coulomb-projection}
 \Pi_{N,T,k}\alpha=\alpha-\zeta_{N,T}(k)
 \frac{\iota_{\overline{\zeta_{N,T}(k)}^\sharp}\alpha}
 {|\zeta_{N,T}(k)|^2}.
\end{equation}

\begin{theorem}[Finite Coulomb spectrum and every mesh regime]
\label{thm:finite-YM-gap}
On the nonzero $k$-Fourier fibre of $\mathscr C_{N,T}$,
$\mathcal L_{N,T}$ has the single eigenvalue
\begin{equation}\label{eq:finite-YM-mode}
 \lambda_{N,T}^{\mathrm{YM}}(k)
 =g_{\mathrm{YM}}^{-2}
 |\mathsf P_z^{-\mathsf T}\mathfrak d_N(k)|^2
 =g_{\mathrm{YM}}^{-2}\mathfrak d_N(k)^*G_T^{-1}
 \mathfrak d_N(k),
\end{equation}
with complex multiplicity $3d_G$.  Put
\begin{equation}\label{eq:finite-rho}
 \rho_N=4N^2\sin^2(\pi/N).
\end{equation}
For all sufficiently large $T$, uniformly in $N\ge2$, the first eigenvalue is
\begin{align}
 \gamma_{N,T}
 &=\min\{\lambda_{N,T}^{\mathrm{YM}}(\mathbf e_\gamma),
          \lambda_{N,T}^{\mathrm{YM}}(\mathbf e_u)\}
 \label{eq:finite-gap-exact}\\
 &=\frac{\rho_N}{g_{\mathrm{YM}}^2T^2}
 \left(1+\frac{2\min\{b_{\mathrm I,0},0\}}T+O(T^{-2})\right).
 \label{eq:finite-gap-expansion}
\end{align}
If $I_{N,T}\subset\{\gamma,u\}$ records the minimizing coordinate modes
and $\eta_N=1$ for $N=2$, $\eta_N=2$ for $N>2$, its real multiplicity is
\begin{equation}\label{eq:finite-gap-multiplicity}
 3d_G\eta_N\#I_{N,T}.
\end{equation}
The following conclusions hold:
\begin{enumerate}[label=(\roman*),leftmargin=2.3em]
\item For fixed $N$, $T^2\gamma_{N,T}\to\rho_N/g_{\mathrm{YM}}^2$ and
$\gamma_{N,T}\to0$.
\item For every integer-valued $N(T)\ge2$, with no growth assumption, there
are constants $0<c<C<\infty$ independent of $N(T)$ such that
\begin{equation}\label{eq:finite-arbitrary-mesh}
 \frac c{T^2}\le\gamma_{N(T),T}\le\frac C{T^2}
\end{equation}
for all large $T$.
\item If $N(T)\to\infty$, then
\begin{equation}\label{eq:finite-growing-mesh}
 \gamma_{N(T),T}=\frac{4\pi^2}{g_{\mathrm{YM}}^2T^2}
 \bigl(1+O(T^{-1})+O(N(T)^{-2})\bigr).
\end{equation}
\item The largest physical elementary edge is
\begin{equation}\label{eq:finite-mesh-size}
 \max_j\frac{\|\mathsf P_z\mathbf e_j\|}{N(T)}
 =\frac{T}{N(T)}(1+O(T^{-1})).
\end{equation}
Thus the varying continuum torus is resolved if and only if
$N(T)/T\to\infty$.  This refinement condition is not needed for the
algebraic soft-mode conclusion in (ii).
\end{enumerate}
\end{theorem}

\begin{proof}
The symbol identities preceding the theorem and the exterior-algebra
identity show that on the transverse space
$\dd_N^*\dd_N=|\zeta_{N,T}(k)|^2\id$, proving
\eqref{eq:finite-YM-mode} and its multiplicity.  Let
$\sigma_N=2N\sin(\pi/N)$, so $\rho_N=\sigma_N^2$.  The least nonzero modulus
of a coordinate of $\mathfrak d_N(k)$ is $\sigma_N$, attained precisely at
frequency $\pm1$ (with the two signs equal when $N=2$).  A frequency supported
in the first two coordinates but not equal to
$\pm\mathbf e_\gamma$ or $\pm\mathbf e_u$ has squared coordinate norm at
least $2\sigma_N^2$.  Indeed, either both coordinates are nonzero, or the
single frequency is not $\pm1$; in the latter case
$\sin(2\pi/N)/\sin(\pi/N)=2\cos(\pi/N)\ge\sqrt2$ for $N\ge4$, and
$N=2,3$ are immediate.

On complex two-vectors $z=(z_1,z_2)$, the relevant soft block is
\[
 \frac1{\mathcal D_T}
 \begin{pmatrix}m&T-B_R\\-T&6m\end{pmatrix}z.
\]
After multiplication by $T$, its singular values are
$1+O(T^{-1})$, uniformly on complex vectors, by
\eqref{eq:MT-J}.  The factor $\sqrt2$ therefore separates every other
soft-plane frequency from the two coordinate frequencies uniformly in $N$
for large $T$.  For a general $z\in\C^4$, solving
$\mathsf P_z^{\mathsf T}x=z$ shows that the first and third components of
$x$ are $z_3,z_4$.  Thus a frequency with nonzero third or fourth coordinate
has image norm at least $\sigma_N$, whereas either coordinate candidate is
$O(\sigma_N/T)$.  This proves the exact minimizing list and
\eqref{eq:finite-gap-exact}; the signed-frequency reality condition proves
\eqref{eq:finite-gap-multiplicity}.  Multiplying
\eqref{eq:gamma-eigenvalue-correction} and
\eqref{eq:u-eigenvalue-correction} by $\rho_N/g_{\mathrm{YM}}^2$ proves
\eqref{eq:finite-gap-expansion}.

The elementary sine bounds give $16\le\rho_N\le4\pi^2$, proving (i) and
(ii); $\rho_N=4\pi^2(1+O(N^{-2}))$ proves (iii).  Finally, the first two
columns of $\mathsf P_z$ have length $T(1+O(T^{-1}))$, while the last two
have length one.  Division by $N(T)$ proves (iv).
\end{proof}

\begin{remark}[Classical linear scope]
For each fixed $T$, the reduced continuum operator is positively gapped; for
each fixed $(T,N)$, so is the reduced finite operator.  What fails is a
positive bound uniform as $T\to\infty$.  Before reduction, gauge-longitudinal
directions and constant harmonic one-forms are separate Hessian-zero
directions.  The spaces \eqref{eq:finite-cochain-spaces} are linear
Lie-algebra-valued cochains, not compact-group link variables, and
\eqref{eq:finite-Hessian} is not a Wilson plaquette action.  Neither operator
contains the nonlinear interaction, a gauge-invariant quantum measure,
renormalization, confinement, or a Hamiltonian reconstruction.
\end{remark}

\section{Full-torus integral magnetic flux}

The elliptic quotient studied in the next section is not required in order to
define a magnetic Laplacian on the full marked four-torus.  We now retain the
original lattice, metric, constants, and flux class and give the exact
self-adjoint operator.  The result is fibrewise along the marked cusp; it does
not assert that the chosen class is fixed by the full monodromy.

Use angular coordinates $u\in\R^4/\Z^4$, so that the Euclidean covering
coordinate is $x=\mathsf P_z u$ and the metric matrix is
\begin{equation}\label{eq:magnetic-Gram}
 G_T=\mathsf P_z^{\mathsf T}\mathsf P_z,
 \qquad \sqrt{\det G_T}=|\det\mathsf P_z|=\mathcal D_T.
\end{equation}
Fix an integral alternating form
\begin{equation}\label{eq:integral-flux}
 \mathsf E=\sum_{1\le j<k\le4}E_{jk}\,\dd u_j\wedge\dd u_k
 \in\bigwedge^2(\Z^4)^*,
 \qquad E_{jk}\in\Z,
\end{equation}
and denote its skew matrix by the same letter.  For $n\in\Z^4$, set
\begin{equation}\label{eq:magnetic-automorphy}
 J_n^{\mathsf E}(u)=
 \exp\left(-2\pi i\sum_{j<k}E_{jk}n_j u_k\right).
\end{equation}
Since $\sum_{j<k}E_{jk}m_jn_k\in\Z$,
\[
 J_m^{\mathsf E}(u+n)J_n^{\mathsf E}(u)=J_{m+n}^{\mathsf E}(u).
\]
Thus \eqref{eq:magnetic-automorphy} defines a Hermitian line bundle
$\mathcal L_{\mathsf E,T}$ on $F_z^{\mathrm{flat}}$.  The connection
\begin{equation}\label{eq:magnetic-connection}
 \nabla^{\mathsf E}=\dd+2\pi i
 \sum_{j<k}E_{jk}u_j\,\dd u_k
\end{equation}
descends because
$J_n^{-1}\dd J_n=-2\pi i\sum_{j<k}E_{jk}n_j\dd u_k$, and its curvature is
\begin{equation}\label{eq:magnetic-curvature}
 F_{\nabla^{\mathsf E}}=2\pi i\,\mathsf E.
\end{equation}

The Hilbert space is concretely
\begin{align}
 \mathcal H_{\mathsf E,T}
 =\{f\in L^2_{\mathrm{loc}}(\R^4):{}&
 f(u+n)=J_n^{\mathsf E}(u)f(u)
 \text{ for every }n\in\Z^4\},
 \label{eq:magnetic-Hilbert-space}\\
 \|f\|^2={}&\int_{[0,1]^4}|f(u)|^2\mathcal D_T\,\dd^4u.
 \nonumber
\end{align}
On smooth functions with the same automorphy, define
\begin{equation}\label{eq:magnetic-form}
 q_{\mathsf E,T}[f]=\int_{[0,1]^4}
 (G_T^{-1})^{jk}\overline{\nabla_j^{\mathsf E}f}
 \nabla_k^{\mathsf E}f\,\mathcal D_T\,\dd^4u.
\end{equation}

\begin{theorem}[Full-torus magnetic oscillator spectrum]
\label{thm:full-torus-magnetic-spectrum}
The form \eqref{eq:magnetic-form} is nonnegative and closable, with closed
form domain $H^1(F_z^{\mathrm{flat}};\mathcal L_{\mathsf E,T})$.  Its
Friedrichs operator is
\[
 \Delta_{\mathsf E,T}=(\nabla^{\mathsf E})^*\nabla^{\mathsf E},
 \qquad
 \operatorname{Dom}(\Delta_{\mathsf E,T})
 =H^2(F_z^{\mathrm{flat}};\mathcal L_{\mathsf E,T}),
\]
and has compact resolvent.

Assume $\mathsf E$ is nondegenerate.  Define $b_1,b_2>0$ by
\begin{equation}\label{eq:magnetic-frequency-definition}
 \Spec\!\left(G_T^{-1/2}(2\pi\mathsf E)G_T^{-1/2}\right)
 =\{ib_1,-ib_1,ib_2,-ib_2\}.
\end{equation}
Then, as a set with finite multiplicities,
\begin{equation}\label{eq:magnetic-oscillator-spectrum}
 \Spec(\Delta_{\mathsf E,T})
 =\{(2n_1+1)b_1+(2n_2+1)b_2:
 n_1,n_2\in\Z_{\ge0}\}.
\end{equation}
In particular, the spacing from the ground level to the first distinct
level is
\begin{equation}\label{eq:magnetic-level-gap}
 \operatorname{gap}(\Delta_{\mathsf E,T})=2\min\{b_1,b_2\}.
\end{equation}
No multiplicity formula is needed for, or asserted in, these identities.
\end{theorem}

\begin{proof}
In local unitary trivializations, the integrand in
\eqref{eq:magnetic-form} is the squared norm of a first-order unitary
covariant derivative.  A finite trivializing cover and a partition of unity
show that $\|f\|^2+q_{\mathsf E,T}[f]$ is equivalent to the usual local
$H^1$ norm.  Its completion is therefore
$H^1(F_z^{\mathrm{flat}};\mathcal L_{\mathsf E,T})$.  The representation
theorem for closed nonnegative forms gives the Friedrichs operator.  Its
principal symbol is $G_T^{-1}(\xi,\xi)$, so it is elliptic; elliptic regularity
on the compact boundaryless torus identifies its domain with $H^2$, and the
compact inclusion $H^1\hookrightarrow L^2$ gives compact resolvent.

The real skew-normal-form theorem gives a constant orthonormal frame
$v_1,\ldots,v_4$ on the Euclidean cover in which $2\pi\mathsf E$ is
\[
 b_1v^1\wedge v^2+b_2v^3\wedge v^4,
\]
after reversing $v^2$ or $v^4$ if necessary.  Let $\nabla_a$ be covariant
differentiation along $v_a$.  On the common smooth invariant core,
\[
 [\nabla_1,\nabla_2]=ib_1,
 \qquad [\nabla_3,\nabla_4]=ib_2,
\]
and all other commutators vanish.  Put
\[
 C_1=\nabla_1-i\nabla_2,\qquad
 C_2=\nabla_3-i\nabla_4,\qquad
 \Delta_j=-(\nabla_{2j-1}^2+\nabla_{2j}^2).
\]
Direct multiplication, retaining the commutator signs, gives
\begin{equation}\label{eq:magnetic-ladder-identities}
 C_j^*C_j=\Delta_j-b_j,
 \qquad C_jC_j^*=\Delta_j+b_j,
 \qquad [\Delta_{\mathsf E,T},C_j]=-2b_jC_j.
\end{equation}

Compact resolvent gives a smooth lowest eigenfunction $\psi_0\ne0$.  If
$C_j\psi_0\ne0$, the last identity in
\eqref{eq:magnetic-ladder-identities} would produce an eigenfunction at a
strictly lower eigenvalue.  Hence $C_1\psi_0=C_2\psi_0=0$, and the first
identity gives
\[
 \Delta_{\mathsf E,T}\psi_0=(b_1+b_2)\psi_0.
\]
Repeated application of $C_1^*$ and $C_2^*$ produces nonzero eigenfunctions
at all values displayed in \eqref{eq:magnetic-oscillator-spectrum}; their
nonvanishing follows from $C_jC_j^*=C_j^*C_j+2b_j$.  Conversely, for an
eigenfunction $\psi$ of eigenvalue $\lambda>b_1+b_2$,
\[
 (\lambda-b_1-b_2)\psi
 =C_1^*C_1\psi+C_2^*C_2\psi.
\]
Each nonzero $C_j\psi$ lies at the strictly smaller eigenvalue
$\lambda-2b_j$.  Induction over the finite set of eigenvalues below
$\lambda$ therefore expresses every eigenspace in raised ground states and
excludes all other spectral values.  Raising once in a least-frequency plane
proves \eqref{eq:magnetic-level-gap}.
\end{proof}

\begin{corollary}[Pfaffian product and cusp gap bound]
\label{cor:full-torus-magnetic-gap}
Let
\begin{equation}\label{eq:magnetic-pfaffian}
 d:=|\operatorname{Pf}(\mathsf E)|
 =|E_{12}E_{34}-E_{13}E_{24}+E_{14}E_{23}|>0.
\end{equation}
Then
\begin{equation}\label{eq:magnetic-frequency-product}
 b_1b_2=\frac{4\pi^2d}{\mathcal D_T},
\end{equation}
and
\begin{equation}\label{eq:full-torus-magnetic-gap-bound}
 \boxed{
 \operatorname{gap}(\Delta_{\mathsf E,T})
 \le4\pi\sqrt{\frac d{\mathcal D_T}}
 =\frac{4\pi\sqrt d}{T}\bigl(1+O(T^{-1})\bigr).}
\end{equation}
Thus the ground-to-first-level spacing tends to zero for every fixed
nondegenerate integral flux on the marked cusp fibres.
\end{corollary}

\begin{proof}
The determinant of the skew matrix in
\eqref{eq:magnetic-frequency-definition} is both $b_1^2b_2^2$ and
\[
 \frac{\det(2\pi\mathsf E)}{\det G_T}
 =\frac{(2\pi)^4\operatorname{Pf}(\mathsf E)^2}{\mathcal D_T^2}.
\]
Taking the positive square root proves
\eqref{eq:magnetic-frequency-product}.  Since
$\min(b_1,b_2)\le\sqrt{b_1b_2}$, Theorem~\ref{thm:full-torus-magnetic-spectrum}
gives the inequality.  Finally,
$\mathcal D_T=T^2(1+O(T^{-1}))$ by
\eqref{eq:volume-asymptotic}, which gives the final asymptotic.
\end{proof}

\begin{remark}[Scope of the magnetic result]
Theorem~\ref{thm:full-torus-magnetic-spectrum} concerns a scalar magnetic
Laplacian on a Hermitian line bundle over each marked flat four-torus.  It is
not a Yang--Mills mass-gap theorem, and it does not claim that
$[\mathsf E]\in H^2(F_z;\Z)$ is invariant under the monodromy of the global
family.  It does show that the failure of the particular elliptic quotient
below is not an obstruction to a fibrewise full-torus magnetic gap-closing
sequence.
\end{remark}

% Standalone notation bridge for the current workbench spectral modules.
% The root paper uses R for the logarithmic radial height and T=Im(tau).

\section{Canonical radial notation for the current operator modules}
\label{cusp-spec-section}

The earlier sections use \(R\) for the logarithmic cusp height and
\(T=\operatorname{Im}\tau\).  The current workbench operator modules use the
logarithmic height as their independent parameter.  In their standalone
adaptations below that parameter is consistently denoted by \(R\), so no
second meaning is assigned to \(T\).  Thus
\begin{equation}
 t_c=e^{-2\pi R},\qquad s=iR,\qquad R\ge R_0,
 \label{eq:current-radial-ray}
\end{equation}
and we set
\begin{equation}
 \begin{aligned}
  \mu_R&:=\mu(e^{-2\pi R})=\alpha_R+i m_R,\\
  h_R&:=h(e^{-2\pi R})=\eta_R+i H_R,\\
  b_R&:=b(e^{-2\pi R})=\xi_R+i B_R,\\
  c_R&:=\xi_R-\eta_R,\qquad
  L_R:=R+H_R,\qquad q_R:=L_R-B_R,\\
  D_R&:=L_Rq_R+6m_R^2.
 \end{aligned}
 \label{cusp-spec-radial-notation}
\end{equation}
All quantities on the right are real.  In particular
\(L_R=\operatorname{Im}\tau_R\),
\(-q_R=\operatorname{Im}\beta_R\), and \(D_R>0\) is the physical covolume,
whereas the determinant function \(D\) in \eqref{eq:D} equals
\(-D_R/L_R<0\).  Within the imported operator modules below, bare \(m\)
denotes the limiting constant, not the varying function in
\eqref{eq:D}:
\begin{equation}
 m:=\operatorname{Im}\mu(0)=m_0,\qquad
 H:=\operatorname{Im}h(0),\qquad
 B:=\operatorname{Im}b(0)=b_{\mathrm I,0}.
 \label{cusp-spec-limiting-imaginary-parts}
\end{equation}
Holomorphy at \(t_c=0\) gives, without identifying any of these three
constants,
\begin{equation}
 m_R=m+O(e^{-2\pi R}),\qquad
 H_R=H+O(e^{-2\pi R}),\qquad
 B_R=B+O(e^{-2\pi R}).
 \label{eq:current-radial-errors}
\end{equation}

We retain the ordered real coordinates
\begin{equation}
 (z_1,z_2)=(x_1+iy_1,x_2+iy_2)
 \longmapsto(x_1,y_1,x_2,y_2).
 \label{cusp-spec-real-coordinate-order}
\end{equation}
They carry the unrescaled Euclidean metric used throughout the operator
modules.

\begin{proposition}[Exact translation to canonical radial variables]
\label{cusp-spec-exact-dual-lattice}
In the ordered marking
\((\widehat\gamma,\widehat u,\widehat w,\widehat\delta)\), the real period
matrix is
\begin{equation}
 P_R=
 \begin{pmatrix}
  6\alpha_R&\eta_R&1&0\\
  6m_R&L_R&0&0\\
  c_R&\alpha_R&0&1\\
  -q_R&m_R&0&0
 \end{pmatrix}.
 \label{cusp-spec-real-period-matrix}
\end{equation}
Its determinant and positive covolume are
\begin{equation}
 \det P_R=-D_R,\qquad
 \operatorname{covol}(P_R\mathbb Z^4)=D_R>0,
 \qquad
 \operatorname{Vol}(F_R)=D_R,
 \label{cusp-spec-exact-volume}
\end{equation}
where \(F_R=\mathbb R^4/P_R\mathbb Z^4\).  The complete inverse transpose is
\begin{equation}
 P_R^{-\mathsf T}=
 \begin{pmatrix}
 0&0&1&0\\[2pt]
 \dfrac{m_R}{D_R}&\dfrac{q_R}{D_R}&
 \dfrac{-6m_R\alpha_R-q_R\eta_R}{D_R}&
 \dfrac{-m_Rc_R-q_R\alpha_R}{D_R}\\[7pt]
 0&0&0&1\\[2pt]
 -\dfrac{L_R}{D_R}&\dfrac{6m_R}{D_R}&
 \dfrac{6(L_R\alpha_R-m_R\eta_R)}{D_R}&
 \dfrac{L_Rc_R-6m_R\alpha_R}{D_R}
 \end{pmatrix}.
 \label{cusp-spec-full-inverse-transpose}
\end{equation}
For \(k=(r,s,n_1,n_2)^{\mathsf T}\in\mathbb Z^4\), define
\begin{equation}
 \mathcal R_R(k)=r-6\alpha_Rn_1-c_Rn_2,
 \qquad
 S_R(k)=s-\eta_Rn_1-\alpha_Rn_2.
 \label{eq:current-shifted-indices}
\end{equation}
Then
\begin{equation}
 P_R^{-\mathsf T}k=
 \left(n_1,
 \frac{m_R\mathcal R_R(k)+q_RS_R(k)}{D_R},n_2,
 \frac{-L_R\mathcal R_R(k)+6m_RS_R(k)}{D_R}\right).
 \label{cusp-spec-exact-dual-lattice-formula}
\end{equation}
The scalar Laplacian defined by the \(H^1\) Dirichlet form has operator
domain \(H^2(F_R)\), compact resolvent, and eigenvalues
\begin{equation}
 \lambda_R(k)=4\pi^2\lVert P_R^{-\mathsf T}k\rVert^2,
 \qquad k\in\mathbb Z^4,
 \label{eq:current-scalar-spectrum}
\end{equation}
with Fourier multiplicity.

Finally,
\begin{equation}
 D_R=R^2+(2H-B)R+H(H-B)+6m^2+O(Re^{-2\pi R}).
 \label{cusp-spec-D-expansion}
\end{equation}
\end{proposition}

\begin{proof}
The identity
\[
 Z=s\begin{pmatrix}0&1\\-1&0\end{pmatrix}
 +\begin{pmatrix}6\mu&h\\b-h&\mu\end{pmatrix},
\]
which follows from \eqref{eq:cusp-s-coordinate} and \(b=\beta+\tau\),
gives \eqref{cusp-spec-real-period-matrix} by taking real and imaginary
parts in the fixed coordinate order.  Expansion along its last two columns
gives
\[
 \det P_R=-L_Rq_R-6m_R^2=-D_R.
\]
The sign \(D_R>0\) is exactly the earlier condition \(D<0\), since
\(D=\operatorname{Im}\beta_R-6m_R^2/L_R=-D_R/L_R\) and \(L_R>0\).

Multiplying the displayed candidate inverse transpose by \(P_R^{\mathsf T}\)
gives the identity matrix entry by entry: the two nontrivial \(2\times2\)
minors reduce to \(L_Rq_R+6m_R^2=D_R\), and the real-part terms cancel by
the definitions of \(c_R\) and \(S_R(k)\).  This proves
\eqref{cusp-spec-full-inverse-transpose}, and multiplication by \(k\)
gives \eqref{cusp-spec-exact-dual-lattice-formula}.

The characters
\(x\mapsto\exp(2\pi i\langle P_R^{-\mathsf T}k,x\rangle)\) form the complete
Fourier basis of the compact quotient.  Applying the nonnegative Euclidean
Laplacian gives \eqref{eq:current-scalar-spectrum}; closed-form theory and
elliptic regularity give the stated form and operator domains, and Rellich
compactness gives compact resolvent.  Finally, substituting
\(L_R=R+H_R\), \(q_R=R+H_R-B_R\), and the three holomorphic expansions into
\(D_R=L_Rq_R+6m_R^2\) gives the last formula with every constant and sign
retained.
\end{proof}

% Standalone-preprint adaptation of cusp_compact_link_gaussian.tex.
% Canonical source SHA-256: 4469781B2386225453DBDD5C12B5B5D9213CD67480055F3F318143941C44FA52.
% The canonical radial-height symbol T was renamed R consistently; transpose
% notation was retained exactly. Bibliography keys were aligned with this preprint.

\section{Cubical holonomy and finite gauge systems}
\label{cusp-link-section}

This section records three constructions on the regular cusp fibres while
keeping two independent integers separate throughout:
\(M_{\mathrm{cov}}\) is the index used to enlarge two periods of a fibre,
whereas \(N_{\mathrm{mesh}}\) is the number of cubical subdivisions in each
coordinate direction.  In particular, no mesh refinement is identified with
a change of covering space.

Let
\begin{equation}
 P_R=
 \begin{pmatrix}
  6\operatorname{Re}\mu_R&\operatorname{Re}\tau_R&1&0\\
  6\operatorname{Im}\mu_R&\operatorname{Im}\tau_R&0&0\\
  \operatorname{Re}\beta_R&\operatorname{Re}\mu_R&0&1\\
  \operatorname{Im}\beta_R&\operatorname{Im}\mu_R&0&0
 \end{pmatrix},
 \qquad
 F_R=\mathbb R^4/P_R\mathbb Z^4.
 \label{cusp-link-period-matrix}
\end{equation}
Write
\begin{equation}
 L_R=\operatorname{Im}\tau_R,
 \qquad m_R=\operatorname{Im}\mu_R,
 \qquad q_R=-\operatorname{Im}\beta_R,
 \qquad D_R=L_Rq_R+6m_R^2>0.
 \label{cusp-link-radial-data}
\end{equation}
Expanding the determinant along the last two columns gives
\[
 \det P_R=L_R\operatorname{Im}\beta_R-6m_R^2=-D_R.
\]
Along the positive radial cusp the source data give
\begin{equation}
 L_R=R+O(1),\qquad q_R=R+O(1),\qquad m_R=O(1),
 \label{cusp-link-cusp-asymptotics}
\end{equation}
and all three real parts occurring in \eqref{cusp-link-period-matrix} are
bounded.  Equivalently, after ordering physical coordinates as
\((x_1,x_3;x_2,x_4)\), the first two columns have block form
\begin{equation}
\binom{\mathsf A_R^{\mathrm{blk}}}{\mathsf B_R^{\mathrm{blk}}},\qquad
\mathsf B_R^{\mathrm{blk}}=
\begin{pmatrix}6m_R&L_R\\-q_R&m_R\end{pmatrix}
=R\begin{pmatrix}0&1\\-1&0\end{pmatrix}
 +\mathsf C_R^{\mathrm{blk}},
\qquad \sup_R\lVert \mathsf C_R^{\mathrm{blk}}\rVert_{\mathrm{op}}<\infty.
 \label{cusp-link-block-form}
\end{equation}

For \(M_{\mathrm{cov}}\geq1\), put
\begin{equation}
 S_{M_{\mathrm{cov}}}
 =\operatorname{diag}(1,1,M_{\mathrm{cov}},M_{\mathrm{cov}}),
 \qquad
 P_{R,M_{\mathrm{cov}}}=P_RS_{M_{\mathrm{cov}}},
 \label{cusp-link-cover-matrix}
\end{equation}
and define the degree-\(M_{\mathrm{cov}}^2\) cover
\begin{equation}
 \widetilde F_{R,M_{\mathrm{cov}}}
 =\mathbb R^4/P_{R,M_{\mathrm{cov}}}\mathbb Z^4
 \longrightarrow F_R.
 \label{cusp-link-cover-torus}
\end{equation}
Indeed,
\(S_{M_{\mathrm{cov}}}\mathbb Z^4\subset\mathbb Z^4\) and
\[
 \mathbb Z^4/S_{M_{\mathrm{cov}}}\mathbb Z^4
 \cong(\mathbb Z/M_{\mathrm{cov}}\mathbb Z)^2,
\]
so the identity on \(\mathbb R^4\) induces the asserted cover and its deck
group has order \(M_{\mathrm{cov}}^2\).
Its Euclidean quotient metric, written on the coordinate torus
\(\mathbb R^4/\mathbb Z^4\), has Gram matrix and covolume
\begin{equation}
 G_{R,M_{\mathrm{cov}}}
 =P_{R,M_{\mathrm{cov}}}^{\mathsf T}P_{R,M_{\mathrm{cov}}},
 \qquad
 \mathcal D_{R,M_{\mathrm{cov}}}
 =\sqrt{\det G_{R,M_{\mathrm{cov}}}}
 =M_{\mathrm{cov}}^2D_R.
 \label{cusp-link-cover-metric}
\end{equation}

\subsection{The transported cubical subdivision and continuum holonomy}

Fix an integer \(N_{\mathrm{mesh}}\geq2\), and let
\(K_{N_{\mathrm{mesh}}}\) be the standard cubical subdivision of
\(\mathbb R^4/\mathbb Z^4\) obtained by dividing every coordinate circle
into \(N_{\mathrm{mesh}}\) equal intervals.  Its vertex set and positively
oriented coordinate edges are
\begin{equation}
 V_{N_{\mathrm{mesh}}}
 =(\mathbb Z/N_{\mathrm{mesh}}\mathbb Z)^4,
 \qquad
 E_{N_{\mathrm{mesh}}}
 =V_{N_{\mathrm{mesh}}}\times\{1,2,3,4\}.
 \label{cusp-link-cubical-sets}
\end{equation}

\begin{proposition}[Exact cubical realization and the torsion vertices]
\label{cusp-link-cubical-realization}
The linear map \(P_{R,M_{\mathrm{cov}}}\) induces a diffeomorphism
\begin{equation}
 \Phi_{R,M_{\mathrm{cov}}}:\mathbb R^4/\mathbb Z^4
 \longrightarrow\widetilde F_{R,M_{\mathrm{cov}}},
 \qquad [y]\longmapsto[P_{R,M_{\mathrm{cov}}}y].
 \label{cusp-link-coordinate-diffeomorphism}
\end{equation}
Transporting \(K_{N_{\mathrm{mesh}}}\) by this diffeomorphism gives an
honest cubical cell decomposition of the whole torus.  On vertices its
restriction is the bijection
\begin{equation}
 \begin{aligned}
 \iota_{R,M_{\mathrm{cov}},N_{\mathrm{mesh}}}:\quad
 &(\mathbb Z/N_{\mathrm{mesh}}\mathbb Z)^4
 \xrightarrow{\ \cong\ }
 \widetilde F_{R,M_{\mathrm{cov}}}[N_{\mathrm{mesh}}],\\
 &[n]\longmapsto
 \left[\frac1{N_{\mathrm{mesh}}}P_{R,M_{\mathrm{cov}}}n\right].
 \end{aligned}
 \label{cusp-link-torsion-bijection}
\end{equation}
Here the codomain is the finite \(N_{\mathrm{mesh}}\)-torsion subgroup,
not the entire torus.  The coordinate edge \((n,j)\) is the path
\begin{equation}
 \gamma_{n,j}(t)=
 \left[
  \frac1{N_{\mathrm{mesh}}}
  P_{R,M_{\mathrm{cov}}}(n+t\mathbf e_j)
 \right],
 \qquad 0\leq t\leq1.
 \label{cusp-link-coordinate-edge}
\end{equation}
More generally, for an ordered subset
\(I=\{i_1,\ldots,i_p\}\subset\{1,2,3,4\}\), the corresponding closed
\(p\)-cell is parametrized by
\begin{equation}
 c_{n,I}(t_1,\ldots,t_p)=
 \left[
  \frac1{N_{\mathrm{mesh}}}P_{R,M_{\mathrm{cov}}}
  \left(n+\sum_{a=1}^pt_a\mathbf e_{i_a}\right)
 \right],
 \qquad (t_1,\ldots,t_p)\in[0,1]^p.
 \label{cusp-link-general-cell}
\end{equation}
\end{proposition}

\begin{proof}
If \(y-y'\in\mathbb Z^4\), then
\(P_{R,M_{\mathrm{cov}}}(y-y')\) belongs to the quotient lattice, so
\eqref{cusp-link-coordinate-diffeomorphism} is well defined.  The matrix
\(P_{R,M_{\mathrm{cov}}}\) is invertible, and its inverse induces the
inverse map on the two quotients.  Hence it is a diffeomorphism and carries
every open cell of the standard subdivision diffeomorphically onto an open
cell of the target.

Every point displayed in \eqref{cusp-link-torsion-bijection} is killed by
\(N_{\mathrm{mesh}}\).  If it is zero, then
\[
 \frac{P_{R,M_{\mathrm{cov}}}n}{N_{\mathrm{mesh}}}
 \in P_{R,M_{\mathrm{cov}}}\mathbb Z^4.
\]
Invertibility gives
\(n\in N_{\mathrm{mesh}}\mathbb Z^4\), proving injectivity.  Conversely,
if \([x]\) is killed by \(N_{\mathrm{mesh}}\), choose a lift
\(x\in\mathbb R^4\).  There is \(n\in\mathbb Z^4\) such that
\(N_{\mathrm{mesh}}x=P_{R,M_{\mathrm{cov}}}n\), and therefore
\([x]=[P_{R,M_{\mathrm{cov}}}n/N_{\mathrm{mesh}}]\).  This proves
surjectivity.  Formula \eqref{cusp-link-coordinate-edge} is simply the
transport of the corresponding coordinate edge; at a wrap-around endpoint
the extra vector \(P_{R,M_{\mathrm{cov}}}\mathbf e_j\) is a period.
\end{proof}

Let \(G\) be a compact matrix Lie group and use the trivial principal bundle over
\(\widetilde F_{R,M_{\mathrm{cov}}}\).  A connection is represented by
\(A\in\Omega^1(\widetilde F_{R,M_{\mathrm{cov}}};\mathfrak g)\).  For an
oriented path \(\gamma\), define \(U_A(\gamma)\) by
\begin{equation}
 \frac{dU}{dt}=-A_{\gamma(t)}(\dot\gamma(t))U(t),
 \qquad U(0)=1,
 \qquad U_A(\gamma)=U(1).
 \label{cusp-link-parallel-transport-ode}
\end{equation}

\begin{proposition}[Typed continuum-to-link holonomy map]
\label{cusp-link-holonomy-map}
The formula
\begin{equation}
 \operatorname{Hol}_{R,M_{\mathrm{cov}},N_{\mathrm{mesh}}}(A)_{n,j}
 =U_A(\gamma_{n,j})
 \label{cusp-link-holonomy-definition}
\end{equation}
defines a map
\begin{equation}
 \operatorname{Hol}_{R,M_{\mathrm{cov}},N_{\mathrm{mesh}}}:
 \Omega^1(\widetilde F_{R,M_{\mathrm{cov}}};\mathfrak g)
 \longrightarrow G^{E_{N_{\mathrm{mesh}}}}.
 \label{cusp-link-holonomy-typed}
\end{equation}
For the convention
\begin{equation}
 A^g=g^{-1}Ag+g^{-1}dg,
 \qquad
 g\in C^\infty(\widetilde F_{R,M_{\mathrm{cov}}},G),
 \label{cusp-link-gauge-convention}
\end{equation}
it is equivariant for evaluation of \(g\) at the vertices:
\begin{equation}
 \operatorname{Hol}(A^g)_{n,j}
 =g(\gamma_{n,j}(1))^{-1}
  \operatorname{Hol}(A)_{n,j}
  g(\gamma_{n,j}(0)).
 \label{cusp-link-holonomy-equivariance}
\end{equation}
\end{proposition}

\begin{proof}
The smooth matrix ODE \eqref{cusp-link-parallel-transport-ode} has a unique
solution on every edge, which proves that \eqref{cusp-link-holonomy-typed}
is defined.  Let \(U\) solve that ODE for \(A\), set
\(g_t=g(\gamma_{n,j}(t))\), and put
\(V(t)=g_t^{-1}U(t)g_0\).  Direct differentiation gives
\[
 V'(t)=-\bigl(g_t^{-1}A(\dot\gamma_{n,j})g_t
                 +g_t^{-1}g_t'\bigr)V(t),
 \qquad V(0)=1.
\]
This is exactly the parallel-transport equation for \(A^g\).  Uniqueness
and evaluation at \(t=1\) prove \eqref{cusp-link-holonomy-equivariance}.
Notice that the endpoint inverse in this formula is forced by the joint
conventions \eqref{cusp-link-parallel-transport-ode} and
\eqref{cusp-link-gauge-convention}.
\end{proof}

\subsection{An explicit compact-link
\texorpdfstring{\(\mathrm{SU}(2)\)}{SU(2)} flux}

The ordered plaquette convention used below is most naturally written for
the inverse-transport link variable
\begin{equation}
 \mathcal U(A)_{n,j}
 =\operatorname{Hol}_{R,M_{\mathrm{cov}},N_{\mathrm{mesh}}}(A)_{n,j}^{-1}.
 \label{cusp-link-inverse-transport-variable}
\end{equation}
Equation \eqref{cusp-link-holonomy-equivariance} then becomes
\begin{equation}
 \mathcal U(A^g)_{n,j}
 =g(\gamma_{n,j}(0))^{-1}\mathcal U(A)_{n,j}g(\gamma_{n,j}(1)).
 \label{cusp-link-inverse-transport-gauge-law}
\end{equation}
Thus the product in \eqref{cusp-link-plaquette-convention} transforms by
conjugation at its base vertex.  The symbols \(U_j(n)\) in the remainder of
this subsection use this inverse-transport convention.

Use
\begin{equation}
 \mathfrak{su}(2)=\{X\in M_2(\mathbb C):X^*=-X,\ \operatorname{tr}X=0\},
 \qquad
 \langle X,Y\rangle_{\mathfrak{su}(2)}
 =-\frac12\operatorname{tr}(XY),
 \label{cusp-link-su2-inner-product}
\end{equation}
on real elements, and fix
\begin{equation}
 H_{\mathrm{link}}=\begin{pmatrix}i&0\\0&-i\end{pmatrix}.
 \qquad
 \lVert H_{\mathrm{link}}\rVert^2=1,
 \qquad \exp(2\pi H_{\mathrm{link}})=1.
 \label{cusp-link-flux-generator}
\end{equation}
Write \(N=N_{\mathrm{mesh}}\) temporarily inside the next displayed link
formulas, and represent \(n_1,n_2\) by integers in
\(\{0,\ldots,N-1\}\).

\begin{theorem}[A periodic compact-link flux and its exact action]
\label{cusp-link-flux-theorem}
On the oriented edges of \(K_N\), define
\begin{align}
 U_1(n)&=\exp\left(-\frac{2\pi n_2}{N^2}H_{\mathrm{link}}\right),
 \label{cusp-link-flux-u1}\\
 U_2(n)&=
 \begin{cases}
  1,&0\leq n_2\leq N-2,\\[1mm]
  \exp(2\pi n_1H_{\mathrm{link}}/N),&n_2=N-1,
 \end{cases}
 \label{cusp-link-flux-u2}\\
 U_3(n)&=U_4(n)=1.
 \label{cusp-link-flux-u34}
\end{align}
For
\begin{equation}
 U_{ij}(n)=U_i(n)U_j(n+\mathbf e_i)
            U_i(n+\mathbf e_j)^{-1}U_j(n)^{-1},
 \label{cusp-link-plaquette-convention}
\end{equation}
one has
\begin{equation}
 \begin{aligned}
 U_{12}(n)&=\exp(2\pi H_{\mathrm{link}}/N^2)
 &&\text{for every }n,\\
 U_{ij}(n)&=1
 &&\text{for }1\leq i<j\leq4,\quad(i,j)\neq(1,2).
 \end{aligned}
 \label{cusp-link-plaquette-values}
\end{equation}

Let \(\operatorname{Log}\) be the principal logarithm on the displayed
plaquette.  It is unambiguous because \(0<2\pi/N^2\leq\pi/2\).  Define
the logarithmic curvature two-form and its metric-weighted action by
\begin{align}
 \mathcal F_N^{\log}(n)
 &=N^2\operatorname{Log}(U_{12}(n))\,dy^1\wedge dy^2,
 \label{cusp-link-log-curvature}\\
 S^{\log}_{R,M_{\mathrm{cov}},N}(U)
 &=\frac{\mathcal D_{R,M_{\mathrm{cov}}}}{2g_{\mathrm{YM}}^2N^4}
   \sum_{n\in V_N}
   \lVert\mathcal F_N^{\log}(n)\rVert^2_{
       G_{R,M_{\mathrm{cov}}}^{-1}\otimes\mathfrak{su}(2)}.
 \label{cusp-link-log-action}
\end{align}
Then
\begin{equation}
 \mathcal F_N^{\log}(n)=2\pi H_{\mathrm{link}}\,dy^1\wedge dy^2,
 \qquad
 S^{\log}_{R,M_{\mathrm{cov}},N}(U)
 =\frac{2\pi^2M_{\mathrm{cov}}^2}{g_{\mathrm{YM}}^2D_R}.
 \label{cusp-link-log-action-value}
\end{equation}
In particular, the native-fibre value \(M_{\mathrm{cov}}=1\) is
\(2\pi^2/(g_{\mathrm{YM}}^2D_R)\), independently of the mesh.

There is no canonical scalar Wilson action for a general nonorthogonal
metric obtained merely by assigning independent character weights to all
coordinate plaquettes.  For this configuration, however, the exact
continuum-matched \((1,2)\)-sector Wilson functional is unambiguous if it is
defined explicitly by the following formula.  Its ordered plaquette
holonomy and real trace are the Wilson lattice-action construction
\cite{Wilson_1974}; the metric coefficient and normalization here are the
present definition, and their consequences are proved below rather than
imported from that source:
\begin{equation}
 S^{W,12}_{R,M_{\mathrm{cov}},N}(U)
 =\frac{\mathcal D_{R,M_{\mathrm{cov}}}\kappa^{12}_{R,M_{\mathrm{cov}}}}
        {2g_{\mathrm{YM}}^2}
   \sum_{n\in V_N}
   \left(2-\operatorname{Re}\operatorname{tr}U_{12}(n)\right),
 \label{cusp-link-wilson-sector-definition}
\end{equation}
where
\begin{equation}
 \kappa^{12}_{R,M_{\mathrm{cov}}}
 =(G_{R,M_{\mathrm{cov}}}^{-1})_{11}
  (G_{R,M_{\mathrm{cov}}}^{-1})_{22}
 -(G_{R,M_{\mathrm{cov}}}^{-1})_{12}^2
 =D_R^{-2}.
 \label{cusp-link-two-form-minor}
\end{equation}
Its exact value and continuum-mesh limit are
\begin{align}
 S^{W,12}_{R,M_{\mathrm{cov}},N}(U)
 &=\frac{M_{\mathrm{cov}}^2N^4}{2g_{\mathrm{YM}}^2D_R}
   \left(2-2\cos\frac{2\pi}{N^2}\right),
 \label{cusp-link-wilson-sector-value}\\
 &=\frac{M_{\mathrm{cov}}^2}{g_{\mathrm{YM}}^2D_R}
   \left(
    2\pi^2-\frac{2\pi^4}{3N^4}
    +\frac{4\pi^6}{45N^8}+O(N^{-12})
   \right),
 \label{cusp-link-wilson-sector-expansion}\\
 \lim_{N\to\infty}S^{W,12}_{R,M_{\mathrm{cov}},N}(U)
 &=\frac{2\pi^2M_{\mathrm{cov}}^2}{g_{\mathrm{YM}}^2D_R}.
 \label{cusp-link-wilson-sector-limit}
\end{align}
\end{theorem}

\begin{proof}
All links lie in the same diagonal circle subgroup, so they commute.  If
\(0\leq n_2\leq N-2\), the two \(U_2\)-terms in
\eqref{cusp-link-plaquette-convention} are trivial and
\[
 U_{12}(n)=
 \exp(-2\pi n_2H_{\mathrm{link}}/N^2)
 \exp(2\pi(n_2+1)H_{\mathrm{link}}/N^2)
 =\exp(2\pi H_{\mathrm{link}}/N^2).
\]
If \(n_2=N-1\), then \(U_1(n+\mathbf e_2)=1\).  Moreover,
\[
 U_2(n+\mathbf e_1)U_2(n)^{-1}=\exp(2\pi H_{\mathrm{link}}/N).
\]
For \(n_1=N-1\), the exponent computed with representatives differs from
\(2\pi H_{\mathrm{link}}/N\) by \(-2\pi H_{\mathrm{link}}\), whose exponential is one; hence the formula
also holds across that wrap boundary.  It follows that
\[
 U_{12}(n)=
 \exp\left(-\frac{2\pi(N-1)}{N^2}H_{\mathrm{link}}\right)
 \exp\left(\frac{2\pi}{N}H_{\mathrm{link}}\right)
 =\exp(2\pi H_{\mathrm{link}}/N^2).
\]
The functions \(U_1,U_2\) are independent of \(n_3,n_4\), and
\(U_3=U_4=1\), proving every remaining assertion in
\eqref{cusp-link-plaquette-values}.  Reversing the orientation gives
\(U_{21}(n)=U_{12}(n)^{-1}\); it is not an additional trivial plaquette.

The cover Gram matrix has lower-right principal block
\(M_{\mathrm{cov}}^2I_2\).  Jacobi's identity for complementary minors
therefore gives
\[
 \kappa^{12}_{R,M_{\mathrm{cov}}}
 =\frac{\det(M_{\mathrm{cov}}^2I_2)}
        {\det G_{R,M_{\mathrm{cov}}}}
 =\frac{M_{\mathrm{cov}}^4}{M_{\mathrm{cov}}^4D_R^2}
 =D_R^{-2},
\]
which proves \eqref{cusp-link-two-form-minor}.  By
\eqref{cusp-link-flux-generator} and the principal-log choice,
\(N^2\operatorname{Log}U_{12}=2\pi H_{\mathrm{link}}\).  Hence every one of the \(N^4\)
summands in \eqref{cusp-link-log-action} equals
\(4\pi^2D_R^{-2}\).  Substitution of
\(\mathcal D_{R,M_{\mathrm{cov}}}=M_{\mathrm{cov}}^2D_R\) proves
\eqref{cusp-link-log-action-value}.

Finally, \(H_{\mathrm{link}}\) has eigenvalues \(\pm i\), so
\[
 \operatorname{Re}\operatorname{tr}\exp(2\pi H_{\mathrm{link}}/N^2)
 =2\cos(2\pi/N^2).
\]
There are \(N^4\) equal plaquettes.  Substitution proves
\eqref{cusp-link-wilson-sector-value}.  Applying
\[
 1-\cos x=\frac{x^2}{2}-\frac{x^4}{24}
           +\frac{x^6}{720}+O(x^8)
\]
with \(x=2\pi/N^2\) proves
\eqref{cusp-link-wilson-sector-expansion} and
\eqref{cusp-link-wilson-sector-limit}.
\end{proof}

\subsection{Finite Gaussian quantization with independent cover and mesh}

Let
\begin{equation}
 \vartheta^j_{R,M_{\mathrm{cov}}}
 =P_{R,M_{\mathrm{cov}}}^{-\mathsf T}\mathbf e_j
 \in(\mathbb R^4)^*,
 \qquad 1\leq j\leq4.
 \label{cusp-link-physical-coframe}
\end{equation}
For \(0\leq p\leq4\), define the real finite cochain spaces
\begin{equation}
 C^p_{N_{\mathrm{mesh}}}(R,M_{\mathrm{cov}})
 =\operatorname{Map}\left(
 V_{N_{\mathrm{mesh}}},
 \bigwedge^p(\mathbb R^4)^*\otimes\mathfrak{su}(2)
 \right).
 \label{cusp-link-cochain-spaces}
\end{equation}
The exterior powers carry the metric induced by the physical coframe
\eqref{cusp-link-physical-coframe}; explicitly, for increasing multi-indices
\(I=(i_1,\ldots,i_p)\) and \(J=(j_1,\ldots,j_p)\),
\begin{equation}
 \left\langle
  \vartheta^{i_1}\wedge\cdots\wedge\vartheta^{i_p},
  \vartheta^{j_1}\wedge\cdots\wedge\vartheta^{j_p}
 \right\rangle
 =
 \det\left((G_{R,M_{\mathrm{cov}}}^{-1})_{i_a j_b}\right)_{a,b=1}^p.
 \label{cusp-link-exterior-metric}
\end{equation}
Together with \eqref{cusp-link-su2-inner-product}, this defines the
pointwise tensor-product metric.  The cochain inner product is
\begin{equation}
 \langle a,b\rangle_p
 =\frac{\mathcal D_{R,M_{\mathrm{cov}}}}{N_{\mathrm{mesh}}^4}
   \sum_{n\in V_{N_{\mathrm{mesh}}}}
   \langle a(n),b(n)\rangle.
 \label{cusp-link-cochain-inner-product}
\end{equation}
Let \(\mathsf T_j a(n)=a(n+\mathbf e_j)\), set
\begin{equation}
 D^+_{N_{\mathrm{mesh}},j}
 =N_{\mathrm{mesh}}(\mathsf T_j-1),
 \qquad
 D^-_{N_{\mathrm{mesh}},j}a(n)
 =N_{\mathrm{mesh}}(a(n)-a(n-\mathbf e_j)),
 \label{cusp-link-differences}
\end{equation}
and define
\begin{equation}
 d_{N_{\mathrm{mesh}}}a
 =\sum_{j=1}^4
  \varepsilon(\vartheta^j_{R,M_{\mathrm{cov}}})
  D^+_{N_{\mathrm{mesh}},j}a.
 \label{cusp-link-finite-exterior-derivative}
\end{equation}
The commuting differences and anticommuting exterior multiplications give
\(d_{N_{\mathrm{mesh}}}^2=0\).  Its adjoint is
\begin{equation}
 d_{N_{\mathrm{mesh}}}^*
 =-\sum_{j=1}^4
  \iota_{(\vartheta^j_{R,M_{\mathrm{cov}}})^\sharp}
  D^-_{N_{\mathrm{mesh}},j}.
 \label{cusp-link-finite-adjoint}
\end{equation}
Define the real reduced Coulomb space
\begin{align}
 \mathscr C_{R,M_{\mathrm{cov}},N_{\mathrm{mesh}}}
 &=
 \ker(d_{N_{\mathrm{mesh}}}^*:C^1\to C^0)
 \cap
 \left(\ker\Delta^{(1)}_{N_{\mathrm{mesh}}}\right)^\perp,
 \label{cusp-link-reduced-coulomb-space}\\
 \Delta^{(1)}_{N_{\mathrm{mesh}}}
 &=d_{N_{\mathrm{mesh}}}d_{N_{\mathrm{mesh}}}^*
 +d_{N_{\mathrm{mesh}}}^*d_{N_{\mathrm{mesh}}}.
 \label{cusp-link-finite-hodge-laplacian}
\end{align}
On this finite-dimensional space put
\begin{equation}
 \mathcal L_{R,M_{\mathrm{cov}},N_{\mathrm{mesh}}}
 =g_{\mathrm{YM}}^{-2}d_{N_{\mathrm{mesh}}}^*
                         d_{N_{\mathrm{mesh}}}.
 \label{cusp-link-finite-hessian}
\end{equation}

\begin{proposition}[Exact finite Hessian spectrum]
\label{cusp-link-finite-spectrum}
For
\begin{equation}
 \mathfrak d_{N_{\mathrm{mesh}}}(k)_j
 =N_{\mathrm{mesh}}
  \left(e^{2\pi i k_j/N_{\mathrm{mesh}}}-1\right),
 \qquad
 k\in(\mathbb Z/N_{\mathrm{mesh}}\mathbb Z)^4,
 \label{cusp-link-fourier-symbol}
\end{equation}
the positive spectrum of \(\mathcal L\) is obtained from
\begin{equation}
 \lambda_{R,M_{\mathrm{cov}},N_{\mathrm{mesh}}}(k)
 =\frac1{g_{\mathrm{YM}}^2}
  \left\lVert
   P_R^{-\mathsf T}S_{M_{\mathrm{cov}}}^{-1}
   \mathfrak d_{N_{\mathrm{mesh}}}(k)
  \right\rVert^2,
 \qquad k\neq0.
 \label{cusp-link-finite-mode-eigenvalue}
\end{equation}
Each complex Fourier fibre has transverse multiplicity
\(3\dim\mathfrak{su}(2)=9\).  In particular,
\begin{equation}
 \lambda_{\min}(\mathcal L)
 =\min_{k\neq0}\lambda_{R,M_{\mathrm{cov}},N_{\mathrm{mesh}}}(k)>0.
 \label{cusp-link-finite-minimum}
\end{equation}
For \(\mathbf e_1=(1,0,0,0)^{\mathsf T}\), the exact mode value is
\begin{equation}
 \lambda_{R,M_{\mathrm{cov}},N_{\mathrm{mesh}}}(\mathbf e_1)
 =\frac{4N_{\mathrm{mesh}}^2
              \sin^2(\pi/N_{\mathrm{mesh}})}{g_{\mathrm{YM}}^2}
  \frac{L_R^2+m_R^2}{D_R^2};
 \label{cusp-link-cover-soft-mode}
\end{equation}
it is independent of \(M_{\mathrm{cov}}\).
\end{proposition}

\begin{proof}
Translation is unitary for \eqref{cusp-link-cochain-inner-product}, so
\((D_j^+)^*=-D_j^-\).  Exterior multiplication by \(\vartheta^j\) has
adjoint contraction by \((\vartheta^j)^\sharp\), proving
\eqref{cusp-link-finite-adjoint}.  On the Fourier character
\(\chi_k(n)=e^{2\pi i k\cdot n/N_{\mathrm{mesh}}}\), the derivative symbol
is exterior multiplication by
\[
 \zeta(k)=P_{R,M_{\mathrm{cov}}}^{-\mathsf T}
           \mathfrak d_{N_{\mathrm{mesh}}}(k)
 =P_R^{-\mathsf T}S_{M_{\mathrm{cov}}}^{-1}
           \mathfrak d_{N_{\mathrm{mesh}}}(k).
\]
The Hermitian exterior-algebra identity
\[
 \iota_{\overline\zeta^\sharp}\varepsilon(\zeta)
 +\varepsilon(\zeta)\iota_{\overline\zeta^\sharp}
 =\lVert\zeta\rVert^2\operatorname{id}
\]
shows that on the Coulomb polarization
\(\ker\iota_{\overline\zeta^\sharp}\subset(\mathbb C^4)^*\), the operator
\(d^*d\) is multiplication by \(\lVert\zeta\rVert^2\).  For \(k\neq0\),
\(\mathfrak d_{N_{\mathrm{mesh}}}(k)\neq0\), so invertibility of
\(P_{R,M_{\mathrm{cov}}}\) gives
\(\zeta(k)\neq0\).  The contraction has rank one, and its kernel has
dimension three.  The zero Fourier fibre consists exactly of constant
harmonic cochains and was removed in
\eqref{cusp-link-reduced-coulomb-space}.  Fourier completeness now proves
\eqref{cusp-link-finite-mode-eigenvalue}--\eqref{cusp-link-finite-minimum}.

For \(k=\mathbf e_1\), the only nonzero component of
\(\mathfrak d_{N_{\mathrm{mesh}}}(k)\) has modulus
\[
 2N_{\mathrm{mesh}}\sin(\pi/N_{\mathrm{mesh}}),
\]
and
\(S_{M_{\mathrm{cov}}}^{-1}\mathbf e_1=\mathbf e_1\).  Solving
\(P_R^{\mathsf T}\xi=\mathbf e_1\), or directly inverting the lower
\(2\times2\) block in \eqref{cusp-link-period-matrix}, gives
\[
 \lVert P_R^{-\mathsf T}\mathbf e_1\rVert^2
 =\frac{L_R^2+m_R^2}{D_R^2}.
\]
Substitution proves \eqref{cusp-link-cover-soft-mode}.
\end{proof}

Let \(da\) denote Lebesgue measure on
\(\mathscr C_{R,M_{\mathrm{cov}},N_{\mathrm{mesh}}}\) induced by
\eqref{cusp-link-cochain-inner-product} and define
\begin{equation}
 \mathcal H^{\mathrm{Gauss}}_{R,M_{\mathrm{cov}},N_{\mathrm{mesh}}}
 =L^2(\mathscr C_{R,M_{\mathrm{cov}},N_{\mathrm{mesh}}},da).
 \label{cusp-link-gaussian-hilbert-space}
\end{equation}

\begin{theorem}[Finite Gaussian oscillator and a resolved cusp sequence]
\label{cusp-link-gaussian-theorem}
The Friedrichs realization on \eqref{cusp-link-gaussian-hilbert-space} of
\begin{equation}
 \widehat H^{\mathrm{Gauss}}_{R,M_{\mathrm{cov}},N_{\mathrm{mesh}}}
 =-\frac12\Delta_{\mathscr C}
  +\frac12\langle a,
     \mathcal L_{R,M_{\mathrm{cov}},N_{\mathrm{mesh}}}a\rangle
 \label{cusp-link-gaussian-hamiltonian}
\end{equation}
has a unique Gaussian ground state and spectral gap
\begin{equation}
 \operatorname{gap}(\widehat H^{\mathrm{Gauss}})
 =\sqrt{\lambda_{\min}(
   \mathcal L_{R,M_{\mathrm{cov}},N_{\mathrm{mesh}}})}.
 \label{cusp-link-gaussian-gap}
\end{equation}
Consequently,
\begin{equation}
 \operatorname{gap}(\widehat H^{\mathrm{Gauss}})
 \leq
 \frac{2N_{\mathrm{mesh}}\sin(\pi/N_{\mathrm{mesh}})}{g_{\mathrm{YM}}}
 \frac{\sqrt{L_R^2+m_R^2}}{D_R}.
 \label{cusp-link-gaussian-gap-bound}
\end{equation}

There is an explicit diagonal sequence for which the cover geometry tends locally to
\(\mathbb R^4\), the cubical mesh resolves that geometry, and the Gaussian
gap tends to zero.  Explicitly, for integers \(j\to\infty\), take
\begin{equation}
 R_j=j^2,\qquad M_{\mathrm{cov},j}=j,\qquad
 N_{\mathrm{mesh},j}=j^3.
 \label{cusp-link-diagonal-sequence}
\end{equation}
Then
\begin{equation}
 \begin{aligned}
 \operatorname{sys}(\widetilde F_{R_j,M_{\mathrm{cov},j}})
 &\longrightarrow\infty,\\
 \max_r
 \frac{\lVert P_{R_j,M_{\mathrm{cov},j}}\mathbf e_r\rVert}
      {N_{\mathrm{mesh},j}}
 &\longrightarrow0,\\
 \operatorname{gap}(\widehat H^{\mathrm{Gauss}})
 &\longrightarrow0,\\
 S^{\log}_{R_j,M_{\mathrm{cov},j},N_{\mathrm{mesh},j}}(U)
 &\longrightarrow0.
 \end{aligned}
 \label{cusp-link-diagonal-conclusions}
\end{equation}
Thus the pointed cover tori converge locally isometrically to \(\mathbb R^4\)
while their finite linearized Gaussian systems are simultaneously refined.
\end{theorem}

\begin{proof}
The real finite-dimensional operator \(\mathcal L\) is positive definite on
the reduced Coulomb space.  Choose an orthonormal eigenbasis with eigenvalues
\(0<\lambda_1\leq\cdots\leq\lambda_r\).  In the resulting coordinates,
\eqref{cusp-link-gaussian-hamiltonian} is the tensor sum
\[
 \sum_{s=1}^r
 \left(-\frac12\frac{\partial^2}{\partial x_s^2}
       +\frac12\lambda_sx_s^2\right).
\]
The normalized Hermite basis has eigenvalues
\(\sum_s(n_s+\tfrac12)\sqrt{\lambda_s}\), with
\(n_s\in\mathbb Z_{\geq0}\).  The all-zero multi-index is the unique ground
state, and the smallest possible positive excitation is
\(\sqrt{\lambda_1}\).  This proves \eqref{cusp-link-gaussian-gap}; applying
the minimum to the explicit \(\mathbf e_1\)-mode in
\eqref{cusp-link-cover-soft-mode} proves
\eqref{cusp-link-gaussian-gap-bound}.

For completeness, reorder physical coordinates as in
\eqref{cusp-link-block-form}.  Every vector of the cover lattice has the
form
\[
 \binom{\mathsf A_R^{\mathrm{blk}}n+M_{\mathrm{cov}}\ell}
       {\mathsf B_R^{\mathrm{blk}}n},
 \qquad n,\ell\in\mathbb Z^2.
\]
Let \(C_{\mathrm{blk}}=\sup_R
\lVert\mathsf C_R^{\mathrm{blk}}\rVert_{\mathrm{op}}\).  If \(n\neq0\), then
\(\lVert\mathsf B_R^{\mathrm{blk}}n\rVert
\geq(R-C_{\mathrm{blk}})\lVert n\rVert\); if \(n=0\) and the lattice
vector is nonzero, its norm is at least \(M_{\mathrm{cov}}\).  Hence
\begin{equation}
 \operatorname{sys}(\widetilde F_{R,M_{\mathrm{cov}}})
 \geq\min\{R-C_{\mathrm{blk}},M_{\mathrm{cov}}\}.
 \label{cusp-link-cover-systole-bound}
\end{equation}
This tends to infinity for the first two choices in
\eqref{cusp-link-diagonal-sequence}.  For any fixed \(r>0\), once the
systole exceeds \(4r\), the quotient map from \(\mathbb R^4\) is injective
and distance-preserving on the radius-\(r\) ball: for a nonzero period
\(\lambda\) and \(x,y\) in that ball,
\[
 \lVert x-y+\lambda\rVert
 \geq\lVert\lambda\rVert-\lVert x-y\rVert
 >4r-2r\geq\lVert x-y\rVert.
\]
This proves the pointed local-isometric \(\mathbb R^4\) limit.

The first two columns of \(P_R\) have norm \(O(R)\), while the last two
columns of \(P_{R,M_{\mathrm{cov}}}\) have norm exactly
\(M_{\mathrm{cov}}\).  Therefore the largest physical edge in the
\(N_{\mathrm{mesh}}\)-subdivision has length
\(O(\max\{R,M_{\mathrm{cov}}\}/N_{\mathrm{mesh}})\), which is \(O(j^{-1})\)
for \eqref{cusp-link-diagonal-sequence}.  Finally,
\eqref{cusp-link-cusp-asymptotics} gives
\(\sqrt{L_R^2+m_R^2}/D_R=R^{-1}(1+O(R^{-1}))\), while
\(2N\sin(\pi/N)\leq2\pi\).  The right-hand side of
\eqref{cusp-link-gaussian-gap-bound} is therefore \(O(j^{-2})\).
Moreover, \eqref{cusp-link-log-action-value} is \(O(j^{-2})\), since
\(M_{\mathrm{cov},j}^2/D_{R_j}=j^2/(j^4+O(j^2))\).
This proves every assertion in \eqref{cusp-link-diagonal-conclusions}.
\end{proof}

\subsection{The exponential map and the exact physical boundary}

Let \(C^1=C^1_{N_{\mathrm{mesh}}}(R,M_{\mathrm{cov}})\), regarded only as a
finite-dimensional Lie-algebra cochain space, and define componentwise
\begin{equation}
 \operatorname{Exp}_{N_{\mathrm{mesh}}}:C^1
 \longrightarrow\mathrm{SU}(2)^{E_{N_{\mathrm{mesh}}}},
 \qquad
 a\longmapsto\left(\exp(a_e/N_{\mathrm{mesh}})\right)_e.
 \label{cusp-link-exponential-map}
\end{equation}
Here, if
\(a(n)=\sum_{j=1}^4\vartheta^j_{R,M_{\mathrm{cov}}}\otimes a_j(n)\),
then \(a_e=a_j(n)\) for the oriented edge \(e=(n,j)\).

\begin{proposition}[Local chart, not a global embedding]
\label{cusp-link-exponential-local-only}
The map \eqref{cusp-link-exponential-map} is smooth and a local
diffeomorphism at zero, but it is not globally injective.  Under left
trivialization of every tangent space at the identity,
\begin{equation}
 D\operatorname{Exp}_{N_{\mathrm{mesh}}}(0)\,a
 =N_{\mathrm{mesh}}^{-1}a.
 \label{cusp-link-exponential-derivative}
\end{equation}
If \(S_W\) is any smooth compact-link action for which the identity
configuration is critical, then
\begin{equation}
 D^2(S_W\circ\operatorname{Exp}_{N_{\mathrm{mesh}}})_0[a,b]
 =N_{\mathrm{mesh}}^{-2}
  D^2(S_W)_{\mathbf1}[a,b]
 \label{cusp-link-exponential-hessian-scaling}
\end{equation}
under the same tangent identification.  Thus, for a prescribed cochain
quadratic form \(Q\),
\[
 D^2(S_W\circ\operatorname{Exp}_{N_{\mathrm{mesh}}})_0=Q
\]
holds exactly when
\[
 D^2(S_W)_{\mathbf1}=N_{\mathrm{mesh}}^2Q
\]
under the coefficient identification.  Both this scaling and the tensor
structure must be verified for the chosen compact-link action; neither
follows from the exponential map alone.
\end{proposition}

\begin{proof}
For the Lie-group exponential, \(D\exp_0\) is the identity.  Applying the
chain rule to \(a_e\mapsto a_e/N_{\mathrm{mesh}}\) gives
\eqref{cusp-link-exponential-derivative}, which is invertible, so the inverse
function theorem gives a local diffeomorphism.  It is not globally
injective: a cochain which is zero on every edge except one, where it equals
\(2\pi N_{\mathrm{mesh}}H_{\mathrm{link}}\), has the same image as the zero cochain by
\eqref{cusp-link-flux-generator}.  Finally, the second-derivative chain rule
contains a term involving \(DS_W(\mathbf1)\), which vanishes by criticality;
the two remaining first derivatives of \(\operatorname{Exp}\) each supply
the factor \(N_{\mathrm{mesh}}^{-1}\).  This proves
\eqref{cusp-link-exponential-hessian-scaling}.
\end{proof}

The Gaussian Hilbert space in \eqref{cusp-link-gaussian-hilbert-space} is
not the compact-link Yang--Mills Hilbert space.  At this finite lattice, the
latter kinematic and gauge-invariant spaces would be
\begin{equation}
 \mathcal H^{\mathrm{link}}_{N_{\mathrm{mesh}}}
 =L^2(\mathrm{SU}(2)^{E_{N_{\mathrm{mesh}}}},dU),
 \qquad
 \mathcal H^{\mathrm{inv}}_{N_{\mathrm{mesh}}}
 =\left(\mathcal H^{\mathrm{link}}_{N_{\mathrm{mesh}}}\right)^{
    \mathrm{SU}(2)^{V_{N_{\mathrm{mesh}}}}},
 \label{cusp-link-physical-hilbert-space}
\end{equation}
where \(dU\) is product normalized Haar measure and the superscript denotes
the nonlinear vertex-gauge action.  Here ``gauge-invariant'' refers only to
the link-configuration space of this four-dimensional cell complex.  Without
a chosen time direction or transfer matrix, it is not asserted to be a
canonical Hamiltonian physical Hilbert space.  L\"uscher's construction of
a selfadjoint, strictly positive transfer matrix
\cite{Luscher_1977} and Menotti--Pelissetto's reflection-positivity theorem
for the Wilson action \cite{MenottiPelissetto_1987} concern the standard
Wilson lattice setting after the requisite reflection/time structure has
been fixed; they do not supply those omitted choices for the nonorthogonal
sector functional \eqref{cusp-link-wilson-sector-definition}.  The
Osterwalder--Schrader reconstruction boundary is likewise the one stated in
the primary Euclidean axioms papers
\cite{OsterwalderSchrader_1973,OsterwalderSchrader_1975}, not a consequence of
a finite Gaussian oscillator calculation.  The differences are
structural:
\begin{enumerate}
\item \(\mathcal H^{\mathrm{Gauss}}\) is an \(L^2\)-space on a noncompact
linear Coulomb slice with Lebesgue measure; \(\mathcal H^{\mathrm{inv}}\)
is a gauge-invariant subspace of an \(L^2\)-space on compact link variables
with Haar measure.
\item The local exponential map is neither a global coordinate chart nor a
measure-preserving bijection.  No unitary map between the two displayed
Hilbert spaces has been constructed.
\item A link configuration \(U\) determines the Dirac distribution
\(\delta_U\), but
\begin{equation}
 \delta_U\in\mathcal D'(\mathrm{SU}(2)^{E_{N_{\mathrm{mesh}}}}),
 \qquad
 \delta_U\notin L^2(\mathrm{SU}(2)^{E_{N_{\mathrm{mesh}}}},dU).
 \label{cusp-link-dirac-not-l2}
\end{equation}
Indeed, product Haar measure is nonatomic, so a point mass is singular with
respect to it and has no square-integrable density.
\item No nonlinear compact-link Hamiltonian, Gauss-law spectral problem,
time direction, reflection-positive transfer matrix, renormalized-scale
limit, or Osterwalder--Schrader reconstruction is supplied here.
\end{enumerate}
Accordingly, \eqref{cusp-link-gaussian-gap} is an exact theorem about the
finite quadratic oscillator defined by the Coulomb-restricted operator
\(g_{\mathrm{YM}}^{-2}d^*d\).  It is not an eigenvalue statement about a
nonperturbative physical Yang--Mills
Hamiltonian, and the small classical action in
\eqref{cusp-link-log-action-value} is not by itself a physical-state Rayleigh
quotient.

% Standalone-preprint adaptation of cusp_su2_magnetic.tex.
% Canonical source SHA-256: 8676F460516102827CCC32DAB613034F2153B8CE83EA87679AC6D6433121BE2A.
% The canonical radial-height symbol T was renamed R consistently; transpose
% notation was retained exactly.

\section{Integral magnetic sectors and a reducible
\texorpdfstring{\(SU(2)\)}{SU(2)} critical point on the radial cusp}
\label{cusp-magnetic-section}

Retain the positive radial flat tori
\(F_R=\mathbb R^4/P_R\mathbb Z^4\), their Euclidean quotient metrics, and
the real quantities \(L_R,q_R,m_R,D_R,H,B,m\) from
\cref{cusp-spec-radial-notation,cusp-spec-exact-dual-lattice,cusp-spec-D-expansion}.
In particular,
\[
 L_R=\operatorname{Im}\tau_R,\qquad
 q_R=-\operatorname{Im}\beta_R,\qquad
 D_R=L_Rq_R+6m_R^2=\operatorname{Vol}(F_R)>0.
\]
The real coordinate order is the already fixed order
\[
 x=(x_1,y_1,x_2,y_2)^{\mathsf t}.
\]

\subsection{The integral coframe}

Let \(e_1,\ldots,e_4\) be the standard basis of \(\mathbb R^4\), put
\(u=P_R^{-1}x\), and define
\begin{equation}
 \theta_R^j:=\mathrm du_j
 =(P_R^{-\mathsf T}e_j)^{\mathsf t}
   (\mathrm dx_1,\mathrm dy_1,\mathrm dx_2,\mathrm dy_2)^{\mathsf t}.
 \label{cusp-magnetic-integral-coframe}
\end{equation}

\begin{proposition}[Integral coframe and the distinguished unit-flux plane]
\label{cusp-magnetic-coframe-proposition}
The forms \(\theta_R^j\) are the integral coframe dual to the ordered
marked lattice basis:
\[
 \int_{[0,1]\ni s\mapsto[x+sP_Re_k]}\theta_R^j=\delta_{jk}.
\]
The first two members are exactly
\begin{equation}
 \theta_R^1=\frac{m_R}{D_R}\,\mathrm dy_1
             -\frac{L_R}{D_R}\,\mathrm dy_2,
 \qquad
 \theta_R^2=\frac{q_R}{D_R}\,\mathrm dy_1
             +\frac{6m_R}{D_R}\,\mathrm dy_2,
 \label{cusp-magnetic-first-two-covectors}
\end{equation}
and hence
\begin{equation}
 \theta_R^1\wedge\theta_R^2
 =\frac1{D_R}\,\mathrm dy_1\wedge\mathrm dy_2,
 \qquad
 \lvert\theta_R^1\wedge\theta_R^2\rvert^2
 =\frac1{D_R^2}.
 \label{cusp-magnetic-wedge-norm}
\end{equation}
\end{proposition}

\begin{proof}
For \(n\in\mathbb Z^4\),
\[
 u(x+P_Rn)=u(x)+n,
\]
which proves the integral-period assertion.  The first two columns of the
already-proved matrix \(P_R^{-\mathsf T}\) in
\cref{cusp-spec-full-inverse-transpose} give
\cref{cusp-magnetic-first-two-covectors}.  Their wedge coefficient is
\[
 \frac{6m_R^2+L_Rq_R}{D_R^2}=\frac1{D_R}.
\]
The coordinate forms \(\mathrm dy_1,\mathrm dy_2\) are orthonormal for the fixed quotient
metric, proving \cref{cusp-magnetic-wedge-norm}.
\end{proof}

\subsection{The unit-flux line and the topologically trivial rank-two sum}

The following construction is made in the marked coordinates \(u\) of
\eqref{cusp-magnetic-integral-coframe}; it therefore applies to the actual
skew lattice \(P_R\mathbb Z^4\), not to a rectangular replacement.

\begin{proposition}[Explicit factor of automorphy and connection]
\label{cusp-magnetic-line-bundle-proposition}
For \(n\in\mathbb Z^4\), set
\begin{equation}
 j_n(u)=\exp(-2\pi i n_1u_2).
 \label{cusp-magnetic-unit-factor}
\end{equation}
The quotient
\[
 \mathscr L_R=(\mathbb R^4\times\mathbb C)/\mathbb Z^4,
 \qquad
 n\cdot(u,z)=(u+n,j_n(u)z),
\]
is a Hermitian line bundle over \(F_R\).  The connection on its pullback
to \(\mathbb R^4\),
\begin{equation}
 \widetilde\nabla^{\mathscr L_R}
 =\mathrm d+2\pi i u_1\,\mathrm du_2,
 \label{cusp-magnetic-line-connection}
\end{equation}
descends, and its curvature is
\begin{equation}
 F_{\mathscr L_R}=2\pi i\,\mathrm du_1\wedge\mathrm du_2
        =2\pi i\,\theta_R^1\wedge\theta_R^2.
 \label{cusp-magnetic-line-curvature}
\end{equation}
With the unitary Chern--Weil convention
\(c_1=[\,iF/(2\pi)\,]\), one has
\(c_1(\mathscr L_R)=-[\theta_R^1\wedge\theta_R^2]\).
\end{proposition}

\begin{proof}
For \(m,n\in\mathbb Z^4\),
\[
 j_m(u+n)j_n(u)
 =j_{m+n}(u)\exp(-2\pi i m_1n_2)
 =j_{m+n}(u).
\]
Thus \eqref{cusp-magnetic-unit-factor} is a factor of automorphy.  Its
values are unitary, so the standard fibre metric descends.  If
\(\mathcal A=2\pi i u_1\,\mathrm du_2\), then
\[
 \mathcal A(u+n)=\mathcal A(u)+2\pi i n_1\,\mathrm du_2,
 \qquad
 j_n^{-1}\mathrm dj_n=-2\pi i n_1\,\mathrm du_2.
\]
Hence \(\mathcal A(u+n)+j_n^{-1}\mathrm dj_n=\mathcal A(u)\), which is precisely
the descent identity for \(\mathrm d+\mathcal A\) under the stated quotient
convention.  Taking \(\mathrm d\mathcal A\) gives
\eqref{cusp-magnetic-line-curvature}.  The Chern--Weil statement follows
without a sign suppression:
\[
 \frac{i}{2\pi}F_{\mathscr L_R}=-\theta_R^1\wedge\theta_R^2.
\]
\end{proof}

Let \(\mathscr L_R^{-1}=\mathscr L_R^*\) carry the dual metric and dual
connection, and set
\begin{equation}
 \mathscr E_R=\mathscr L_R\oplus\mathscr L_R^{-1}.
 \label{cusp-magnetic-rank-two-bundle}
\end{equation}

\begin{proposition}[Determinant, \(c_2\), and explicit
\(SU(2)\)-triviality]
\label{cusp-magnetic-su2-triviality}
The factors of automorphy of \(\mathscr E_R\) are
\[
 G_n(u)=\begin{pmatrix}j_n(u)&0\\0&j_n(u)^{-1}\end{pmatrix}\in SU(2).
\]
Evaluation gives a canonical parallel determinant trivialization
\[
 \det\mathscr E_R=\mathscr L_R\otimes\mathscr L_R^{-1}
 \xrightarrow{\ \mathrm{ev}\ }F_R\times\mathbb C.
\]
Moreover
\[
 c_1(\mathscr E_R)=0,
 \qquad
 c_2(\mathscr E_R)=0\quad\text{in }H^4(F_R;\mathbb Z),
\]
and the underlying principal \(SU(2)\)-bundle is smoothly trivial.  The
last assertion follows from the explicit nullhomotopy below and does not
rely only on the vanishing of a de Rham Chern form.
\end{proposition}

\begin{proof}
The determinant assertion follows pointwise from
\(j_nj_n^{-1}=1\).  The induced determinant connection has connection
form
\[
 \operatorname{tr}
 \begin{pmatrix}
  2\pi i u_1\,\mathrm du_2&0\\0&-2\pi i u_1\,\mathrm du_2
 \end{pmatrix}=0,
\]
so the evaluation trivialization is parallel.

Let \(c=c_1(\mathscr L_R)\).  The integral Whitney identity gives
\[
 c(\mathscr E_R)=(1+c)(1-c)=1-c^2.
\]
The factor \(j_n\) depends only on \((u_1,u_2)\), so \(\mathscr L_R\) is the
pullback of a line bundle \(L_0\) on
\(\mathbb R^2/\mathbb Z^2\) under the marked projection
\([u_1,u_2,u_3,u_4]\mapsto[u_1,u_2]\).  Hence
\[
 c^2=\pi_{12}^*\bigl(c_1(L_0)^2\bigr)=0
\]
integrally because \(H^4(\mathbb T^2;\mathbb Z)=0\).  This proves the two
Chern-class identities on the four-torus.

It remains to trivialize the \(SU(2)\)-bundle itself.  On the fundamental
square \(0\le u_1,u_2\le1\) for the first two marked coordinates, the only
nontrivial seam is
\[
 (0,t,v)\sim(1,t,\gamma(t)v),
 \qquad
 \gamma(t)=
 \begin{pmatrix}e^{-2\pi it}&0\\0&e^{2\pi it}\end{pmatrix}.
\]
For \(0\le r\le1\), define
\[
 a(r,t)=(1-r)e^{-2\pi it}+r,
 \qquad
 s(r)=\sin(\pi r),
 \qquad
 N(r,t)=\bigl(|a(r,t)|^2+s(r)^2\bigr)^{1/2},
\]
and
\begin{equation}
 U(r,t)=\frac1{N(r,t)}
 \begin{pmatrix}
  a(r,t)&-s(r)\\
  s(r)&\overline{a(r,t)}
 \end{pmatrix}.
 \label{cusp-magnetic-explicit-nullhomotopy}
\end{equation}
If \(0<r<1\), then \(s(r)>0\); at \(r=0,1\), one has
\(|a(r,t)|=1\).  Thus \(N\) never vanishes.  The two columns in
\eqref{cusp-magnetic-explicit-nullhomotopy} are orthonormal and its
determinant is one, so \(U(r,t)\in SU(2)\).  It is smooth and periodic in
\(t\), and
\[
 U(0,t)=\gamma(t),
 \qquad
 U(1,t)=I_2.
\]
This is an explicit nullhomotopy of the clutching loop.  To see directly
that it produces a smooth global frame, choose a smooth
\(\chi:[0,1]\to[0,1]\) which is zero near \(0\) and one near \(1\), and put
\[
 K(\sigma,t)=U(1-\chi(\sigma),t).
\]
Then \(K=I_2\) near \(\sigma=0\), \(K=\gamma\) near \(\sigma=1\), and
\(K(\sigma,0)=K(\sigma,1)\).  The required trivialization, written in the
direction from the clutched bundle to the product bundle, is
\begin{equation}
 [\sigma,t,v]\longmapsto
 \bigl([\sigma,t],K(\sigma,t)^{-1}v\bigr).
 \label{cusp-magnetic-explicit-global-frame}
\end{equation}
Indeed, the quotient convention gives
\([0,t,v]=[1,t,\gamma(t)v]\), and the two representatives have the same
image because
\[
 K(1,t)^{-1}\gamma(t)v
 =\gamma(t)^{-1}\gamma(t)v
 =v=K(0,t)^{-1}v.
\]
At the other seam,
\(K(\sigma,0)^{-1}=K(\sigma,1)^{-1}\).  The constant collars used in the
definition of \(\chi\) make the local matrix in
\cref{cusp-magnetic-explicit-global-frame} constant after applying the
seam transition, so the map and its inverse are smooth across the seam.
It is fibrewise unitary and determinant one.  Extending it independently
of \(u_3,u_4\) and
pulling it through the marking gives the required global
\(SU(2)\)-frame of \(\mathscr E_R\).
\end{proof}

\subsection{A non-flat reducible Yang--Mills critical connection}

The symbol \(H=\operatorname{Im}h(0)\) is already fixed by
\cref{cusp-spec-limiting-imaginary-parts}.  For the requested Lie-algebra
element, fix instead the unambiguous notation
\begin{equation}
 \langle X,Y\rangle_{\mathfrak{su}(2)}
   =-\frac12\operatorname{Tr}(XY),
 \qquad
 H_{\mathrm{mag}}=\begin{pmatrix}i&0\\0&-i\end{pmatrix}.
 \label{cusp-magnetic-su2-normalization}
\end{equation}
\begin{theorem}[Reducible critical point and exact action]
\label{cusp-magnetic-ym-critical-point}
The direct-sum connection on \(\mathscr E_R\) has, in the lifted splitting,
\begin{equation}
 \nabla^R=\mathrm d+2\pi u_1\,\mathrm du_2\,H_{\mathrm{mag}},
 \qquad
 F_{\nabla^R}
 =2\pi H_{\mathrm{mag}}\,\theta_R^1\wedge\theta_R^2
 =\frac{2\pi}{D_R}H_{\mathrm{mag}}\,\mathrm dy_1\wedge\mathrm dy_2.
 \label{cusp-magnetic-su2-curvature}
\end{equation}
After any of the global \(SU(2)\)-trivializations from
 \cref{cusp-magnetic-su2-triviality}, this is a connection on the
trivial principal \(SU(2)\)-bundle.  It is non-flat and satisfies
\[
 \mathrm d_{\nabla^R}^*F_{\nabla^R}=0.
\]
For the normalization
\begin{equation}
 \operatorname{YM}_R(\nabla)
 =\frac1{2g_{\mathrm{YM}}^2}
  \int_{F_R}\lvert F_\nabla\rvert^2\,
  \mathrm d\operatorname{vol}_{F_R},
 \label{cusp-magnetic-action-normalization}
\end{equation}
its action is exactly
\begin{equation}
 \operatorname{YM}_R(\nabla^R)
 =\frac{2\pi^2}{g_{\mathrm{YM}}^2D_R}.
 \label{cusp-magnetic-exact-action}
\end{equation}
Its first inverse-power correction, in the canonical scalar notation
\(H=\operatorname{Im}h(0)\), \(B=\operatorname{Im}b(0)\), is
\begin{equation}
 \operatorname{YM}_R(\nabla^R)
 =\frac{2\pi^2}{g_{\mathrm{YM}}^2R^2}
  \left(1+\frac{B-2H}{R}
  +O(R^{-2})\right).
 \label{cusp-magnetic-action-first-correction}
\end{equation}
\end{theorem}

\begin{proof}
The connection matrix is
\[
 \begin{pmatrix}
  2\pi i u_1\,\mathrm du_2&0\\
  0&-2\pi i u_1\,\mathrm du_2
 \end{pmatrix}
 =2\pi u_1\,\mathrm du_2\,H_{\mathrm{mag}},
\]
so \cref{cusp-magnetic-su2-curvature} follows from
\cref{cusp-magnetic-line-bundle-proposition,cusp-magnetic-wedge-norm}.
The diagonal endomorphism \(H_{\mathrm{mag}}\) is a
globally defined parallel section of \(\operatorname{ad}\mathscr E_R\),
because every transition \(G_n\) commutes with it.  The forms
\(\theta_R^j\) are parallel constant forms for the descended flat metric.
Thus \(\nabla^R F_{\nabla^R}=0\), which implies
\(\mathrm d_{\nabla^R}^*F_{\nabla^R}=0\).  The curvature is nonzero because
\(D_R>0\).

The chosen Lie-algebra normalization gives
\[
 \lvert H_{\mathrm{mag}}\rvert^2
 =-\frac12\operatorname{Tr}(H_{\mathrm{mag}}^2)=1.
\]
Equations \cref{cusp-magnetic-wedge-norm,cusp-spec-exact-volume} therefore
give
\[
 \int_{F_R}\lvert F_{\nabla^R}\rvert^2\,
   \mathrm d\operatorname{vol}_{F_R}
 =4\pi^2D_R\frac1{D_R^2}=\frac{4\pi^2}{D_R},
\]
which proves \cref{cusp-magnetic-exact-action}.  Finally,
\cref{cusp-spec-D-expansion} and the elementary reciprocal expansion
\[
 \frac1{R^2+aT+O(1)}
 =\frac1{R^2}\left(1-\frac aT+O(R^{-2})\right)
\]
with \(a=2H-B\) prove
\cref{cusp-magnetic-action-first-correction}.
\end{proof}

\subsection{Constant-curvature magnetic Laplacians for arbitrary integral
nondegenerate flux}

We next give the Hilbert space, closed form, self-adjoint realization, and
oscillator spectrum in a form valid for every nondegenerate integral
alternating form on the marked four-torus.

Let
\[
 \mathsf E=\sum_{1\le j<k\le4}E_{jk}\,\mathrm du_j\wedge\mathrm du_k
 \in\bigwedge^2(\mathbb Z^4)^*,
 \qquad E_{jk}\in\mathbb Z,
\]
and let the same symbol denote its skew matrix.  For \(n\in\mathbb Z^4\),
put
\begin{equation}
 J_n^{\mathsf E}(u)
 =\exp\!\left(-2\pi i\sum_{j<k}E_{jk}n_ju_k\right).
 \label{cusp-magnetic-general-factor}
\end{equation}
The integer in
\[
 \sum_{j<k}E_{jk}m_jn_k
\]
shows directly that
\(J_m^{\mathsf E}(u+n)J_n^{\mathsf E}(u)=J_{m+n}^{\mathsf E}(u)\).
Thus \eqref{cusp-magnetic-general-factor} defines a Hermitian line bundle
\(\mathscr L_{\mathsf E,R}\).  On its pullback, set
\begin{equation}
 \nabla^{\mathsf E}
 =\mathrm d+2\pi i\sum_{j<k}E_{jk}u_j\,\mathrm du_k.
 \label{cusp-magnetic-general-connection}
\end{equation}
Exactly as in \cref{cusp-magnetic-line-bundle-proposition}, the
connection descends and
\begin{equation}
 F_{\nabla^{\mathsf E}}=2\pi i\,\mathsf E.
 \label{cusp-magnetic-general-curvature}
\end{equation}

In \(u\)-coordinates the flat metric matrix is
\begin{equation}
 G_R=P_R^{\mathsf T}P_R,
 \qquad
 \sqrt{\det G_R}=D_R.
 \label{cusp-magnetic-flat-Gram}
\end{equation}
Define the Hilbert space concretely by
\begin{align}
 \mathcal H_{\mathsf E,R}
 =\{\,&f\in L^2_{\mathrm{loc}}(\mathbb R^4):
 f(u+n)=J_n^{\mathsf E}(u)f(u)\text{ for every }n\in\mathbb Z^4\},
 \label{cusp-magnetic-Hilbert-space}\\
 \lVert f\rVert^2
 ={}&\int_{[0,1]^4}\lvert f(u)\rvert^2\sqrt{\det G_R}\,\mathrm d^4u.
 \nonumber
\end{align}
Let \(\mathcal C_{\mathsf E,R}\) be the smooth functions with the same
automorphy.  On this dense subspace define
\begin{equation}
 q_{\mathsf E,R}[f]
 =\int_{[0,1]^4}
  (G_R^{-1})^{jk}
  \overline{\nabla^{\mathsf E}_jf}
  \nabla^{\mathsf E}_kf
  \sqrt{\det G_R}\,\mathrm d^4u.
 \label{cusp-magnetic-quadratic-form}
\end{equation}

\begin{theorem}[Friedrichs realization and oscillator spectrum]
\label{cusp-magnetic-oscillator-theorem}
The form \eqref{cusp-magnetic-quadratic-form} is nonnegative and closable.
Its closed form domain is
\[
 \operatorname{Dom}\overline q_{\mathsf E,R}
 =H^1(F_R;\mathscr L_{\mathsf E,R}).
\]
Its Friedrichs operator is the
magnetic Laplacian
\[
 \Delta_{\mathsf E,R}=(\nabla^{\mathsf E})^*\nabla^{\mathsf E},
 \qquad
 \operatorname{Dom}(\Delta_{\mathsf E,R})
 =H^2(F_R;\mathscr L_{\mathsf E,R}),
\]
and it has compact resolvent.

Assume now that \(\mathsf E\) is nondegenerate.  Let \(b_1,b_2>0\) be
defined by the eigenvalues
\begin{equation}
 \operatorname{Spec}\!\left(
  G_R^{-1/2}(2\pi\mathsf E)G_R^{-1/2}\right)
 =\{ib_1,-ib_1,ib_2,-ib_2\}.
 \label{cusp-magnetic-frequency-definition}
\end{equation}
Then, as a set (with finite multiplicities),
\begin{equation}
 \operatorname{Spec}(\Delta_{\mathsf E,R})
 =\left\{(2n_1+1)b_1+(2n_2+1)b_2:
 n_1,n_2\in\mathbb Z_{\ge0}\right\}.
 \label{cusp-magnetic-oscillator-spectrum}
\end{equation}
In particular, the ground-to-first-distinct-level spacing is
\begin{equation}
 \operatorname{gap}(\Delta_{\mathsf E,R})=2\min(b_1,b_2).
 \label{cusp-magnetic-level-gap}
\end{equation}
No multiplicity formula is needed for, or asserted in, these identities.
\end{theorem}

\begin{proof}
In local unitary trivializations, the integrand in
\eqref{cusp-magnetic-quadratic-form} is the squared norm of a first-order
unitary covariant derivative.  A finite trivializing cover and a partition
of unity show that the form norm
\(\lVert f\rVert^2+q_{\mathsf E,R}[f]\) is equivalent to the usual local \(H^1\)
norm.  Its completion is therefore
\(H^1(F_R;\mathscr L_{\mathsf E,R})\), proving closability and identifying
the closed domain.  The representation theorem for closed nonnegative
forms gives the Friedrichs operator.  Its local principal symbol is
\(G_R^{-1}(\xi,\xi)\), so it is uniformly elliptic.  Standard interior
elliptic regularity on the compact manifold without boundary identifies
the operator domain with \(H^2\).  The compact inclusion
\(H^1\hookrightarrow L^2\) implies compact resolvent.

The real skew-normal-form theorem supplies a constant orthonormal frame
\(v_1,\ldots,v_4\) on the Euclidean cover in which the real two-form
\(2\pi\mathsf E\) is
\[
 b_1v^1\wedge v^2+b_2v^3\wedge v^4.
\]
Changing the sign of \(v^2\) or \(v^4\), if necessary, makes both
frequencies positive.  Let \(\nabla_a\) denote covariant differentiation
along \(v_a\).  Each \(\nabla_a\) is formally skew-adjoint on the common
invariant core \(\mathcal C_{\mathsf E,R}\); the closures used below are the
ones determined by the closed magnetic form.  On that smooth core one has
\[
 [\nabla_1,\nabla_2]=ib_1,
 \qquad
 [\nabla_3,\nabla_4]=ib_2,
\]
and all other commutators vanish.  Put
\[
 C_1=\nabla_1-i\nabla_2,
 \qquad
 C_2=\nabla_3-i\nabla_4,
 \qquad
 \Delta_j=-(\nabla_{2j-1}^2+\nabla_{2j}^2).
\]
Direct multiplication, retaining the commutator signs, gives
\begin{equation}
 C_j^*C_j=\Delta_j-b_j,
 \qquad
 C_jC_j^*=\Delta_j+b_j,
 \qquad
 [\Delta_{\mathsf E,R},C_j]=-2b_jC_j.
 \label{cusp-magnetic-ladder-identities}
\end{equation}

Because the resolvent is compact, the lowest eigenvalue has a smooth
eigenfunction \(\psi_0\neq0\).  If \(C_j\psi_0\neq0\), the last identity
in \eqref{cusp-magnetic-ladder-identities} would make it an eigenfunction
with a strictly smaller eigenvalue.  Hence \(C_1\psi_0=C_2\psi_0=0\), and
the first identity gives
\[
 \Delta_{\mathsf E,R}\psi_0=(b_1+b_2)\psi_0.
\]
Conversely, repeated application of \(C_1^*,C_2^*\) produces nonzero
eigenfunctions with eigenvalues
\((2n_1+1)b_1+(2n_2+1)b_2\).  Nonvanishing follows from
\(C_jC_j^*=C_j^*C_j+2b_j\).

It remains to exclude other eigenvalues.  For an eigenfunction \(\psi\)
with eigenvalue \(\lambda>b_1+b_2\), the operator identity
\[
 (\lambda-b_1-b_2)\psi
 =C_1^*C_1\psi+C_2^*C_2\psi
\]
expresses \(\psi\) as a sum of raising operators applied to the lower
eigenfunctions \(C_j\psi\).  Induction over the finite set of eigenvalues
below \(\lambda\) reduces every eigenfunction to raised ground states.
This proves \eqref{cusp-magnetic-oscillator-spectrum}.  Raising once in a
least-frequency plane proves \eqref{cusp-magnetic-level-gap}.
\end{proof}

\begin{corollary}[Frequency product and cusp spacing bound]
\label{cusp-magnetic-frequency-product}
Let
\[
 d=|\operatorname{Pf}(\mathsf E)|
 =\lvert E_{12}E_{34}-E_{13}E_{24}+E_{14}E_{23}\rvert>0.
\]
Then
\begin{equation}
 b_1b_2=\frac{4\pi^2d}{D_R},
 \label{cusp-magnetic-product-identity}
\end{equation}
and
\begin{align}
 \operatorname{gap}(\Delta_{\mathsf E,R})
 &\le 4\pi\sqrt{\frac d{D_R}}
 \label{cusp-magnetic-gap-bound}\\
 &=\frac{4\pi\sqrt d}{R}
 \left(1+\frac{B-2H}{2R}
 +O(R^{-2})\right).
 \nonumber
\end{align}
\end{corollary}

\begin{proof}
The determinant of the skew matrix in
\eqref{cusp-magnetic-frequency-definition} is both
\(b_1^2b_2^2\) and
\[
 \frac{\det(2\pi\mathsf E)}{\det G_R}
 =\frac{(2\pi)^4\operatorname{Pf}(\mathsf E)^2}{D_R^2}.
\]
Taking the positive square root proves
\eqref{cusp-magnetic-product-identity}.  Since
\(\min(b_1,b_2)\le\sqrt{b_1b_2}\),
\Cref{cusp-magnetic-oscillator-theorem} gives
\eqref{cusp-magnetic-gap-bound}.  Expanding the square root of
\cref{cusp-spec-D-expansion} gives the displayed first correction.
\end{proof}

\subsection{A background-gauge negative mode}

We finish with the classical background-gauge Jacobi operator at the
reducible critical connection.  Every sign below uses exactly the inner
product and the element \(H_{\mathrm{mag}}\) fixed in
\cref{cusp-magnetic-su2-normalization}.

Put
\[
 \mathscr H_R^{\mathrm{ad},1}
 :=L^2\Omega^1(F_R;\operatorname{ad}\mathscr E_R).
\]
Put also
\[
 \mathscr H_R^{\mathrm{ad},0}
 :=L^2\Omega^0(F_R;\operatorname{ad}\mathscr E_R),
 \qquad
 D_R^{\max}:=\mathrm d_{\nabla^R}^{*,\max}.
\]
For the resolvent argument these real Hilbert spaces are complexified;
all operators commute with the real structure, so their restrictions to
real sections are unchanged.
Here the maximal distributional adjoint is the closed, densely defined
operator whose domain and value are specified without an implicit
regularity assumption by
\begin{align}
 \operatorname{Dom}(D_R^{\max})
 :=\{a\in\mathscr H_R^{\mathrm{ad},1}:{}&\text{ there is a }u\in
 \mathscr H_R^{\mathrm{ad},0}\text{ such that}\notag\\[-1mm]
 &\langle a,\mathrm d_{\nabla^R}\phi\rangle_{L^2}
 =\langle u,\phi\rangle_{L^2}
 \text{ for every }\phi\in
 C^\infty(F_R;\operatorname{ad}\mathscr E_R)\},
 \label{cusp-magnetic-maximal-codifferential-domain}\\
 D_R^{\max}a:={}&u.\notag
\end{align}
The test space is dense, so \(u\) is unique.  The definition also proves
closedness directly: if \(a_n\to a\) and \(D_R^{\max}a_n\to u\) in
\(L^2\), passage to the limit in the displayed testing identity gives
\(a\in\operatorname{Dom}(D_R^{\max})\) and \(D_R^{\max}a=u\).
Consequently the background slice
\begin{equation}
 \mathscr S_R
 :=\ker D_R^{\max}
 \subset\mathscr H_R^{\mathrm{ad},1}.
 \label{cusp-magnetic-closed-background-slice}
\end{equation}
is a closed subspace.

The scalar covariant Laplacian used below is fixed as follows.  On
\(C^\infty(F_R;\operatorname{ad}\mathscr E_R)\) start with the
nonnegative form
\[
 q_R^{(0)}[\phi,\psi]
 :=\langle\mathrm d_{\nabla^R}\phi,
             \mathrm d_{\nabla^R}\psi\rangle_{L^2}.
\]
Its closure has domain
\(H^1\Omega^0(F_R;\operatorname{ad}\mathscr E_R)\), and its associated
Friedrichs operator is
\begin{equation}
 \Delta_{\nabla^R}^{(0)}
 :=\mathrm d_{\nabla^R}^*\mathrm d_{\nabla^R},
 \qquad
 \operatorname{Dom}(\Delta_{\nabla^R}^{(0)})
 =H^2\Omega^0(F_R;\operatorname{ad}\mathscr E_R).
 \label{cusp-magnetic-scalar-covariant-laplacian-domain}
\end{equation}
The \(H^2\) domain follows from elliptic regularity on the compact flat
torus; in particular, this notation always denotes the nonnegative
self-adjoint realization, not merely its differential expression.

For an adjoint-valued one-form \(a\), define the smooth bundle
endomorphism
\[
 (V_Ra)_\nu:=-2\sum_{\mu=1}^4[F_{\nu\mu},a_\mu].
\]
It is symmetric: after interchanging \(\mu,\nu\), the identities
\(F_{\mu\nu}=-F_{\nu\mu}\) and skew-adjointness of
\(\operatorname{ad}(F_{\nu\mu})\) give
\(\langle V_Ra,b\rangle=\langle a,V_Rb\rangle\).
On the full one-form Hilbert space define
\begin{equation}
 \mathfrak j_R^{\mathrm{gf}}[a,b]
 :=\langle\nabla^Ra,\nabla^Rb\rangle_{L^2}
   +\langle V_Ra,b\rangle_{L^2},
 \qquad
 \operatorname{Dom}(\mathfrak j_R^{\mathrm{gf}})
 =H^1\Omega^1(F_R;\operatorname{ad}\mathscr E_R).
 \label{cusp-magnetic-full-jacobi-form}
\end{equation}
Because \(V_R\) is a bounded symmetric endomorphism on the compact
torus, this form is closed and bounded below.  Its associated
self-adjoint operator \(\mathcal J_R^{\mathrm{gf}}\) has
\begin{equation}
 \operatorname{Dom}(\mathcal J_R^{\mathrm{gf}})
 =H^2\Omega^1(F_R;\operatorname{ad}\mathscr E_R)
 \label{cusp-magnetic-full-jacobi-domain}
\end{equation}
and differential expression
\[
 (\mathcal J_R^{\mathrm{gf}}a)_\nu
 =-\sum_{\mu=1}^4\nabla_{e_\mu}\nabla_{e_\mu}a_\nu
  -2\sum_{\mu=1}^4[F_{\nu\mu},a_\mu].
\]

\begin{proposition}[Intertwining and the reducing Coulomb slice]
\label{cusp-magnetic-jacobi-reduction-proposition}
At the Yang--Mills connection \(\nabla^R\), the following formal
identities hold on smooth sections:
\begin{equation}
 \mathcal J_R^{\mathrm{gf}}\mathrm d_{\nabla^R}
 =\mathrm d_{\nabla^R}\Delta_{\nabla^R}^{(0)},
 \qquad
 \mathrm d_{\nabla^R}^*\mathcal J_R^{\mathrm{gf}}
 =\Delta_{\nabla^R}^{(0)}\mathrm d_{\nabla^R}^*.
 \label{cusp-magnetic-jacobi-intertwining}
\end{equation}
For every \(z\in\mathbb C\setminus\mathbb R\), the correctly typed
closed-operator resolvent identity is
\begin{equation}
 D_R^{\max}(\mathcal J_R^{\mathrm{gf}}-z)^{-1}f
 =(\Delta_{\nabla^R}^{(0)}-z)^{-1}D_R^{\max}f,
 \qquad f\in\operatorname{Dom}(D_R^{\max}).
 \label{cusp-magnetic-jacobi-resolvent-intertwining}
\end{equation}
In particular, \(\mathscr S_R\) reduces
\(\mathcal J_R^{\mathrm{gf}}\).  The part of the full operator in this
subspace and its closed form are therefore
\begin{equation}
 \mathcal J_R
 :=\mathcal J_R^{\mathrm{gf}}\big|_{\mathscr S_R},
 \quad
 \operatorname{Dom}(\mathcal J_R)
 =H^2\Omega^1(F_R;\operatorname{ad}\mathscr E_R)\cap\mathscr S_R,
 \quad
 \operatorname{Dom}(\mathfrak j_R)
 =H^1\Omega^1(F_R;\operatorname{ad}\mathscr E_R)\cap\mathscr S_R
 \label{cusp-magnetic-reduced-jacobi-domain}
\end{equation}
with
\(\mathfrak j_R=\mathfrak j_R^{\mathrm{gf}}|_{\operatorname{Dom}(\mathfrak j_R)}\).
The operator \(\mathcal J_R\) is self-adjoint on \(\mathscr S_R\), and
\(\mathfrak j_R\) is its closed, lower-bounded form.
\end{proposition}

\begin{proof}
Write \(\nabla_\mu=\nabla^R_{e_\mu}\) in the constant orthonormal frame,
so
\([\nabla_\mu,\nabla_\nu]\phi=[F_{\mu\nu},\phi]\).
For a smooth adjoint-valued scalar \(\phi\), commuting the derivatives
without suppressing either curvature contribution gives
\begin{align*}
 &\bigl(\mathcal J_R^{\mathrm{gf}}\mathrm d_{\nabla^R}\phi
       -\mathrm d_{\nabla^R}\Delta_{\nabla^R}^{(0)}\phi\bigr)_\nu\\
 &\quad=-\nabla_\mu\nabla_\mu\nabla_\nu\phi
         +\nabla_\nu\nabla_\mu\nabla_\mu\phi
         -2[F_{\nu\mu},\nabla_\mu\phi]\\
 &\quad=-[\nabla_\mu F_{\mu\nu},\phi]
   =[ (\mathrm d_{\nabla^R}^*F_{\nabla^R})_\nu,\phi].
\end{align*}
The last member vanishes by
\cref{cusp-magnetic-ym-critical-point}.  This proves the first identity
in \eqref{cusp-magnetic-jacobi-intertwining}.  Since both Laplacians and
\(\mathcal J_R^{\mathrm{gf}}\) are formally symmetric, pairing it against
arbitrary smooth zero- and one-forms proves the second identity.

Fix \(z\notin\mathbb R\), take
\(f\in\operatorname{Dom}(D_R^{\max})\), and put
\(u=(\mathcal J_R^{\mathrm{gf}}-z)^{-1}f\).
The full operator domain gives \(u\in H^2\), hence
\(D_R^{\max}u\in H^1\).  Moreover
\((\mathcal J_R^{\mathrm{gf}}-z)u=f\) and the distributional extension of
the second formal identity imply
\begin{equation}
 (\Delta_{\nabla^R}^{(0)}-z)D_R^{\max}u=D_R^{\max}f
 \quad\text{in }\mathcal D'(F_R;\operatorname{ad}\mathscr E_R).
 \label{cusp-magnetic-resolvent-distributional-equation}
\end{equation}
The right side is in \(L^2\).  Elliptic regularity for the scalar
Laplacian upgrades \(D_R^{\max}u\) from \(H^1\) to \(H^2\), so it belongs
to \(\operatorname{Dom}(\Delta_{\nabla^R}^{(0)})\).  Applying the bounded
scalar resolvent to
\eqref{cusp-magnetic-resolvent-distributional-equation} proves
\eqref{cusp-magnetic-jacobi-resolvent-intertwining}; no extension of
\(D_R^{\max}\) to all of \(L^2\) is asserted or needed.

If \(f\in\mathscr S_R\), the right side of
\eqref{cusp-magnetic-jacobi-resolvent-intertwining} is zero, hence both
resolvents at \(z\) and at \(\overline z\) preserve \(\mathscr S_R\).
For \(g\in\mathscr S_R^\perp\) and \(s\in\mathscr S_R\), self-adjointness
of the full operator gives
\[
 \langle(\mathcal J_R^{\mathrm{gf}}-z)^{-1}g,s\rangle
 =\langle g,(\mathcal J_R^{\mathrm{gf}}-\overline z)^{-1}s\rangle=0.
\]
Thus the resolvent preserves both \(\mathscr S_R\) and its orthogonal
complement, which is precisely the reducing property.  The standard
part of a self-adjoint operator on a reducing subspace is self-adjoint,
and its domain is the intersection in
\eqref{cusp-magnetic-reduced-jacobi-domain}.  The reducing orthogonal
projection commutes with the spectral calculus of the lower-bounded full
operator, including its form square root; this proves the asserted form
domain and completes the proof.
\end{proof}

\begin{theorem}[Exact charged Landau negative mode]
\label{cusp-magnetic-negative-mode}
For the self-adjoint Yang--Mills Jacobi operator \(\mathcal J_R\) on the
closed background slice \(\mathscr S_R\) defined above, the differential
expression in a constant orthonormal frame \((e_\mu)_{\mu=1}^4\) is
\begin{equation}
 (\mathcal J_Ra)_\nu
 =-\sum_{\mu=1}^4\nabla_{e_\mu}\nabla_{e_\mu}a_\nu
  -2\sum_{\mu=1}^4[F_{\nu\mu},a_\mu].
 \label{cusp-magnetic-jacobi-operator}
\end{equation}
There exists a nonzero smooth real
\(\mathfrak{su}(2)\)-valued one-form \(a\) such that
\begin{equation}
 \lVert a\rVert_{L^2}=1,
 \qquad
 \mathrm d_{\nabla^R}^*a=0,
 \qquad
 \mathcal J_Ra=-\frac{4\pi}{D_R}a.
 \label{cusp-magnetic-negative-eigen-equation}
\end{equation}
Equivalently, for the action normalization
\cref{cusp-magnetic-action-normalization},
\begin{equation}
 \operatorname{YM}_R''\big|_{\nabla^R}(a,a)
 =-\frac{4\pi}{g_{\mathrm{YM}}^2D_R}.
 \label{cusp-magnetic-negative-hessian}
\end{equation}
Thus the reducible non-flat critical connection has a strictly negative
classical second-variation direction for every finite \(R\).
\end{theorem}

\begin{proof}
First isolate the exact two-dimensional magnetic problem.  Let
\[
 W=\{x_1=x_2=0\}\subset\mathbb R^4,
\]
so that \((y_1,y_2)\) are oriented orthonormal coordinates on \(W\).
By \cref{cusp-spec-real-period-matrix}, orthogonal projection to \(W\)
sends the first two marked lattice generators to
\[
 v_1=(6m_R,-q_R),
 \qquad
 v_2=(L_R,m_R),
\]
and sends the last two to zero.  Since
\(\det(v_1,v_2)=D_R>0\), projection descends to
\[
 \varpi_R:F_R\longrightarrow
 \Sigma_R:=W/(\mathbb Zv_1+\mathbb Zv_2).
\]
The factors and connection of \(\mathscr L_R^2\) depend only on
\((u_1,u_2)\).  They therefore define a Hermitian line bundle with unitary
connection \((\mathscr M_R,\nabla^{\mathscr M_R})\) on \(\Sigma_R\) such
that
\begin{equation}
 (\mathscr L_R^2,\nabla^{\mathscr L_R^2})
 =\varpi_R^*(\mathscr M_R,\nabla^{\mathscr M_R}).
 \label{cusp-magnetic-descended-root-line}
\end{equation}
Let \(\Delta_+=(\nabla^{\mathscr M_R})^*\nabla^{\mathscr M_R}\) be the
nonnegative self-adjoint Bochner Laplacian on
\(L^2(\Sigma_R;\mathscr M_R)\), defined by its closed \(H^1\) form and
with operator domain \(H^2(\Sigma_R;\mathscr M_R)\).  More precisely, if
\(c_R>0\) is the constant Riemannian volume of a fibre of \(\varpi_R\),
then
\[
 U_R:L^2(\Sigma_R;\mathscr M_R)\longrightarrow
       L^2(F_R;\mathscr L_R^2),
 \qquad U_Rs=c_R^{-1/2}\varpi_R^*s,
\]
is an isometry.  Its image
\[
 \mathscr B_R:=U_RL^2(\Sigma_R;\mathscr M_R)
\]
is the closed fibre-basic subspace.  The pulled-back realization is,
by definition,
\[
 \Delta_{+,\mathrm{bas}}:=U_R\Delta_+U_R^{-1}
 \quad\text{on }\mathscr B_R,
 \qquad
 \operatorname{Dom}(\Delta_{+,\mathrm{bas}})
 =U_RH^2(\Sigma_R;\mathscr M_R),
\]
and its closed form domain is \(U_RH^1(\Sigma_R;\mathscr M_R)\).
Only this operator on \(\mathscr B_R\) is self-adjoint.  In the formulas
below we abbreviate \(\Delta_{+,\mathrm{bas}}\) again to \(\Delta_+\);
no self-adjoint operator on the full
\(L^2(F_R;\mathscr L_R^2)\) is thereby defined or asserted.  On this
root line the curvature is
\begin{equation}
 F_{\mathscr L_R^2}
 =4\pi i\,\theta_R^1\wedge\theta_R^2
 =i\kappa_R\,\mathrm dy_1\wedge\mathrm dy_2,
 \qquad
 \kappa_R=\frac{4\pi}{D_R}>0.
 \label{cusp-magnetic-root-field-strength}
\end{equation}

We now construct an exact lowest Landau section on this actual projected
lattice.  In the marked coordinates put
\[
 \mathcal D_1:=\partial_{u_1},
 \qquad
 \mathcal D_2:=\partial_{u_2}+4\pi i u_1,
\]
and define
\begin{equation}
 \zeta_R:=\frac{q_R-6im_R}{m_R+iL_R},
 \qquad
 v_R:=-\operatorname{Im}\zeta_R
     =\frac{D_R}{m_R^2+L_R^2}>0.
 \label{cusp-magnetic-theta-parameter}
\end{equation}
Indeed,
\[
 \operatorname{Im}\bigl((q_R-6im_R)(m_R-iL_R)\bigr)
 =-(L_Rq_R+6m_R^2)=-D_R,
\]
which also checks the sign \(v_R>0\).
The series
\begin{equation}
 \Theta_R(u_1,u_2)
 =\sum_{\ell\in\mathbb Z}
 \exp\!\left(
  -2\pi i\zeta_R(\ell+u_1)^2+4\pi i\ell u_2
 \right)
 \label{cusp-magnetic-theta-series}
\end{equation}
converges normally with all derivatives, because the absolute value of
its \(\ell\)-th summand is
\(\exp(-2\pi v_R(\ell+u_1)^2)\).  Reindexing the series and using
\(\ell\in\mathbb Z\) give
\begin{equation}
 \Theta_R(u_1+1,u_2)=e^{-4\pi i u_2}\Theta_R(u_1,u_2),
 \qquad
 \Theta_R(u_1,u_2+1)=\Theta_R(u_1,u_2).
 \label{cusp-magnetic-theta-automorphy}
\end{equation}
Thus \(\Theta_R\) is a smooth section of \(\mathscr L_R^2\).  Termwise
differentiation gives
\begin{equation}
 (\mathcal D_1+\zeta_R\mathcal D_2)\Theta_R=0.
 \label{cusp-magnetic-theta-marked-equation}
\end{equation}
The exact period block on the \((y_1,y_2)\)-plane is
\[
 \begin{pmatrix}6m_R&L_R\\-q_R&m_R\end{pmatrix}.
\]
Its inverse transpose gives the physical covariant derivatives
\begin{equation}
 \nabla_{y_1}^+
  =\frac{m_R\mathcal D_1+q_R\mathcal D_2}{D_R},
 \qquad
 \nabla_{y_2}^+
  =\frac{-L_R\mathcal D_1+6m_R\mathcal D_2}{D_R}.
 \label{cusp-magnetic-physical-root-derivatives}
\end{equation}
Consequently
\begin{equation}
 \nabla_{y_1}^+-i\nabla_{y_2}^+
 =\frac{m_R+iL_R}{D_R}
   (\mathcal D_1+\zeta_R\mathcal D_2).
 \label{cusp-magnetic-annihilator-conversion}
\end{equation}
Orthogonality of the \(u_2\)-Fourier modes also gives the nonzero norm
\[
 \int_0^1\!\!\int_0^1\lvert\Theta_R\rvert^2
 \,\mathrm du_1\,\mathrm du_2
 =\sum_{\ell\in\mathbb Z}\int_0^1
   e^{-4\pi v_R(\ell+u_1)^2}\,\mathrm du_1
 =\frac1{2\sqrt{v_R}}.
\]
Set
\[
 \phi_R:=\left(\frac{2\sqrt{v_R}}{D_R}\right)^{1/2}\Theta_R
\]
and pull it back by \(\varpi_R\).  Since the marked-coordinate volume
element is \(D_R\,\mathrm d^4u\), it has unit \(L^2(F_R)\)-norm.  Its
covariant derivatives in the \(x_1,x_2\) directions vanish, and
\begin{equation}
 (\nabla_{y_1}^+-i\nabla_{y_2}^+)\phi_R=0,
 \qquad
 \Delta_+\phi_R=\kappa_R\phi_R,
 \qquad
 \Delta_+\phi_R
 =-\bigl((\nabla_{y_1}^+)^2+(\nabla_{y_2}^+)^2\bigr)\phi_R.
 \label{cusp-magnetic-ground-equation}
\end{equation}
Indeed, \cref{cusp-magnetic-physical-root-derivatives} and
\([\mathcal D_1,\mathcal D_2]=4\pi i\) give
\[
 [\nabla_{y_1}^+,\nabla_{y_2}^+]
 =\frac{m_R(6m_R)-q_R(-L_R)}{D_R^2}\,4\pi i
 =i\kappa_R.
\]
Direct multiplication then gives
\[
 (\nabla_{y_1}^+-i\nabla_{y_2}^+)^*
 (\nabla_{y_1}^+-i\nabla_{y_2}^+)=\Delta_+-\kappa_R,
\]
and the second equation in \eqref{cusp-magnetic-ground-equation} follows
from the first.

In \(\mathfrak{sl}_2(\mathbb C)\), put
\[
 X_+=\begin{pmatrix}0&1\\0&0\end{pmatrix},
 \qquad
 X_-=\begin{pmatrix}0&0\\1&0\end{pmatrix}.
\]
Then
\[
 [H_{\mathrm{mag}},X_+]=2iX_+,
 \qquad
 [H_{\mathrm{mag}},X_-]=-2iX_-.
\]
The \(X_+\)-root line in
\(\operatorname{ad}(\mathscr L_R\oplus\mathscr L_R^{-1})
 \otimes\mathbb C\) is \(\mathscr L_R^2\),
so
\[
 a^+=\phi_RX_+(\mathrm dy_1-i\,\mathrm dy_2)
\]
is a globally defined complexified adjoint-valued one-form.  The first
equation in \eqref{cusp-magnetic-ground-equation} gives
\[
 \mathrm d_{\nabla^R}^*a^+
 =-(\nabla_{y_1}^+-i\nabla_{y_2}^+)\phi_R\,X_+=0.
\]

Put \(\mathcal B_R:=2\pi/D_R\), so that
\(\kappa_R=2\mathcal B_R\) and
\(F_{y_1y_2}=\mathcal B_RH_{\mathrm{mag}}\).  The two nonzero components
of \(a^+\) are \(a_{y_1}^+=\phi_RX_+\) and
\(a_{y_2}^+=-i\phi_RX_+\).  Substitution into
\eqref{cusp-magnetic-jacobi-operator}, with no suppressed charge factor,
gives
\begin{align*}
 (\mathcal J_Ra^+)_{y_1}
 &=\kappa_R\phi_RX_+
   -2[\mathcal B_RH_{\mathrm{mag}},-i\phi_RX_+]
  =-\kappa_R\phi_RX_+,\\
 (\mathcal J_Ra^+)_{y_2}
 &=-i\kappa_R\phi_RX_+
   -2[-\mathcal B_RH_{\mathrm{mag}},\phi_RX_+]
  =i\kappa_R\phi_RX_+.
\end{align*}
Thus \(\mathcal J_Ra^+=-\kappa_Ra^+\).

For completeness, formula \eqref{cusp-magnetic-jacobi-operator} follows by
expanding
\(F_{\nabla^R+ta}=F_{\nabla^R}+t\,\mathrm d_{\nabla^R}a+
\tfrac12t^2[a\wedge a]\), integrating the quadratic form by parts on
the background slice, and using the flat-base Weitzenbock identity.  The
Hodge Laplacian contributes
\(-[F_{\nu\mu},a_\mu]\), and the curvature term in the second variation
contributes the identical term, giving the coefficient \(-2\).

Finally take the anti-Hermitian real form
\begin{equation}
 a:=\frac1{\sqrt2}\bigl(a^+-(a^+)^\dagger\bigr)
  =\frac1{\sqrt2}\left(
   \phi_RX_+(\mathrm dy_1-i\,\mathrm dy_2)
   -\overline{\phi_R}X_-(\mathrm dy_1+i\,\mathrm dy_2)\right).
 \label{cusp-magnetic-real-negative-mode}
\end{equation}
It is nonzero, traceless, and anti-Hermitian.  For
\(X(z):=zX_+-\overline zX_-\), the fixed inner product gives
\(-\tfrac12\operatorname{Tr}(X(z)^2)=\lvert z\rvert^2\).  The two
orthonormal one-form components in
\cref{cusp-magnetic-real-negative-mode} therefore have total pointwise
norm \(\lvert a\rvert^2=\lvert\phi_R\rvert^2\), so
\(\lVert a\rVert_{L^2}=1\).  Complex conjugation and
Hermitian adjunction preserve both the gauge equation and the real Jacobi
operator, so \eqref{cusp-magnetic-negative-eigen-equation} follows from
the calculation for \(a^+\).  The Hessian identity follows from
\(\operatorname{YM}_R''(a,a)
=g_{\mathrm{YM}}^{-2}\langle a,\mathcal J_Ra\rangle_{L^2}\)
on the background slice.
\end{proof}



% Standalone-preprint adaptation of cusp_native_full_flux_magnetic.tex.
% Canonical source SHA-256: C31D956283574111C88B6C102C2BD8D8060FE46E263C7F05B98FEA232CD7857D.
% The canonical radial-height symbol T was renamed R consistently; transpose notation
% and the finite-monodromy generators T_1,T_2 were retained exactly.  The dual-action
% bridge to the root preprint's homology/angular monodromy matrices is explicit below.

\section{The native full-flux magnetic sector}
\label{cusp-native-full-flux-section}

This section retains the canonical radial parameter, period matrix, Euclidean
quotient metric, marked lattice, and every sign already fixed in
\cref{cusp-spec-section}.  It treats the monodromy-invariant integral class
that uses both elementary alternating blocks.  It is therefore distinct from
the unit-flux plane treated in \cref{cusp-magnetic-section}.  No metric
rescaling, finite covering, rectangular replacement of the lattice, or
orientation suppression occurs below.

For \(R>R_0\), retain
\begin{equation}
 \begin{gathered}
 L_R=R+H_R=\operatorname{Im}\tau_R,
 \qquad q_R=L_R-B_R=-\operatorname{Im}\beta_R,\\
 m_R=\operatorname{Im}\mu_R,
 \qquad D_R=L_Rq_R+6m_R^2>0,
 \end{gathered}
 \label{cusp-native-full-flux-canonical-scalars}
\end{equation}
and
write
\[
 \eta_R^{\mathrm{re}}:=\operatorname{Re}h_R.
\]
This scalar is denoted by \(\eta_R\) in
\cref{cusp-spec-section}; the superscript is retained throughout this
section so that it cannot be confused with the native integral two-form
introduced below.  The complete period matrix is
\begin{equation}
 P_R=
 \begin{pmatrix}
  6\alpha_R&\eta_R^{\mathrm{re}}&1&0\\
  6m_R&L_R&0&0\\
  c_R&\alpha_R&0&1\\
  -q_R&m_R&0&0
 \end{pmatrix},
 \qquad \det P_R=-D_R.
 \label{cusp-native-full-flux-period-matrix}
\end{equation}
Thus \(F_R=\mathbb R^4/P_R\mathbb Z^4\) has the unrescaled Euclidean quotient
metric and volume \(D_R\).  In the marking coordinates \(y\in\mathbb R^4\),
the map \(y\mapsto x=P_Ry\) pulls the Euclidean metric back to
\(G_R=P_R^{\mathsf T}P_R\).

The limiting cusp constants remain
\begin{equation}
 m=\operatorname{Im}\mu(0),\qquad
 H=\operatorname{Im}h(0),\qquad
 B=\operatorname{Im}b(0),
 \label{cusp-native-full-flux-limiting-constants}
\end{equation}
so holomorphy at \(e^{-2\pi R}=0\) gives
\begin{equation}
 \begin{aligned}
 m_R&=m+O(e^{-2\pi R}),&
 H_R&=H+O(e^{-2\pi R}),\\
 B_R&=B+O(e^{-2\pi R}),&
 D_R&=R^2+(2H-B)R\\
 &&&\quad +(H^2-BH+6m^2)+O(Re^{-2\pi R}).
 \end{aligned}
 \label{cusp-native-full-flux-cusp-expansions}
\end{equation}

\subsection{The invariant class, its cocycle, and both orientations}

The matrices \(A_1,A_2\) in \eqref{eq:monodromy-matrices} act on the marked
homology, equivalently on the angular lattice used there.  On the dual
integral cohomology lattice put
\begin{equation}
 T_j:=A_j^{-\mathsf T},\qquad j=1,2.
 \label{cusp-native-full-flux-dual-monodromy}
\end{equation}
Direct inversion of those displayed matrices gives the actions written
below.  In the resulting ordered cohomological marking
\(\mathbb Z\langle\gamma,u,w,\delta\rangle\), put
\begin{equation}
 \eta=u\wedge w+6\gamma\wedge\delta.
 \label{cusp-native-full-flux-class}
\end{equation}
The two finite monodromy generators act by
\begin{align*}
T_1\gamma&=\gamma,&T_1u&=-u-w,&
T_1w&=-6\gamma+u,&T_1\delta&=2\gamma+u+w+\delta,\\
T_2\gamma&=\gamma,&T_2u&=6\gamma+w,&
T_2w&=-u,&T_2\delta&=-3\gamma+u+\delta.
\end{align*}
Directly, without passing to a normalized representative,
\begin{align*}
(\bigwedge^2T_1)\eta
&=(-u-w)\wedge(-6\gamma+u)
 +6\gamma\wedge(2\gamma+u+w+\delta)\\
&=u\wedge w+6\gamma\wedge\delta,\\
(\bigwedge^2T_2)\eta
&=(6\gamma+w)\wedge(-u)
 +6\gamma\wedge(-3\gamma+u+\delta)\\
&=u\wedge w+6\gamma\wedge\delta.
\end{align*}
Hence \(\eta\) is invariant under both generators.

On \(Y=\mathbb R^4/\mathbb Z^4\), with ordered coordinates
\(y^1,y^2,y^3,y^4\), its constant representative and alternating matrix are
\begin{equation}
 \eta_R=\mathrm dy^2\wedge\mathrm dy^3
       +6\,\mathrm dy^1\wedge\mathrm dy^4,
 \qquad
 E_\eta=
 \begin{pmatrix}
 0&0&0&6\\
 0&0&1&0\\
 0&-1&0&0\\
 -6&0&0&0
 \end{pmatrix}.
 \label{cusp-native-full-flux-form-matrix}
\end{equation}
In the paired basis \((e_2,e_3;e_1,e_4)\), the alternating elementary
divisors are \((1,6)\); consequently
\begin{equation}
 \operatorname{SNF}(E_\eta)=\operatorname{diag}(1,1,6,6),
 \qquad \operatorname{Pf}(E_\eta)=6,
 \qquad \det E_\eta=36.
 \label{cusp-native-full-flux-smith-pfaffian}
\end{equation}

The actual Hermitian line bundle is the quotient of
\(\mathbb R^4\times\mathbb C\) by
\begin{align}
 (y,z)&\longmapsto(y+e_1,e^{12\pi i y^4}z),&
 (y,z)&\longmapsto(y+e_2,e^{2\pi i y^3}z),
 \label{cusp-native-full-flux-cocycle}\\
 (y,z)&\longmapsto(y+e_3,z),&
 (y,z)&\longmapsto(y+e_4,z).
 \notag
\end{align}
The only possibly nontrivial generator commutators acquire the factors
\(e^{12\pi i}=1\) and \(e^{2\pi i}=1\), so these maps define a
\(\mathbb Z^4\)-action.  Denote the quotient by
\(\mathscr L_{\eta,R}\to F_R\), using the marking
\([y]\mapsto[P_Ry]\).  The connection
\begin{equation}
 \nabla_{\eta,R}
 =\mathrm d-2\pi i(6y^1\,\mathrm dy^4+y^2\,\mathrm dy^3)
 \label{cusp-native-full-flux-connection}
\end{equation}
descends through the displayed cocycle.  For instance translation by \(e_1\)
changes its connection form by \(-12\pi i\,\mathrm dy^4\), which equals the
negative logarithmic derivative of \(e^{12\pi i y^4}\); the \(e_2\) check is
identical.  Therefore
\begin{equation}
 F_{\nabla_{\eta,R}}=-2\pi i\eta_R,
 \qquad c_1(\mathscr L_{\eta,R})
 =\frac{i}{2\pi}F_{\nabla_{\eta,R}}=\eta_R.
 \label{cusp-native-full-flux-curvature}
\end{equation}

There are two opposite orientations which must not be conflated.  Write
\begin{equation}
 \operatorname{or}_{\mathrm{mark}}
 =\mathrm dy^1\wedge\mathrm dy^2\wedge
  \mathrm dy^3\wedge\mathrm dy^4
 \label{cusp-native-full-flux-marking-orientation}
\end{equation}
and let \(\operatorname{or}_{x}\) be the native complex orientation of
\((z_1,z_2)=(x_1+iy_1,x_2+iy_2)\).  Then
\begin{equation}
 P_R^*\operatorname{or}_{x}
 =-D_R\operatorname{or}_{\mathrm{mark}},
 \label{cusp-native-full-flux-orientation-reversal}
\end{equation}
whereas
\begin{equation}
 \frac12\eta_R^2=6\operatorname{or}_{\mathrm{mark}}.
 \label{cusp-native-full-flux-square}
\end{equation}
Consequently
\begin{equation}
 \frac12\langle\eta_R^2,[F_R]_{\mathrm{mark}}\rangle=6,
 \qquad
 \frac12\langle\eta_R^2,[F_R]_{x\text{-}\mathrm{complex}}\rangle=-6.
 \label{cusp-native-full-flux-two-orientation-pairings}
\end{equation}

\subsection{The exact physical frequencies}

In the physical Euclidean coordinates the two-form matrix is
\begin{equation}
 \Omega_R=P_R^{-\mathsf T}E_\eta P_R^{-1}
 =\frac1{D_R}
 \begin{pmatrix}
 0&-q_R&0&-6m_R\\
 q_R&0&6m_R&0\\
 0&-6m_R&0&6L_R\\
 6m_R&0&-6L_R&0
 \end{pmatrix}.
 \label{cusp-native-full-flux-physical-matrix}
\end{equation}
After the exact permutation
\((x^1,x^2,x^3,x^4)\mapsto(x^1,x^3;x^2,x^4)\), its off-diagonal block is
\begin{equation}
 K_R=\begin{pmatrix}-q_R&-6m_R\\-6m_R&6L_R\end{pmatrix},
 \qquad \det K_R=-6D_R.
 \label{cusp-native-full-flux-K}
\end{equation}
Define the magnetic scalars
\begin{equation}
 S_R:=6L_R-q_R,
 \qquad
 \Xi_R:=\sqrt{(6L_R+q_R)^2+144m_R^2}.
 \label{cusp-native-full-flux-S-Xi}
\end{equation}
Here \(S_R\) has no argument and is distinct from the affine residual
\(S_R(k)\) used in the dual-lattice calculation of
\cref{cusp-spec-exact-dual-lattice}.  The signed eigenvalues of \(K_R\) are
\begin{equation}
 \kappa_{R,\pm}=\frac{S_R\pm\Xi_R}{2}.
 \label{cusp-native-full-flux-kappas}
\end{equation}
Direct expansion, with no asymptotic replacement, gives
\begin{equation}
 \Xi_R^2-S_R^2
 =24(L_Rq_R+6m_R^2)=24D_R>0.
 \label{cusp-native-full-flux-Xi-identity}
\end{equation}
Thus \(\kappa_{R,+}>0>\kappa_{R,-}\).  The two positive charge-one
magnetic frequencies are
\begin{equation}
 \boxed{
 \mathcal B_{R,+}=\frac{\pi(S_R+\Xi_R)}{D_R},
 \qquad
 \mathcal B_{R,-}=\frac{\pi(\Xi_R-S_R)}{D_R}.}
 \label{cusp-native-full-flux-frequencies}
\end{equation}
They obey the exact identities
\begin{align}
 \mathcal B_{R,+}\mathcal B_{R,-}
 &=\frac{24\pi^2}{D_R},
 \label{cusp-native-full-flux-frequency-product}\\
 \mathcal B_{R,+}+\mathcal B_{R,-}
 &=\frac{2\pi\Xi_R}{D_R},&
 \mathcal B_{R,+}-\mathcal B_{R,-}
 &=\frac{2\pi S_R}{D_R},
 \label{cusp-native-full-flux-frequency-sum-difference}\\
 \mathcal B_{R,+}&=\frac{24\pi}{\Xi_R-S_R},&
 \mathcal B_{R,-}&=\frac{24\pi}{\Xi_R+S_R}.
 \label{cusp-native-full-flux-rationalized-frequencies}
\end{align}
The last line follows by multiplying by the conjugate factor and using
\cref{cusp-native-full-flux-Xi-identity}.

Keeping the three independent limiting coefficients \(H,B,m\), one obtains
\begin{align}
 \mathcal B_{R,-}
 &=\frac{2\pi}{R}-\frac{2\pi H}{R^2}
 +\frac{2\pi}{R^3}\left(H^2-\frac67m^2\right)+O(R^{-4}),
 \label{cusp-native-full-flux-Bminus-expansion}\\
 \mathcal B_{R,+}
 &=\frac{12\pi}{R}+\frac{12\pi(B-H)}{R^2}
 +\frac{12\pi}{R^3}
 \left((B-H)^2-\frac{36}{7}m^2\right)+O(R^{-4}),
 \label{cusp-native-full-flux-Bplus-expansion}\\
 \rho_R:=\frac{\mathcal B_{R,+}}{\mathcal B_{R,-}}
 &=6+\frac{6B}{R}
 +\frac{6B^2-6BH-180m^2/7}{R^2}+O(R^{-3}).
 \label{cusp-native-full-flux-ratio-expansion}
\end{align}
These follow by substituting \cref{cusp-native-full-flux-cusp-expansions}
into \cref{cusp-native-full-flux-rationalized-frequencies} and expanding the
reciprocals; no sign of \(H,B,m\) is assumed.  In particular, on a sufficiently
deep cusp tail,
\begin{equation}
 S_R>0,
 \qquad \mathcal B_{R,+}>\mathcal B_{R,-}>0,
 \qquad 5<\rho_R<7.
 \label{cusp-native-full-flux-deep-tail-window}
\end{equation}

\subsection{Joint Landau levels and their exact multiplicities}

Let \(\Delta_{\eta,R}=\nabla_{\eta,R}^*\nabla_{\eta,R}\) be the Friedrichs
magnetic Laplacian on \(L^2(F_R,\mathscr L_{\eta,R})\).  Its form domain is
\(H^1(F_R;\mathscr L_{\eta,R})\), its operator domain is
\(H^2(F_R;\mathscr L_{\eta,R})\), and compactness of the torus gives compact
resolvent.  We now prove the multiplicity rather than infer it from the
Pfaffian alone.

Define \(A_R^{\mathrm{pol}}\) on the marked tangent space by
\begin{equation}
 G_R(A_R^{\mathrm{pol}}v,w)=2\pi E_\eta(v,w),
 \qquad
 Q_R^{\mathrm{pol}}=\bigl(-(A_R^{\mathrm{pol}})^2\bigr)^{1/2},
 \qquad
 J_{\eta,R}=A_R^{\mathrm{pol}}(Q_R^{\mathrm{pol}})^{-1}.
 \label{cusp-native-full-flux-polar-complex-structure}
\end{equation}
The endomorphism \(A_R^{\mathrm{pol}}\) is \(G_R\)-skew-adjoint and invertible,
so \(Q_R^{\mathrm{pol}}\) is positive, commutes with
\(A_R^{\mathrm{pol}}\), and
\begin{equation}
 J_{\eta,R}^2=-1,
 \qquad
 E_\eta(J_{\eta,R}v,J_{\eta,R}w)=E_\eta(v,w),
 \qquad
 2\pi E_\eta(v,J_{\eta,R}v)
 =G_R(Q_R^{\mathrm{pol}}v,v)>0.
 \label{cusp-native-full-flux-polar-properties}
\end{equation}
For every signed charge \(q\in\mathbb Z\setminus\{0\}\), use the convention
\begin{equation}
 \mathscr L_{\eta,R}^{\otimes q}=
 \begin{cases}
  (\mathscr L_{\eta,R})^{\otimes q},&q>0,\\
  (\mathscr L_{\eta,R}^{\vee})^{\otimes(-q)},&q<0,
 \end{cases}
 \qquad
 A_{q,R}^{\mathrm{pol}}=qA_R^{\mathrm{pol}},
 \qquad
 Q_{q,R}^{\mathrm{pol}}
 =\bigl(-(A_{q,R}^{\mathrm{pol}})^2\bigr)^{1/2}
 =|q|Q_R^{\mathrm{pol}},
 \label{cusp-native-full-flux-signed-charge-bundle}
\end{equation}
and define
\begin{equation}
 J_{q,R}=A_{q,R}^{\mathrm{pol}}(Q_{q,R}^{\mathrm{pol}})^{-1}
 =\operatorname{sgn}(q)J_{\eta,R}.
 \label{cusp-native-full-flux-signed-charge-complex-structure}
\end{equation}
All these coefficients are constant, hence the Nijenhuis tensor of
\(J_{q,R}\) vanishes and the marking transports it to a complex structure on
\(F_R\).  The induced connection \(\nabla_{\eta,R}^{(q)}\) has
\begin{equation}
 F_{\nabla_{\eta,R}^{(q)}}=-2\pi i q\eta_R,
 \qquad
 2\pi qE_\eta(v,J_{q,R}v)
 =|q|G_R(Q_R^{\mathrm{pol}}v,v)>0.
 \label{cusp-native-full-flux-signed-charge-positivity}
\end{equation}
Thus its curvature is of type \((1,1)\), and
\((\nabla_{\eta,R}^{(q)})^{0,1}\) makes
\(\mathscr L_{\eta,R}^{\otimes q}\) a positive Appell--Humbert line bundle on
the auxiliary complex torus \((F_R,J_{q,R})\).

The auxiliary orientation is fixed, rather than repaired by an absolute value.
For \(q>0\), compatibility identifies the \(J_{q,R}\)-complex orientation
with the \(E_\eta\)-symplectic orientation; for \(q<0\), replacing \(J\) by
\(-J\) preserves complex orientation in complex dimension two.  Hence
\begin{equation}
 [F_R]_{J_{q,R}}=[F_R]_{\mathrm{mark}}
 =-[F_R]_{x\text{-}\mathrm{complex}}.
 \label{cusp-native-full-flux-auxiliary-native-orientations}
\end{equation}
This auxiliary positive structure is not the native complex structure inherited
from the \(S^6\) construction.  Indeed, with \(E_\eta=\operatorname{Im}H_{\eta,R}\),
the native Appell--Humbert calculation is
\begin{equation}
 G_{\eta,R}=
 \begin{pmatrix}
  2L_R&2m_R\\
  2m_R&-q_R/3
 \end{pmatrix},
 \qquad
 \det G_{\eta,R}=-\frac23D_R<0,
 \qquad H_{\eta,R}=2G_{\eta,R}^{-1}.
 \label{cusp-native-full-flux-native-hermitian-signature}
\end{equation}
Consequently \(H_{\eta,R}\), and also \(qH_{\eta,R}\) for every nonzero
\(q\), has signature \((1,1)\).  The native line bundle is holomorphic because
\(q\eta_R\) is of type \((1,1)\), but it is indefinite and is not the positive
bundle just constructed.

The tangent and canonical bundles of the auxiliary complex torus are trivial.
Kodaira vanishing and Hirzebruch--Riemann--Roch therefore give
\begin{equation}
 \begin{aligned}
 \dim_{\mathbb C}H^0\bigl((F_R,J_{q,R}),
                   \mathscr L_{\eta,R}^{\otimes q}\bigr)
 &=\int_{[F_R]_{J_{q,R}}}\frac{(q\eta_R)^2}{2}
  =\int_{[F_R]_{\mathrm{mark}}}\frac{(q\eta_R)^2}{2}
  =6q^2,\\
 \dim_{\mathbb C}H^0\bigl((F_R,J_{\eta,R}),\mathscr L_{\eta,R}\bigr)
 &=6,\\
 \dim_{\mathbb C}H^0\bigl((F_R,J_{2,R}),
                   \mathscr L_{\eta,R}^{\otimes2}\bigr)
 &=24.
 \end{aligned}
 \label{cusp-native-full-flux-vacuum-dimensions}
\end{equation}
No absolute value occurs in this auxiliary Hirzebruch--Riemann--Roch integral.
Equivalently, the positive polarization has elementary divisors
\((|q|,6|q|)\), and the Appell--Humbert theta theorem gives their product
\(6q^2\); see
\href{https://math.univ-cotedazur.fr/~beauvill/pubs/thetaon.pdf}%
{Beauville, Theorems~2.6 (pp.~107--108) and~3.5 (p.~112)}.
On the native complex torus the same signed Hirzebruch--Riemann--Roch
calculation instead gives
\begin{equation}
 \chi\bigl((F_R,J_{\mathrm{native},R}),
           \mathscr L_{\eta,R}^{\otimes q}\bigr)
 =\int_{[F_R]_{x\text{-}\mathrm{complex}}}
       \frac{(q\eta_R)^2}{2}
 =-6q^2,
 \label{cusp-native-full-flux-native-signed-hrr}
\end{equation}
and positivity vanishing is unavailable there.  Thus the positive auxiliary
section count and the native signed Euler characteristic concern distinct
Dolbeault problems and distinct complex orientations.

Choose a \(G_R\)-orthonormal frame
\(e_+,f_+=J_{\eta,R}e_+,e_-,f_-=J_{\eta,R}e_-\) which diagonalizes
\(Q_R^{\mathrm{pol}}\), labelled so its positive eigenvalues are
\(\mathcal B_{R,+},\mathcal B_{R,-}\).  Put
\begin{equation}
 a_\sigma=
 \frac{\nabla_{e_\sigma}+i\nabla_{f_\sigma}}
      {\sqrt{2\mathcal B_{R,\sigma}}},
 \qquad
 a_\sigma^*=
 \frac{-\nabla_{e_\sigma}+i\nabla_{f_\sigma}}
      {\sqrt{2\mathcal B_{R,\sigma}}},
 \qquad \sigma\in\{+,-\}.
 \label{cusp-native-full-flux-ladders}
\end{equation}
The curvature commutators give
\begin{equation}
 [a_\sigma,a_\tau^*]=\delta_{\sigma\tau},
 \qquad [a_\sigma,a_\tau]=[a_\sigma^*,a_\tau^*]=0,
 \label{cusp-native-full-flux-CCR}
\end{equation}
and, for \(N_\sigma=a_\sigma^*a_\sigma\),
\begin{equation}
 \Delta_{\eta,R}
 =\mathcal B_{R,+}(2N_++1)+\mathcal B_{R,-}(2N_-+1).
 \label{cusp-native-full-flux-number-operator}
\end{equation}
The simultaneous vacuum
\(\ker a_+\cap\ker a_-=H^0((F_R,J_{\eta,R}),\mathscr L_{\eta,R})\)
has complex dimension six by
\cref{cusp-native-full-flux-vacuum-dimensions}.

\begin{theorem}[Complete native joint Landau decomposition]
\label{cusp-native-full-flux-landau-theorem}
For every \(n_+,n_-\in\mathbb Z_{\ge0}\), the joint eigenspace
\begin{equation}
 \mathcal H_{R;n_+,n_-}
 =\ker(N_+-n_+)\cap\ker(N_--n_-)
 \label{cusp-native-full-flux-joint-space}
\end{equation}
has complex dimension six and energy
\begin{equation}
 E_{R;n_+,n_-}
 =(2n_++1)\mathcal B_{R,+}
 +(2n_-+1)\mathcal B_{R,-}.
 \label{cusp-native-full-flux-joint-energy}
\end{equation}
Moreover
\begin{equation}
 L^2(F_R,\mathscr L_{\eta,R})
 =\widehat{\bigoplus}_{n_+,n_-\ge0}\mathcal H_{R;n_+,n_-},
 \label{cusp-native-full-flux-complete-decomposition}
\end{equation}
and the map from the vacuum to the joint level is the unitary
\begin{equation}
 \frac{(a_+^*)^{n_+}(a_-^*)^{n_-}}{\sqrt{n_+!n_-!}}.
 \label{cusp-native-full-flux-ladder-isomorphism}
\end{equation}
For a numerical eigenvalue, the multiplicity is six times the number of
joint pairs in \cref{cusp-native-full-flux-joint-energy} producing that
number.
\end{theorem}

\begin{proof}
The elliptic operator has compact resolvent, and the two commuting number
operators preserve its finite-dimensional eigenspaces; they can therefore be
diagonalized simultaneously.  If \(N_\sigma\psi=\lambda_\sigma\psi\), then
\(\lambda_\sigma\ge0\).  Repeated application of \(a_\sigma\) lowers the
Laplacian energy by \(2\mathcal B_{R,\sigma}>0\), so it terminates.  If
\(a_\sigma^r\psi\ne0\) and \(a_\sigma^{r+1}\psi=0\), then
\[
 0=\|a_\sigma^{r+1}\psi\|^2
  =(\lambda_\sigma-r)\|a_\sigma^r\psi\|^2,
\]
and hence \(\lambda_\sigma=r\in\mathbb Z_{\ge0}\).  Lowering both indices
reaches the simultaneous vacuum.

The commutator identities give, for a vacuum vector \(v\),
\[
 \|(a_\sigma^*)^nv\|^2=n!\|v\|^2,
 \qquad a_\sigma^n(a_\sigma^*)^nv=n!v.
\]
On the \(n\)-th level they also give
\((a_\sigma^*)^na_\sigma^n=n!\).  Operators in the two planes commute, so
\cref{cusp-native-full-flux-ladder-isomorphism} is unitary with inverse
\(a_+^{n_+}a_-^{n_-}/\sqrt{n_+!n_-!}\).  This proves the dimension, energy,
and completeness assertions.  Orthogonal joint eigenspaces with coincident
numerical energies must be added, which proves the final sentence.
\end{proof}

The compact-resolvent lowering argument proves that no additional joint
eigenspaces occur.  Independently, this is exactly the case \(m=2\), \(k=1\)
of
\href{https://www.math.hu-berlin.de/~schueth/quantumtori.pdf}%
{Gordon--Kirwin--Schueth--Webb, Proposition~4.5, article pp.~9--10
(volume pp.~175--176)},
\href{https://doi.org/10.1090/pspum/084/1355}%
{doi:10.1090/pspum/084/1355}.  Their signed tensor-power statement gives
every oscillator tuple multiplicity
\(|k|^mV_{E_\eta}=|k|^2\operatorname{Pf}(E_\eta)=6k^2\), so it also supplies
the external completeness check for every \(k\ne0\).

The exact ground energy and first distinct spacing are therefore
\begin{align}
 E_0(R)&=\mathcal B_{R,+}+\mathcal B_{R,-}
 =\frac{2\pi\Xi_R}{D_R}
 =\frac{14\pi}{R}+\frac{\pi(12B-14H)}{R^2}+O(R^{-3}),
 \label{cusp-native-full-flux-ground-energy}\\
 g_0(R)&=2\min\{\mathcal B_{R,+},\mathcal B_{R,-}\}
 =\frac{2\pi(\Xi_R-|S_R|)}{D_R}
 =\frac{48\pi}{\Xi_R+|S_R|}.
 \label{cusp-native-full-flux-global-gap}
\end{align}
On the deep tail \cref{cusp-native-full-flux-deep-tail-window}, this is
\(g_0(R)=2\mathcal B_{R,-}\), with expansion
\begin{equation}
 g_0(R)=\frac{4\pi}{R}-\frac{4\pi H}{R^2}
 +\frac{4\pi}{R^3}\left(H^2-\frac67m^2\right)+O(R^{-4}).
 \label{cusp-native-full-flux-gap-expansion}
\end{equation}

\begin{corollary}[Heat trace and exact fixed-energy count]
\label{cusp-native-full-flux-heat-count}
For \(t>0\),
\begin{equation}
 \operatorname{Tr}(e^{-t\Delta_{\eta,R}})
 =\frac{3}{2\sinh(t\mathcal B_{R,+})
                 \sinh(t\mathcal B_{R,-})}
 =\frac{D_R}{16\pi^2t^2}\bigl(1+O_t(R^{-2})\bigr).
 \label{cusp-native-full-flux-heat-trace}
\end{equation}
For \(E>0\), let \(N_{\eta,R}(E)\) count complex eigenvalues not exceeding
\(E\), with multiplicity.  Then
\begin{equation}
 N_{\eta,R}(E)=
 \begin{cases}
 0,&E<E_0(R),\\[1mm]
 6\displaystyle\sum_{n_+=0}^{\left\lfloor
 (E-E_0(R))/(2\mathcal B_{R,+})\right\rfloor}
 \left(1+\left\lfloor
 \dfrac{E-E_0(R)-2\mathcal B_{R,+}n_+}
       {2\mathcal B_{R,-}}
 \right\rfloor\right),&E\ge E_0(R).
 \end{cases}
 \label{cusp-native-full-flux-counting-function}
\end{equation}
For fixed \(E>0\), as \(R\to\infty\),
\begin{equation}
 N_{\eta,R}(E)
 =\frac{D_R}{32\pi^2}(E-E_0(R))^2+O_E(R)
 =\frac{E^2D_R}{32\pi^2}+O_E(R).
 \label{cusp-native-full-flux-counting-asymptotic}
\end{equation}
\end{corollary}

\begin{proof}
The heat trace is the product of the two geometric series in
\cref{cusp-native-full-flux-joint-energy}, multiplied by six.  Expanding both
hyperbolic sines and using
\cref{cusp-native-full-flux-frequency-product} proves its asymptotic.
The counting formula enumerates the lattice points
\[
 (n_+,n_-)\in\mathbb R_{\ge0}^2,
 \qquad
 2\mathcal B_{R,+}n_++2\mathcal B_{R,-}n_-
 \le E-E_0(R).
\]
Its triangular area is
\((E-E_0(R))^2/(8\mathcal B_{R,+}\mathcal B_{R,-})\).  The difference
between the lattice count and the area is bounded by a constant times the
perimeter plus one, hence is \(O_E(R)\).  Multiplying by six and using the
exact frequency product gives the first asymptotic; the second follows from
\(E_0(R)=O(R^{-1})\) and \(D_R=O(R^2)\).
\end{proof}

\subsection{The split \texorpdfstring{\(SU(2)\)}{SU(2)} critical connection}

Put
\begin{equation}
 \mathscr E_{\eta,R}=\mathscr L_{\eta,R}\oplus
 \mathscr L_{\eta,R}^{-1},
 \qquad
 A_{\eta,R}=\nabla_{\eta,R}\oplus\nabla_{\eta,R}^{-1}.
 \label{cusp-native-full-flux-split-bundle-connection}
\end{equation}
The determinant is canonically trivial and
\begin{equation}
 c(\mathscr E_{\eta,R})=(1+\eta_R)(1-\eta_R)=1-\eta_R^2,
 \qquad c_2(\mathscr E_{\eta,R})=-\eta_R^2.
 \label{cusp-native-full-flux-Chern-class}
\end{equation}
Thus the orientation-sensitive Chern numbers are
\begin{equation}
 \langle c_2(\mathscr E_{\eta,R}),[F_R]_{\mathrm{mark}}\rangle=-12,
 \qquad
 \langle c_2(\mathscr E_{\eta,R}),
 [F_R]_{x\text{-}\mathrm{complex}}\rangle=+12.
 \label{cusp-native-full-flux-Chern-numbers}
\end{equation}
The cohomology class \(c_2=-\eta_R^2\) is nonzero, so this split
\(SU(2)\)-bundle is not smoothly trivial.  This is another exact distinction
from the unit-flux split bundle of \cref{cusp-magnetic-su2-triviality}.

With
\begin{equation}
 H_{\mathrm{mag}}=\begin{pmatrix}i&0\\0&-i\end{pmatrix},
 \qquad
 \langle X,Y\rangle_{\mathfrak{su}(2)}
 =-\frac12\operatorname{Tr}(XY),
 \label{cusp-native-full-flux-Lie-normalization}
\end{equation}
one has \(\lvert H_{\mathrm{mag}}\rvert=1\) and
\begin{equation}
 F_{A_{\eta,R}}=-2\pi\eta_R H_{\mathrm{mag}}.
 \label{cusp-native-full-flux-SU2-curvature}
\end{equation}
This curvature is parallel; hence
\(d_{A_{\eta,R}}^*F_{A_{\eta,R}}=0\), and \(A_{\eta,R}\) is Yang--Mills.
For
\begin{equation}
 \operatorname{YM}_R(A)=\frac1{2g_{\mathrm{YM}}^2}
 \int_{F_R}|F_A|^2\,\mathrm d\operatorname{vol}_{g_R},
 \label{cusp-native-full-flux-action-normalization}
\end{equation}
diagonalizing \cref{cusp-native-full-flux-K} gives the exact action
\begin{align}
 \operatorname{YM}_R(A_{\eta,R})
 &=\frac{D_R}{2g_{\mathrm{YM}}^2}
   (\mathcal B_{R,+}^2+\mathcal B_{R,-}^2)
 \label{cusp-native-full-flux-action-frequency}\\
 &=\frac{2\pi^2}{g_{\mathrm{YM}}^2D_R}
   (q_R^2+36L_R^2+72m_R^2)
 \label{cusp-native-full-flux-action-coordinate}\\
 &=\frac{24\pi^2}{g_{\mathrm{YM}}^2}
 +\frac{2\pi^2S_R^2}{g_{\mathrm{YM}}^2D_R}.
 \label{cusp-native-full-flux-action-excess}
\end{align}
Indeed
\(q_R^2+36L_R^2+72m_R^2=12D_R+(6L_R-q_R)^2\).
Consequently
\begin{equation}
 \operatorname{YM}_R(A_{\eta,R})
 =\frac{74\pi^2}{g_{\mathrm{YM}}^2}
 +\frac{70\pi^2B}{g_{\mathrm{YM}}^2R}+O(R^{-2}).
 \label{cusp-native-full-flux-action-asymptotic}
\end{equation}

\subsection{The closed Jacobi operator and exact Coulomb reduction}

We next define the operator before computing its spectrum.  Let
\begin{equation}
 \mathscr H_R^1=L^2\Omega^1(F_R;\operatorname{ad}\mathscr E_{\eta,R}),
 \qquad
 \mathscr H_R^0=L^2\Omega^0(F_R;\operatorname{ad}\mathscr E_{\eta,R})
 \label{cusp-native-full-flux-real-Hilbert-spaces}
\end{equation}
be real Hilbert spaces.  Their complexifications carry the root
decomposition used below, and every operator constructed here commutes with
the anti-linear real structure.

Define the maximal distributional codifferential by
\begin{align}
 \operatorname{Dom}D_R^{\max}
 =\{a\in\mathscr H_R^1:\;&\text{there is }u\in\mathscr H_R^0\text{ such that}
 \notag\\[-1mm]
 &\langle a,d_{A_{\eta,R}}\phi\rangle_{L^2}
 =\langle u,\phi\rangle_{L^2}
 \text{ for every smooth }\phi\},
 \label{cusp-native-full-flux-max-codiff-domain}\\
 D_R^{\max}a={}&u,
 \qquad \mathscr S_R:=\ker D_R^{\max}.
 \label{cusp-native-full-flux-Coulomb-slice}
\end{align}
Uniqueness follows from density of the test space.  Passing to the limit in
the defining identity proves that \(D_R^{\max}\) is closed, hence the Coulomb
slice \(\mathscr S_R\) is a closed subspace.

On adjoint-valued scalars, let \(\Delta_{A_{\eta,R}}^{(0)}\) be the
nonnegative self-adjoint operator of the closed form
\begin{equation}
 \mathfrak q_R^{(0)}[\phi,\psi]
 =\langle d_{A_{\eta,R}}\phi,d_{A_{\eta,R}}\psi\rangle_{L^2},
 \qquad
 \operatorname{Dom}\mathfrak q_R^{(0)}=H^1\Omega^0.
 \label{cusp-native-full-flux-scalar-form}
\end{equation}
Elliptic regularity gives
\begin{equation}
 \Delta_{A_{\eta,R}}^{(0)}=d_{A_{\eta,R}}^*d_{A_{\eta,R}},
 \qquad
 \operatorname{Dom}\Delta_{A_{\eta,R}}^{(0)}=H^2\Omega^0.
 \label{cusp-native-full-flux-scalar-domain}
\end{equation}

In a constant orthonormal frame, define the symmetric zeroth-order
endomorphism
\begin{equation}
 (V_Ra)_\nu=-2\sum_{\mu=1}^4
 [F_{A_{\eta,R},\nu\mu},a_\mu].
 \label{cusp-native-full-flux-spin-endomorphism}
\end{equation}
Symmetry follows from invariance of the compact Lie-algebra inner product,
antisymmetry of \(F_{\nu\mu}\), and an exchange of \(\nu,\mu\).  Consider
the gauge-fixed form
\begin{equation}
 \mathfrak j_R^{\mathrm{gf}}[a,b]
 =\langle\nabla_{A_{\eta,R}}a,\nabla_{A_{\eta,R}}b\rangle_{L^2}
  +\langle V_Ra,b\rangle_{L^2},
 \qquad
 \operatorname{Dom}\mathfrak j_R^{\mathrm{gf}}=H^1\Omega^1.
 \label{cusp-native-full-flux-Jacobi-form}
\end{equation}

\begin{proposition}[Closed gauge-fixed realization]
\label{cusp-native-full-flux-closed-Jacobi}
The form \cref{cusp-native-full-flux-Jacobi-form} is densely defined, closed,
symmetric, and bounded below.  Its self-adjoint operator
\(\mathcal J_R^{\mathrm{gf}}\) has compact resolvent, domain \(H^2\Omega^1\),
and expression
\begin{equation}
 (\mathcal J_R^{\mathrm{gf}}a)_\nu
 =-\sum_{\mu=1}^4
   \nabla_{A_{\eta,R},e_\mu}\nabla_{A_{\eta,R},e_\mu}a_\nu
  -2\sum_{\mu=1}^4[F_{A_{\eta,R},\nu\mu},a_\mu].
 \label{cusp-native-full-flux-Jacobi-expression}
\end{equation}
It is exactly the Hessian operator with the background-gauge term, before
the overall factor \(g_{\mathrm{YM}}^{-2}\).
\end{proposition}

\begin{proof}
The first summand in \cref{cusp-native-full-flux-Jacobi-form}, together with
the \(L^2\)-norm, is an equivalent \(H^1\)-norm.  The smooth endomorphism
\(V_R\) is bounded and symmetric.  The bounded perturbation theorem for
closed forms gives closedness and a lower bound, and the representation
theorem gives a self-adjoint operator.  Its principal symbol is
\(|\xi|^2\) times the identity.  Elliptic regularity identifies its maximal
\(L^2\)-domain with \(H^2\Omega^1\); direct substitution supplies the reverse
inclusion.  The compact embedding \(H^1\hookrightarrow L^2\) gives compact
resolvent.

For the final assertion, define
\([a\wedge b]_{\mu\nu}=[a_\mu,b_\nu]-[a_\nu,b_\mu]\).  The exact
two-parameter curvature expansion is
\[
 F_{A_{\eta,R}+sa+tb}=F_{A_{\eta,R}}+s d_Aa+t d_Ab
 +\frac{s^2}{2}[a\wedge a]+st[a\wedge b]
 +\frac{t^2}{2}[b\wedge b].
\]
It yields
\begin{equation}
 \operatorname{YM}_R''\big|_{A_{\eta,R}}(a,b)
 =\frac1{g_{\mathrm{YM}}^2}
 \left(\langle d_Aa,d_Ab\rangle
 +\langle F_A,[a\wedge b]\rangle\right).
 \label{cusp-native-full-flux-raw-Hessian}
\end{equation}
On the flat torus, commuting the covariant derivatives with the retained
curvature convention gives
\begin{align*}
 \langle d_Aa,d_Ab\rangle+\langle d_A^*a,d_A^*b\rangle
 &=\langle\nabla_Aa,\nabla_Ab\rangle
 -\sum_{\nu,\mu}\langle a_\nu,[F_{\nu\mu},b_\mu]\rangle,\\
 \langle F_A,[a\wedge b]\rangle
 &=-\sum_{\nu,\mu}\langle a_\nu,[F_{\nu\mu},b_\mu]\rangle.
\end{align*}
The factor \(1/2\) in the two-form inner product is included in the second
identity.  Adding
\(g_{\mathrm{YM}}^{-2}\langle d_A^*a,d_A^*b\rangle\) to
\cref{cusp-native-full-flux-raw-Hessian} gives exactly
\(g_{\mathrm{YM}}^{-2}\mathfrak j_R^{\mathrm{gf}}[a,b]\), including the
coefficient \(-2\) in \cref{cusp-native-full-flux-spin-endomorphism}.
\end{proof}

\begin{proposition}[Resolvent intertwining and the reducing Coulomb slice]
\label{cusp-native-full-flux-Coulomb-reduction}
On smooth sections,
\begin{equation}
 \mathcal J_R^{\mathrm{gf}}d_{A_{\eta,R}}
 =d_{A_{\eta,R}}\Delta_{A_{\eta,R}}^{(0)},
 \qquad
 d_{A_{\eta,R}}^*\mathcal J_R^{\mathrm{gf}}
 =\Delta_{A_{\eta,R}}^{(0)}d_{A_{\eta,R}}^*.
 \label{cusp-native-full-flux-formal-intertwining}
\end{equation}
For \(z\in\mathbb C\setminus\mathbb R\) and
\(f\in\operatorname{Dom}D_R^{\max}\),
\begin{equation}
 D_R^{\max}(\mathcal J_R^{\mathrm{gf}}-z)^{-1}f
 =(\Delta_{A_{\eta,R}}^{(0)}-z)^{-1}D_R^{\max}f.
 \label{cusp-native-full-flux-resolvent-intertwining}
\end{equation}
Thus \(\mathscr S_R\) reduces \(\mathcal J_R^{\mathrm{gf}}\), and its part
\begin{equation}
 \mathcal J_R=\mathcal J_R^{\mathrm{gf}}\big|_{\mathscr S_R},
 \qquad
 \operatorname{Dom}\mathcal J_R=H^2\Omega^1\cap\mathscr S_R
 \label{cusp-native-full-flux-Coulomb-operator}
\end{equation}
is self-adjoint.  Its closed form is the restriction of
\(\mathfrak j_R^{\mathrm{gf}}\) to \(H^1\Omega^1\cap\mathscr S_R\), and the
Yang--Mills Hessian on this slice is
\begin{equation}
 \operatorname{YM}_R''\big|_{A_{\eta,R}}
 =g_{\mathrm{YM}}^{-2}\mathfrak j_R.
 \label{cusp-native-full-flux-Coulomb-Hessian}
\end{equation}
\end{proposition}

\begin{proof}
For a smooth adjoint-valued scalar \(\phi\), write
\(\nabla_\mu=\nabla_{A_{\eta,R},e_\mu}\).  Retaining both curvature terms
when commuting three derivatives gives
\begin{align*}
 &(\mathcal J_R^{\mathrm{gf}}d_A\phi
       -d_A\Delta_A^{(0)}\phi)_\nu\\
 &\quad=-\nabla_\mu\nabla_\mu\nabla_\nu\phi
       +\nabla_\nu\nabla_\mu\nabla_\mu\phi
       -2[F_{\nu\mu},\nabla_\mu\phi]\\
 &\quad=-[\nabla_\mu F_{\mu\nu},\phi]
       =[(d_A^*F_A)_\nu,\phi]=0.
\end{align*}
This proves the first identity in
\cref{cusp-native-full-flux-formal-intertwining}; formal adjunction gives the
second.

Let \(u=(\mathcal J_R^{\mathrm{gf}}-z)^{-1}f\).  The operator domain gives
\(u\in H^2\), hence \(D_R^{\max}u\in H^1\).  Distributionally applying the
second intertwining identity to
\((\mathcal J_R^{\mathrm{gf}}-z)u=f\) yields
\[
 (\Delta_A^{(0)}-z)D_R^{\max}u=D_R^{\max}f.
\]
The right side is in \(L^2\); scalar elliptic regularity upgrades
\(D_R^{\max}u\) to \(H^2\), and the scalar resolvent proves
\cref{cusp-native-full-flux-resolvent-intertwining}.  Both resolvents of the
self-adjoint gauge-fixed operator preserve \(\mathscr S_R\); adjunction shows
that they also preserve its orthogonal complement.  Therefore the slice is
reducing and its operator part is self-adjoint with the displayed domain.
The reducing projection commutes with the spectral calculus, so the form
domain is exactly \(H^1\Omega^1\cap\mathscr S_R\).  The Hessian identity is
the restriction of the last assertion of
\cref{cusp-native-full-flux-closed-Jacobi}.
\end{proof}

\subsection{Charge two, all polarizations, and the Morse index}

Choose root vectors with
\([H_{\mathrm{mag}},E_\pm]=\pm2iE_\pm\).  Then
\begin{equation}
 (\operatorname{ad}\mathscr E_{\eta,R})_{\mathbb C}
 =\underline{\mathbb C}H_{\mathrm{mag}}
 \oplus\mathscr L_{\eta,R}^{\otimes2}E_+
 \oplus\mathscr L_{\eta,R}^{\otimes(-2)}E_-.
 \label{cusp-native-full-flux-root-decomposition}
\end{equation}
The \(E_+\) root connection has curvature \(-4\pi i\eta_R\).  Replacing
\(E_\eta\) by \(2E_\eta\) gives \(J_{2,R}=J_{\eta,R}\), doubles both
frequencies, and changes the vacuum dimension from six to twenty-four by
\cref{cusp-native-full-flux-vacuum-dimensions}.  The \(E_-\) root has
\(q=-2\), uses the conjugate auxiliary structure
\(J_{-2,R}=-J_{\eta,R}\), and has the same positive frequencies and the same
dimension \(6(-2)^2=24\).  Repeating the complete ladder proof of
\cref{cusp-native-full-flux-landau-theorem} therefore gives, for either root,
the joint charge-two scalar levels
\begin{equation}
 E^{(2)}_{R;p,q}
 =2(2p+1)\mathcal B_{R,+}+2(2q+1)\mathcal B_{R,-},
 \qquad p,q\in\mathbb Z_{\ge0},
 \label{cusp-native-full-flux-charge-two-levels}
\end{equation}
each of complex dimension
\begin{equation}
 |\operatorname{Pf}(2E_\eta)|=24.
 \label{cusp-native-full-flux-charge-two-multiplicity}
\end{equation}
Coincident joint pairs must again be summed.

Choose \(O_R\in SO(2)\) with
\(O_R^{\mathsf T}K_RO_R=\operatorname{diag}(\kappa_{R,+},\kappa_{R,-})\),
and define the constant orthonormal coframe
\begin{equation}
 \binom{\mathrm dX_+}{\mathrm dX_-}
 =O_R^{\mathsf T}\binom{\mathrm dx^1}{\mathrm dx^3},
 \qquad
 \binom{\mathrm dY_+}{\mathrm dY_-}
 =O_R^{\mathsf T}\binom{\mathrm dx^2}{\mathrm dx^4}.
 \label{cusp-native-full-flux-orthogonal-coframe}
\end{equation}
Then
\begin{equation}
 \eta_R=\frac{\mathcal B_{R,+}}{2\pi}
          \mathrm dX_+\wedge\mathrm dY_+
         -\frac{\mathcal B_{R,-}}{2\pi}
          \mathrm dX_-\wedge\mathrm dY_-.
 \label{cusp-native-full-flux-normal-form}
\end{equation}
The odd grouping permutation in
\cref{cusp-native-full-flux-orthogonal-coframe} reverses orientation; it is
used only to expose the two blocks and never to infer a Chern sign.

On the \(E_+\) root line, the four constant circular one-forms diagonalize
the spin endomorphism.  Direct substitution into
\cref{cusp-native-full-flux-spin-endomorphism}, with the charge-two root
normalization retained, gives
\begin{center}
\begin{tabular}{>{$}c<{$} >{$}c<{$} >{$}c<{$}}
\toprule
\text{polarization}&V_R\text{-shift}
 &\mathcal J_R^{\mathrm{gf}}\text{-eigenvalue on }(p,q)\\
\midrule
\mathrm dX_+ + i\mathrm dY_+&-4\mathcal B_{R,+}
 &2[(2p-1)\mathcal B_{R,+}+(2q+1)\mathcal B_{R,-}]\\
\mathrm dX_+ - i\mathrm dY_+&+4\mathcal B_{R,+}
 &2[(2p+3)\mathcal B_{R,+}+(2q+1)\mathcal B_{R,-}]\\
\mathrm dX_- - i\mathrm dY_-&-4\mathcal B_{R,-}
 &2[(2p+1)\mathcal B_{R,+}+(2q-1)\mathcal B_{R,-}]\\
\mathrm dX_- + i\mathrm dY_-&+4\mathcal B_{R,-}
 &2[(2p+1)\mathcal B_{R,+}+(2q+3)\mathcal B_{R,-}]\\
\bottomrule
\end{tabular}
\end{center}
The \(E_-\) table is its anti-linear conjugate.  The Cartan sector is the
ordinary nonnegative coefficientwise Laplacian on one-forms.

\begin{theorem}[Exact negative spectrum, walls, and real Morse index]
\label{cusp-native-full-flux-Morse-theorem}
Put
\begin{equation}
 \mathcal B_{R,\max}=\max\{\mathcal B_{R,+},\mathcal B_{R,-}\},
 \quad
 \mathcal B_{R,\min}=\min\{\mathcal B_{R,+},\mathcal B_{R,-}\},
 \quad
 \rho_R^*=\frac{\mathcal B_{R,\max}}{\mathcal B_{R,\min}}
 =\frac{\Xi_R+|S_R|}{\Xi_R-|S_R|}.
 \label{cusp-native-full-flux-ordered-ratio}
\end{equation}
For every regular \(R\), the negative spectrum of the self-adjoint Coulomb
operator is exactly
\begin{equation}
 \left\{
 \nu_{R,n}=2[(2n+1)\mathcal B_{R,\min}-\mathcal B_{R,\max}]
 =\frac{4\pi}{D_R}[n\Xi_R-(n+1)|S_R|]:
 2n+1<\rho_R^*
 \right\}.
 \label{cusp-native-full-flux-general-negative-spectrum}
\end{equation}
Every indicated band has real multiplicity \(48\), and
\begin{equation}
 \operatorname{ind}_{\mathbb R}(\mathcal J_R)
 =48\left\lceil\frac{\rho_R^*-1}{2}\right\rceil.
 \label{cusp-native-full-flux-general-index}
\end{equation}
At the \(n\)-th wall the strict inequality becomes equality.  Equivalently,
\begin{equation}
 \rho_R^*=2n+1
 \iff n\Xi_R=(n+1)|S_R|
 \iff 24n^2D_R=(2n+1)S_R^2.
 \label{cusp-native-full-flux-wall-equations}
\end{equation}
The corresponding branch is then a zero branch and is not counted in the
Morse index.  If one removes the absolute value by choosing the
\(\mathcal B_{R,+}>\mathcal B_{R,-}\) side, the polynomial equation in
\cref{cusp-native-full-flux-wall-equations} must be accompanied by \(S_R>0\);
otherwise squaring introduces the opposite side as an extraneous solution.

On the sufficiently deep tail of
\cref{cusp-native-full-flux-deep-tail-window}, the negative branches are
exactly \(n=0,1,2\),
\begin{equation}
 \nu_{R,n}=\frac{4\pi}{D_R}[n\Xi_R-(n+1)S_R],
 \qquad n=0,1,2,
 \label{cusp-native-full-flux-deep-negative-branches}
\end{equation}
each of real multiplicity \(48\).  Therefore
\begin{equation}
 \operatorname{ind}_{\mathbb R}(\mathcal J_R)=144,
 \label{cusp-native-full-flux-deep-index}
\end{equation}
and
\begin{align}
 \nu_{R,0}&=-\frac{20\pi}{R}+O(R^{-2}),&
 \nu_{R,1}&=-\frac{12\pi}{R}+O(R^{-2}),\notag\\
 \nu_{R,2}&=-\frac{4\pi}{R}+O(R^{-2}),&
 \nu_{R,3}&=+\frac{4\pi}{R}+O(R^{-2}).
 \label{cusp-native-full-flux-negative-asymptotics}
\end{align}
The Hessian operator for
\cref{cusp-native-full-flux-action-normalization} has eigenvalues
\(\nu_{R,n}/g_{\mathrm{YM}}^2\) and the same index.
\end{theorem}

\begin{proof}
Assume first \(\mathcal B_{R,+}>\mathcal B_{R,-}\).  In the first row of
the polarization table, \(p\ge1\) makes the eigenvalue positive.  At \(p=0\)
it is
\[
 2[-\mathcal B_{R,+}+(2q+1)\mathcal B_{R,-}],
\]
which is negative exactly when \(2q+1<\rho_R^*\).  The second and fourth
rows are positive, and the third has minimum
\(2(\mathcal B_{R,+}-\mathcal B_{R,-})>0\).  When
\(\mathcal B_{R,-}>\mathcal B_{R,+}\), the identical argument with the two
planes exchanged shows that only the \(q=0\) part of the third row can be
negative and gives the same formula with the ordered frequencies.  If the
frequencies agree, both candidate lowest branches are zero and none is
negative.  This proves the exhaustive list for the gauge-fixed operator.

Every negative gauge-fixed eigenform is automatically Coulomb.  Elliptic
regularity makes it smooth; if
\(\mathcal J_R^{\mathrm{gf}}a=\lambda a\) with \(\lambda<0\), then
\cref{cusp-native-full-flux-formal-intertwining} gives
\[
 \Delta_{A_{\eta,R}}^{(0)}(d_{A_{\eta,R}}^*a)
 =\lambda d_{A_{\eta,R}}^*a.
\]
Nonnegativity of the scalar operator forces \(d_A^*a=0\).  Hence the full and
Coulomb negative spectral spaces coincide.

One joint charge-two scalar space contributes complex dimension \(24\) in
the \(E_+\) sector.  Its conjugate contributes complex dimension \(24\) in
the complexification of the real adjoint problem.  A complexified real
eigenspace has complex dimension equal to the original real dimension, so
the real multiplicity is \(24+24=48\), not \(24\) or \(96\).

The count of integers \(n\ge0\) satisfying \(2n+1<\rho_R^*\) is
\(\lceil(\rho_R^*-1)/2\rceil\), including the correct exclusion at an odd
integer wall.  Substitution of the ordered frequency formulas gives
\cref{cusp-native-full-flux-general-negative-spectrum}.  The wall equation
is
\(n\Xi_R=(n+1)|S_R|\); squaring and using
\(\Xi_R^2-S_R^2=24D_R\) gives the final polynomial form, and nonnegativity of
both sides gives the converse.  Finally,
\cref{cusp-native-full-flux-deep-tail-window} leaves precisely
\(n=0,1,2\), and the leading frequency terms give
\(\nu_{R,n}=(8n-20)\pi/R+O(R^{-2})\).
\end{proof}

The lowest branch and the action excess are tied by the exact identities
\begin{equation}
 \nu_{R,0}=-2|\mathcal B_{R,+}-\mathcal B_{R,-}|
 =-\frac{4\pi|S_R|}{D_R},
 \qquad
 \operatorname{YM}_R(A_{\eta,R})-\frac{24\pi^2}{g_{\mathrm{YM}}^2}
 =\frac{D_R}{8g_{\mathrm{YM}}^2}\nu_{R,0}^2.
 \label{cusp-native-full-flux-action-instability}
\end{equation}



% Standalone-preprint adaptation of cusp_native_full_flux_radial.tex.
% Canonical source SHA-256: 362E218ADEF6D2E6F7A549B14DE21A4FB2DF9DF42037391A0554B841CF139516.
% The canonical radial-height symbol T was renamed R consistently; transpose
% notation was retained exactly.

% Canonical integration of agent_logs/native_direct_integral_complete_canonical_20260903.tex.
% The fixed and varying fields remain separately typed; R is a multiplication
% parameter and never a hidden radial derivative.

\subsection{The fixed and varying radial direct integrals}
\label{cusp-native-full-flux-radial-section}

Choose \(R_1>R_0\) so that the complete deep-tail inequalities in
\cref{cusp-native-full-flux-deep-tail-window} hold on
\[
 I:=(R_1,\infty).
\]
Let \(Y=\mathbb R^4/\mathbb Z^4\).  Denote by
\(\mathscr L_\eta\to Y\) the fixed marked Hermitian line bundle defined
by the factor of automorphy in
\cref{cusp-native-full-flux-cocycle}.  The marking diffeomorphism and its
transported bundle isomorphism are
\[
 \Phi_R:Y\longrightarrow F_R,\qquad [y]\longmapsto[P_Ry],
 \qquad
 \widehat\Phi_R:\mathscr L_\eta\longrightarrow\mathscr L_{\eta,R}.
\]
They are fixed by requiring
\[
 \widehat\Phi_R^*\nabla_{\eta,R}
 =\mathrm d-2\pi i(6y^1\,\mathrm dy^4+y^2\,\mathrm dy^3),
\]
so no gauge, metric, lattice, or orientation normalization is hidden in the
identification.  Put
\[
 \mathscr K_R
 :=L^2(F_R,\mathscr L_{\eta,R};
       \mathrm d\operatorname{vol}_{g_R}).
\]

The pullback of the physical volume is
\begin{equation}
 \Phi_R^*(\mathrm d\operatorname{vol}_{g_R})
 =\sqrt{\det(P_R^{\mathsf T}P_R)}\,\mathrm d^4y
 =D_R\,\mathrm d^4y.
\label{native-di-volume-pullback}
\end{equation}
Put
\[
 \mathscr H_*:=L^2(Y,\mathscr L_\eta;\mathrm d^4y),
 \qquad
 \mathscr D_*:=H^2(Y;\mathscr L_\eta),
\]
where the Sobolev space is defined using any one fixed smooth reference
metric and connection on the fixed compact bundle.  Define
\begin{equation}
 U_R:\mathscr K_R\longrightarrow\mathscr H_*,
 \qquad
 U_Rs:=D_R^{1/2}\widehat\Phi_R^*s.
\label{native-di-marking-unitary}
\end{equation}
Then
\[
 \|U_Rs\|_{\mathscr H_*}^2
 =D_R\int_Y|\widehat\Phi_R^*s|^2\,\mathrm d^4y
 =\|s\|_{\mathscr K_R}^2,
\]
and
\[
 U_R^{-1}v
 =D_R^{-1/2}(\widehat\Phi_R^{-1})^*v.
\]
Thus \(U_R\) is an onto unitary, not merely an isometric injection.

Transport the fibre operator to the fixed space:
\begin{equation}
 A_R:=U_R\Delta_{\eta,R}U_R^{-1},
 \qquad
 \operatorname{Dom}(A_R)=\mathscr D_*.
\label{native-di-fixed-fibre}
\end{equation}
Its closed quadratic form is
\begin{equation}
 \widetilde{\mathfrak a}_R[v,w]
 =\int_Y(G_R^{-1})^{jk}
  \langle\nabla_{\eta,j}v,\nabla_{\eta,k}w\rangle\,\mathrm d^4y,
 \qquad G_R=P_R^{\mathsf T}P_R.
\label{native-di-fixed-form}
\end{equation}
Here \(\nabla_{\eta,j}\) is the \(\mathrm dy^j\)-component of the fixed
marked connection displayed above.  The factor \(D_R^{1/2}\) in \(U_R\) is
constant in \(y\), and it cancels
the factor \(D_R\) in \eqref{native-di-volume-pullback}; hence no unrecorded
factor appears in \eqref{native-di-fixed-form}.

\begin{lemma}[Common domain and measurable resolvents]
\label{native-di-common-domain-lemma}
For every compact interval \(K\Subset I\), the family
\[
 R\longmapsto A_R\in\mathcal B(\mathscr D_*,\mathscr H_*)
\]
is norm-continuous, in fact real analytic.  For every
\(z\in\mathbb C\setminus\mathbb R\), the resolvents
\((A_R-z)^{-1}\) are locally norm-continuous as operators on
\(\mathscr H_*\), and hence form a measurable operator field.
\end{lemma}

\begin{proof}
On \(K\), the matrices \(G_R^{-1}\) are uniformly positive definite, and
their entries are real analytic in \(R\).  In fixed local bundle
trivializations, \eqref{native-di-fixed-form} therefore gives a uniformly
elliptic second-order expression whose coefficients depend real
analytically on \(R\).  Elliptic regularity identifies every operator
domain with the same set \(\mathscr D_*\), and the standard local
coordinate estimate for a differential operator of order two gives
norm-continuity in
\(\mathcal B(\mathscr D_*,\mathscr H_*)\).

Uniform ellipticity on \(K\) gives
\[
 \|v\|_{\mathscr D_*}
 \le C_K\bigl(\|A_Rv\|_{\mathscr H_*}
                  +\|v\|_{\mathscr H_*}\bigr),
 \qquad R\in K.
\]
For \(v=(A_R-z)^{-1}f\), self-adjointness gives
\(\|v\|\le|\operatorname{Im}z|^{-1}\|f\|\), while
\(A_Rv=f+zv\).  Consequently
\begin{equation}
 \sup_{R\in K}
 \|(A_R-z)^{-1}\|_{\mathscr H_*\to\mathscr D_*}<\infty.
\label{native-di-elliptic-resolvent-bound}
\end{equation}
The resolvent identity, with its exact sign, for \(R,S\in K\),
\[
 (A_R-z)^{-1}-(A_S-z)^{-1}
 =-(A_R-z)^{-1}(A_R-A_S)(A_S-z)^{-1},
\]
together with \eqref{native-di-elliptic-resolvent-bound}, proves local
norm-resolvent continuity.  A countable compact exhaustion of \(I\) gives
measurability on all of \(I\).
\end{proof}

We now define the varying measurable Hilbert field rather than presuppose
one.  A section \(R\mapsto s_R\in\mathscr K_R\) is declared measurable
precisely when
\[
 R\longmapsto U_Rs_R\in\mathscr H_*
\]
is strongly measurable.  This is the measurable field pulled back from the
constant separable field \(I\times\mathscr H_*\).  Its direct integral is
\[
 \mathscr H_{\mathrm{var}}
 :=\int_I^\oplus\mathscr K_R\,\mathrm d R.
\]
The fixed radial Hilbert space is
\[
 \mathscr H_{\mathrm{rad}}
 :=L^2(I,\mathrm d R;\mathscr H_*).
\]
By construction,
\begin{equation}
 \mathscr U:=\int_I^\oplus U_R\,\mathrm d R:
 \mathscr H_{\mathrm{var}}\longrightarrow\mathscr H_{\mathrm{rad}},
 \qquad
 (\mathscr Us)(R)=U_Rs_R,
\label{native-di-direct-integral-unitary}
\end{equation}
is an isometry.  It is onto: if
\(\Psi\in\mathscr H_{\mathrm{rad}}\), then
\(s_R=U_R^{-1}\Psi(R)\) is measurable by the very definition of the
varying field and has the same squared norm integral.

On the fixed direct integral, define the maximal multiplication operator
\begin{equation}
 (\widetilde{\mathscr M}_\eta\Psi)(R):=A_R\Psi(R)
\label{native-di-fixed-multiplication}
\end{equation}
with the exact domain
\begin{equation}
\begin{split}
 \operatorname{Dom}(\widetilde{\mathscr M}_\eta)
 :=\{\Psi\in\mathscr H_{\mathrm{rad}}:\;&
 \Psi(R)\in\mathscr D_*\text{ for almost every }R,\\
 &R\mapsto A_R\Psi(R)\text{ is strongly measurable},\\
 &\int_I\|A_R\Psi(R)\|_{\mathscr H_*}^2\,\mathrm d R<\infty\}.
\end{split}
\label{native-di-fixed-maximal-domain}
\end{equation}
On the varying direct integral, define
\begin{equation}
\begin{split}
 \operatorname{Dom}(\mathscr M_\eta)
 :=\{s\in\mathscr H_{\mathrm{var}}:\;&
 s_R\in\operatorname{Dom}(\Delta_{\eta,R})\text{ for almost every }R,\\
 &R\mapsto U_R\Delta_{\eta,R}s_R
       \text{ is strongly measurable},\\
 &\int_I\|\Delta_{\eta,R}s_R\|_{\mathscr K_R}^2\,\mathrm d R<\infty\},
\end{split}
\label{native-di-varying-maximal-domain}
\end{equation}
and
\begin{equation}
 (\mathscr M_\eta s)_R:=\Delta_{\eta,R}s_R
 \quad\text{for almost every }R.
\label{native-di-varying-multiplication}
\end{equation}
Every expression in \eqref{native-di-fixed-multiplication} and
\eqref{native-di-varying-multiplication} is now typed on its actual fibre.
Equations \eqref{native-di-fixed-fibre},
\eqref{native-di-direct-integral-unitary},
\eqref{native-di-fixed-maximal-domain}, and
\eqref{native-di-varying-maximal-domain} give both the domain equality and
the exact intertwining identity
\begin{equation}
 \mathscr U\operatorname{Dom}(\mathscr M_\eta)
 =\operatorname{Dom}(\widetilde{\mathscr M}_\eta),
 \qquad
 \mathscr U\mathscr M_\eta\mathscr U^{-1}
 =\widetilde{\mathscr M}_\eta.
\label{native-di-intertwining}
\end{equation}

\subsubsection{Density and self-adjointness on the maximal domain}

\begin{proposition}[Maximal-domain realization]
\label{native-di-self-adjoint-proposition}
The operators \(\widetilde{\mathscr M}_\eta\) and
\(\mathscr M_\eta\) are densely defined, nonnegative, and self-adjoint on
the domains \eqref{native-di-fixed-maximal-domain} and
\eqref{native-di-varying-maximal-domain}, respectively.
\end{proposition}

\begin{proof}
For \(\chi\in C_c(I)\) and \(v\in\mathscr D_*\), local continuity of
\(R\mapsto A_Rv\) makes it bounded on \(\operatorname{supp}\chi\).
Thus \(\chi v\in\operatorname{Dom}(\widetilde{\mathscr M}_\eta)\).  Since
\(\mathscr D_*\) is dense in \(\mathscr H_*\), the algebraic tensor product
\[
 C_c(I)\mathbin{\odot}\mathscr D_*
\]
is dense in \(L^2(I;\mathscr H_*)\).  Hence the fixed operator is densely
defined.  Fibrewise self-adjointness and Cauchy--Schwarz give, for
\(\Psi,\Theta\) in its domain,
\[
 \langle\widetilde{\mathscr M}_\eta\Psi,\Theta\rangle
 =\int_I\langle A_R\Psi(R),\Theta(R)\rangle\,\mathrm d R
 =\langle\Psi,\widetilde{\mathscr M}_\eta\Theta\rangle,
\]
so it is symmetric.

Fix \(z\in\mathbb C\setminus\mathbb R\).  For
\(F\in\mathscr H_{\mathrm{rad}}\), define pointwise
\begin{equation}
 (\mathcal R_zF)(R):=(A_R-z)^{-1}F(R).
\label{native-di-radial-resolvent}
\end{equation}
The locally norm-continuous resolvent field, approximation of \(F\) by
simple functions, and
\[
 \|(A_R-z)^{-1}\|_{\mathscr H_*\to\mathscr H_*}
 \le|\operatorname{Im}z|^{-1}
\]
show that \(\mathcal R_zF\) is strongly measurable and
\begin{equation}
 \|\mathcal R_zF\|_{\mathscr H_{\mathrm{rad}}}
 \le|\operatorname{Im}z|^{-1}
    \|F\|_{\mathscr H_{\mathrm{rad}}}.
\label{native-di-radial-resolvent-L2}
\end{equation}
Pointwise, \((A_R-z)^{-1}F(R)\in\mathscr D_*\).  The exact identity
\begin{equation}
 A_R(A_R-z)^{-1}=I+z(A_R-z)^{-1}
\label{native-di-maximal-domain-identity}
\end{equation}
gives the global uniform bound
\[
 \|A_R(A_R-z)^{-1}\|
 \le1+\frac{|z|}{|\operatorname{Im}z|}.
\]
The right-hand side of \eqref{native-di-maximal-domain-identity}, applied
to \(F(R)\), is strongly measurable and square-integrable.  Therefore
\(\mathcal R_zF\) belongs to the full maximal domain
\eqref{native-di-fixed-maximal-domain}; this is the domain step that
pointwise inversion alone does not prove.

Fibrewise identities now yield
\[
 (\widetilde{\mathscr M}_\eta-z)\mathcal R_z=I
 \quad\text{on }\mathscr H_{\mathrm{rad}},
 \qquad
 \mathcal R_z(\widetilde{\mathscr M}_\eta-z)=I
 \quad\text{on }\operatorname{Dom}(\widetilde{\mathscr M}_\eta).
\]
In particular, the ranges of
\(\widetilde{\mathscr M}_\eta-i\) and
\(\widetilde{\mathscr M}_\eta+i\) are both all of
\(\mathscr H_{\mathrm{rad}}\).  The range criterion for a densely defined
symmetric operator proves self-adjointness.  Moreover,
\[
 \langle\widetilde{\mathscr M}_\eta\Psi,\Psi\rangle
 =\int_I\langle A_R\Psi(R),\Psi(R)\rangle\,\mathrm d R\ge0,
\]
so it is nonnegative.  The unitary equivalence
\eqref{native-di-intertwining} transfers density, self-adjointness, and
nonnegativity to the varying-fibre operator.
\end{proof}

\subsubsection{Measurable eigenvectors, including at band crossings}

\begin{lemma}[Explicit measurable eigenvector selection]
\label{native-di-measurable-eigenvector-lemma}
For every fixed
\(\mathbf n=(n_+,n_-)\in\mathbb Z_{\ge0}^2\), there is a strongly
measurable map
\[
 I\longrightarrow\mathscr D_*,
 \qquad R\longmapsto\phi_{\mathbf n}(R),
\]
such that
\begin{equation}
 \|\phi_{\mathbf n}(R)\|_{\mathscr H_*}=1,
 \qquad
 A_R\phi_{\mathbf n}(R)
 =E_{\mathbf n}(R)\phi_{\mathbf n}(R)
\label{native-di-measurable-eigenvector}
\end{equation}
for every \(R\in I\).  The construction remains valid when
\(E_{\mathbf n}(R)\) crosses one or more other bands.
\end{lemma}

\begin{proof}
Set
\[
 C_{\mathbf n,R}:=A_R-E_{\mathbf n}(R),
 \qquad \operatorname{Dom}(C_{\mathbf n,R})=\mathscr D_*.
\]
The band \(E_{\mathbf n}\) is continuous.  The proof of
Lemma~\ref{native-di-common-domain-lemma}, applied also to the scalar
shift, shows that
\((C_{\mathbf n,R}-z)^{-1}\) is a measurable resolvent field for every
nonreal \(z\).

Let
\[
 \mathsf P_{\mathbf n}(R)
 :=\mathbf 1_{\{0\}}(C_{\mathbf n,R})
 =\mathbf 1_{\{E_{\mathbf n}(R)\}}(A_R).
\]
This singleton projection is not assumed continuous at a crossing.
Instead, the spectral theorem gives the pointwise strong limit
\begin{equation}
 \mathsf P_{\mathbf n}(R)
 =\operatorname*{s-lim}_{k\to\infty}
  (I+kC_{\mathbf n,R}^2)^{-1},
\label{native-di-kernel-projection}
\end{equation}
and the exact resolvent formula
\begin{equation}
 (I+kC_{\mathbf n,R}^2)^{-1}
 =\frac1{2i\sqrt{k}}
 \left[
 (C_{\mathbf n,R}-i/\sqrt{k})^{-1}
 -(C_{\mathbf n,R}+i/\sqrt{k})^{-1}
 \right]
\label{native-di-kernel-projection-resolvent}
\end{equation}
shows that every approximating field is strongly measurable.  Therefore
\(R\mapsto\mathsf P_{\mathbf n}(R)v\) is strongly measurable for every
\(v\in\mathscr H_*\).  Proposition
\ref{cusp-native-full-flux-landau-theorem} says that
\(\mathsf P_{\mathbf n}(R)\ne0\) for every \(R\), including crossings.

Choose a countable dense sequence \((v_k)_{k\ge1}\) in the separable
space \(\mathscr H_*\), and define the measurable partition
\[
 X_k=
 \{R:\|\mathsf P_{\mathbf n}(R)v_k\|>0\}
 \cap
 \bigcap_{\ell<k}
 \{R:\|\mathsf P_{\mathbf n}(R)v_\ell\|=0\}.
\]
The sets \(X_k\) are pairwise disjoint and cover \(I\): a nonzero bounded
projection cannot annihilate a dense set.  On \(X_k\), set
\[
 \phi_{\mathbf n}(R)
 =\frac{\mathsf P_{\mathbf n}(R)v_k}
 {\|\mathsf P_{\mathbf n}(R)v_k\|}.
\]
Countable measurable patching gives a strongly measurable unit vector.
The range of the singleton spectral projection lies in
\(\operatorname{Dom}(A_R)\) and satisfies
\((A_R-E_{\mathbf n}(R))\phi_{\mathbf n}(R)=0\), which proves
\eqref{native-di-measurable-eigenvector}.
\end{proof}

\subsubsection{Full spectrum and spectral type}

\begin{theorem}[Purely continuous full half-line spectrum]
\label{native-di-main-theorem}
For both the fixed-fibre and varying-fibre maximal radial multiplication
operators,
\begin{equation}
\begin{aligned}
 \sigma(\widetilde{\mathscr M}_\eta)
 &=\sigma_{\mathrm{ess}}(\widetilde{\mathscr M}_\eta)
 =\sigma_{\mathrm c}(\widetilde{\mathscr M}_\eta)
 =[0,\infty),\\
 \sigma_{\mathrm p}(\widetilde{\mathscr M}_\eta)
 &=\varnothing,
 \qquad
 \ker\widetilde{\mathscr M}_\eta=\{0\},
\end{aligned}
\label{native-di-fixed-spectrum-result}
\end{equation}
and
\begin{equation}
\begin{aligned}
 \sigma(\mathscr M_\eta)
 &=\sigma_{\mathrm{ess}}(\mathscr M_\eta)
 =\sigma_{\mathrm c}(\mathscr M_\eta)
 =[0,\infty),\\
 \sigma_{\mathrm p}(\mathscr M_\eta)
 &=\varnothing,
 \qquad
 \ker\mathscr M_\eta=\{0\}.
\end{aligned}
\label{native-di-varying-spectrum-result}
\end{equation}
Here \(\sigma_{\mathrm{ess}}\) is the Weyl essential spectrum.
\end{theorem}

\begin{proof}
It is enough to prove \eqref{native-di-fixed-spectrum-result}, because
\eqref{native-di-intertwining} is a unitary equivalence.
Proposition~\ref{native-di-self-adjoint-proposition} and nonnegativity give
\[
 \sigma(\widetilde{\mathscr M}_\eta)\subset[0,\infty).
\]

\paragraph{The zero-energy singular sequence.}
Let \(a_j=R_1+2j+1\).  Using the measurable normalized ground eigenvector
from Lemma~\ref{native-di-measurable-eigenvector-lemma}, define
\[
 \Psi_j(R)
 :=\mathbf 1_{[a_j,a_j+1]}(R)\phi_{(0,0)}(R).
\]
These vectors have norm one and have pairwise disjoint supports, hence are
orthonormal and converge weakly to zero.  They belong to the maximal
domain because
\[
 A_R\Psi_j(R)=
 \mathbf 1_{[a_j,a_j+1]}(R)E_0(R)\phi_{(0,0)}(R)
\]
is strongly measurable and square-integrable.  From
\eqref{cusp-native-full-flux-ground-energy},
\[
 \|\widetilde{\mathscr M}_\eta\Psi_j\|^2
 =\int_{a_j}^{a_j+1}E_0(R)^2\,\mathrm d R
 =\frac{196\pi^2}{a_j^2}+O(a_j^{-3})
 \longrightarrow0.
\]
Thus \((\Psi_j)\) is a singular Weyl sequence at zero, and
\[
 0\in\sigma_{\mathrm{ess}}(\widetilde{\mathscr M}_\eta).
\]

\paragraph{Positive energy: escaping crossings before disjoint supports.}
Fix \(\lambda>0\), and fix once and for all
\(R_\bullet=R_1+1\in I\).  For \(n\in\mathbb Z_{\ge0}\), retain the
specific canonical band
\begin{equation}
 e_n(R):=E_{(0,n)}(R)
 =\mathcal B_{R,+}+(2n+1)\mathcal B_{R,-}.
\label{native-di-crossing-band}
\end{equation}
Since \(\mathcal B_{R_\bullet,-}>0\),
\[
 e_n(R_\bullet)\longrightarrow\infty
 \quad(n\to\infty).
\]
For each fixed \(n\), equations
\eqref{cusp-native-full-flux-Bminus-expansion} and
\eqref{cusp-native-full-flux-Bplus-expansion} give
\[
 e_n(R)\longrightarrow0
 \quad(R\to\infty).
\]
Hence, for every sufficiently large \(n\), the intermediate value theorem
gives at least one crossing parameter \(r_n>R_\bullet\) satisfying
\begin{equation}
 e_n(r_n)=\lambda.
\label{native-di-crossing}
\end{equation}

We prove explicitly that every such choice of crossing parameters escapes
to infinity.  Fix \(M>R_\bullet\).  Positivity and continuity give
\[
 c_M:=\min_{R\in[R_\bullet,M]}\mathcal B_{R,-}>0.
\]
For every \(R\in[R_\bullet,M]\),
\[
 e_n(R)\ge(2n+1)c_M.
\]
Thus, once \((2n+1)c_M>\lambda\), equation
\eqref{native-di-crossing} has no solution in
\([R_\bullet,M]\), and the chosen \(r_n\) must satisfy \(r_n>M\).
Since \(M\) was arbitrary,
\begin{equation}
 r_n\longrightarrow\infty.
\label{native-di-crossings-escape}
\end{equation}

Only now, before choosing any support, pass to a subsequence
\(n_j\to\infty\) such that, with \(r_j:=r_{n_j}\),
\begin{equation}
 r_{j+1}>r_j+2.
\label{native-di-separated-crossings}
\end{equation}
Continuity of \(e_{n_j}\) at \(r_j\) permits the subsequent choice of
\(\varepsilon_j\in(0,\tfrac14)\) such that
\[
 J_j:=(r_j-\varepsilon_j,r_j+\varepsilon_j)\Subset I,
 \qquad
 \sup_{R\in J_j}|e_{n_j}(R)-\lambda|<\frac1j.
\]
The order of construction matters: the separation
\eqref{native-di-separated-crossings} makes the intervals \(J_j\)
pairwise disjoint.

Take the measurable normalized eigenvector
\(\phi_{(0,n_j)}(R)\) from
Lemma~\ref{native-di-measurable-eigenvector-lemma}; the singleton
projection used there remains valid if the selected band crosses another
one inside \(J_j\).  Define
\begin{equation}
 \Theta_j(R)
 :=|J_j|^{-1/2}\mathbf 1_{J_j}(R)\phi_{(0,n_j)}(R).
\label{native-di-positive-Weyl-sequence}
\end{equation}
The sections \(\Theta_j\) have norm one and are orthogonal by disjointness
of their supports.  They lie in the maximal domain because
\[
 A_R\Theta_j(R)
 =|J_j|^{-1/2}\mathbf 1_{J_j}(R)
   e_{n_j}(R)\phi_{(0,n_j)}(R)
\]
is strongly measurable and square-integrable.  Moreover,
\[
 \|(\widetilde{\mathscr M}_\eta-\lambda)\Theta_j\|^2
 =\frac1{|J_j|}\int_{J_j}
   |e_{n_j}(R)-\lambda|^2\,\mathrm d R
 <\frac1{j^2}.
\]
Every orthonormal sequence is weakly null, so
\((\Theta_j)\) is a singular Weyl sequence at \(\lambda\).
Thus every \(\lambda>0\) belongs to the Weyl essential spectrum.  Together
with the zero sequence and nonnegativity, this proves
\begin{equation}
 \sigma(\widetilde{\mathscr M}_\eta)
 =\sigma_{\mathrm{ess}}(\widetilde{\mathscr M}_\eta)
 =[0,\infty).
\label{native-di-spectrum-essential}
\end{equation}

\paragraph{Real analyticity, absence of flat bands, and no eigenvalues.}
Equation \eqref{cusp-native-full-flux-Xi-identity} writes
\[
 \Xi_R=\sqrt{S_R^2+24D_R}
\]
with a strictly positive real-analytic radicand.  The positive square-root
function is real analytic on \((0,\infty)\).  Hence
\(\Xi_R\), \(\mathcal B_{R,+}\), \(\mathcal B_{R,-}\), and every
\(E_{\mathbf n}(R)\) are real analytic on the connected interval \(I\).
For each fixed \(\mathbf n\),
\[
 E_{\mathbf n}(R)>0\quad(R<\infty),
 \qquad
 \lim_{R\to\infty}E_{\mathbf n}(R)=0.
\]
It follows that no band is constant: a constant band would have constant
value equal to its limit zero, contradicting strict positivity at every
finite \(R\).  This proves the absence of flat bands, rather than assuming
it.

Fix \(\lambda\in\mathbb R\).  For every \(\mathbf n\), the real-analytic
function \(E_{\mathbf n}-\lambda\) is not identically zero, so its zero set
in \(I\) is discrete and has Lebesgue measure zero.  Since
\(\mathbb Z_{\ge0}^2\) is countable,
\begin{equation}
 Z_\lambda
 :=\{R\in I:\lambda\in\sigma_{\mathrm p}(A_R)\}
 =\bigcup_{\mathbf n\in\mathbb Z_{\ge0}^2}
   \{R:E_{\mathbf n}(R)=\lambda\}
\label{native-di-null-fibre-level-set}
\end{equation}
is a null set.

If
\(\widetilde{\mathscr M}_\eta\Psi=\lambda\Psi\), then
\((A_R-\lambda)\Psi(R)=0\) for almost every \(R\).  At every point where
\(\Psi(R)\ne0\), Proposition
\ref{cusp-native-full-flux-landau-theorem} forces \(R\in Z_\lambda\).
Thus the essential support of \(\Psi\) is contained in the null set
\eqref{native-di-null-fibre-level-set}, and \(\Psi=0\).  A self-adjoint
operator has only real eigenvalues, so
\[
 \sigma_{\mathrm p}(\widetilde{\mathscr M}_\eta)=\varnothing,
 \qquad
 \ker\widetilde{\mathscr M}_\eta=\{0\}.
\]
Self-adjoint operators have empty residual spectrum.  Combining this with
\eqref{native-di-spectrum-essential} proves that the entire spectrum is
continuous and gives \eqref{native-di-fixed-spectrum-result}.  Unitary
equivalence gives \eqref{native-di-varying-spectrum-result}.
\end{proof}

\subsubsection{Exact interpretive boundary}

\begin{remark}[No radial derivative and no interacting Yang--Mills claim]
The variable \(R\) in
\eqref{native-di-fixed-multiplication} and
\eqref{native-di-varying-multiplication} is a direct-integral parameter.
Neither operator contains a \(R\)-derivative.  In particular, there is no
term such as \(-\partial_R^2\), no radial Sobolev condition, no cusp
boundary condition, and no interaction between distinct \(R\)-fibres.
This is why characteristic functions of radial intervals legitimately
occur in the maximal-domain Weyl sequences above.

The theorem is therefore a complete spectral statement for the specified
auxiliary decomposable family of native full-flux magnetic Laplacians.  It
does not construct, identify, or determine the spectrum or mass gap of an
interacting quantum Yang--Mills Hamiltonian, and it makes no
Osterwalder--Schrader reconstruction claim.
\end{remark}



% Standalone-preprint adaptation of cusp_native_full_flux_suN.tex.
% Canonical source SHA-256: 165EA6F8009E4BE6371DD438076C70CDEA1D6890D013469968D6AFEE16DE5255.
% The canonical radial-height symbol T was renamed R consistently; transpose
% notation was retained exactly.

% Canonical split-Cartan SU(r) extension module, 2026-09-03.
% This is a delta to cusp_native_full_flux_magnetic.tex.  It imports that
% module's native torus, full-flux line bundle, orientations, exact magnetic
% frequencies, and SU(2) calculation.  It deliberately does not duplicate
% the radial direct-integral operator or the charge-two SU(2) proof.

\section{The native full-flux split-Cartan \texorpdfstring{\(SU(r)\)}{SU(r)} extension}
\label{cusp-native-full-flux-suN-section}

This section fixes an integer \(r\geq2\) and extends the full-flux
construction of \cref{cusp-native-full-flux-section} from its already proved
split \(SU(2)\) case to arbitrary split-Cartan \(SU(r)\) backgrounds.  Every
object remains on the same regular flat four-torus
\(F_R=\mathbb R^4/P_R\mathbb Z^4\), with the unrescaled Euclidean quotient
metric.  In particular, the integral two-form \(\eta_R\), the line bundle
\(\mathscr L_{\eta,R}\), the two orientations, and the exact quantities
\[
 \begin{gathered}
  L_R,\quad q_R,\quad m_R,\quad D_R=L_R q_R+6m_R^2>0,\\
  S_R=6L_R-q_R,\qquad
  \Xi_R=\sqrt{(6L_R+q_R)^2+144m_R^2}.
 \end{gathered}
\]
are precisely those of
\cref{cusp-native-full-flux-canonical-scalars,%
cusp-native-full-flux-S-Xi}.  The charge-one frequencies are
\begin{equation}
 \mathcal B_{R,+}=\frac{\pi(S_R+\Xi_R)}{D_R},
 \qquad
 \mathcal B_{R,-}=\frac{\pi(\Xi_R-S_R)}{D_R},
 \label{cusp-native-full-flux-suN-signed-frequencies}
\end{equation}
and we shall use the globally ordered notation
\begin{equation}
 \begin{aligned}
  \mathcal B_{R,\max}
  &=\max\{\mathcal B_{R,+},\mathcal B_{R,-}\},&
  \mathcal B_{R,\min}
  &=\min\{\mathcal B_{R,+},\mathcal B_{R,-}\},\\
  \rho_R^*&=\frac{\mathcal B_{R,\max}}{\mathcal B_{R,\min}}.
 \end{aligned}
 \label{cusp-native-full-flux-suN-ordered-frequencies}
\end{equation}
Thus every assertion below remains valid on either side of \(S_R=0\).

\subsection{The split bundle, determinant, and Chern number}

Fix a nonzero trace-zero integral weight vector
\begin{equation}
 \mathbf m=(m_1,\ldots,m_r)\in\mathbb Z^r,
 \qquad \sum_{j=1}^r m_j=0,
 \qquad s_2(\mathbf m):=\sum_{j=1}^r m_j^2.
 \label{cusp-native-full-flux-suN-weight-vector}
\end{equation}
Define
\begin{equation}
 \mathscr E_{\mathbf m,R}:=
 \bigoplus_{j=1}^r\mathscr L_{\eta,R}^{\otimes m_j},
 \qquad
 H_{\mathbf m}:=i\operatorname{diag}(m_1,\ldots,m_r)
 \in\mathfrak{su}(r).
 \label{cusp-native-full-flux-suN-split-bundle}
\end{equation}

\begin{proposition}[Parallel determinant trivialization and exact \(c_2\)]
\label{cusp-native-full-flux-suN-topology}
The ordered splitting in
\cref{cusp-native-full-flux-suN-split-bundle} has a canonical parallel
determinant trivialization.  With \(c=[\eta_R]\),
\begin{equation}
 c_1(\mathscr E_{\mathbf m,R})=0,
 \qquad
 c_2(\mathscr E_{\mathbf m,R})
 =-\frac{s_2(\mathbf m)}2c^2.
 \label{cusp-native-full-flux-suN-Chern-classes}
\end{equation}
The marked and physical-complex Chern numbers are respectively
\begin{equation}
 \left\langle c_2,[F_R]_{\mathrm{mark}}\right\rangle=-6s_2(\mathbf m),
 \qquad
 k_{\mathbf m}:=
 \left\langle c_2,[F_R]_{x\text{-}\mathrm{complex}}\right\rangle
 =6s_2(\mathbf m)>0.
 \label{cusp-native-full-flux-suN-Chern-numbers}
\end{equation}
Moreover \(s_2(\mathbf m)\) is even.  Hence this direct-sum family realizes
only \(k_{\mathbf m}\in12\mathbb Z_{>0}\), not all topological \(SU(r)\)
sectors.
\end{proposition}

\begin{proof}
Let \(j_n^\eta(y)\) denote the factor of automorphy in
\cref{cusp-native-full-flux-cocycle}.  The transition matrix of the ordered
sum is
\[
 \operatorname{diag}\bigl((j_n^\eta)^{m_1},\ldots,
 (j_n^\eta)^{m_r}\bigr).
\]
Its determinant is \((j_n^\eta)^{\sum_j m_j}=1\) identically.  The tensor
product of the ordered unit frames therefore defines a global determinant
section, and the determinant connection has connection form proportional to
\(\sum_j m_j=0\); the section is parallel.

The splitting principle gives
\[
 c(\mathscr E_{\mathbf m,R})=\prod_{j=1}^r(1+m_j c),
 \qquad
 c_2=\sum_{i<j}m_i m_j c^2.
\]
Since
\[
 2\sum_{i<j}m_i m_j
 =\left(\sum_j m_j\right)^2-\sum_j m_j^2=-s_2(\mathbf m),
\]
this proves \cref{cusp-native-full-flux-suN-Chern-classes}.  The two
evaluations follow from
\cref{cusp-native-full-flux-two-orientation-pairings}; retaining its
orientation reversal is essential for the sign.  Finally
\[
 s_2(\mathbf m)=\sum_j m_j^2\equiv\sum_j m_j\equiv0\pmod2,
\]
because \(a^2\equiv a\pmod2\).  This proves the divisibility statement.
\end{proof}

The diagonal connection induced by the summands is
\begin{equation}
 \nabla^{\mathbf m,R}
 =\mathrm d-2\pi(6y^1\,\mathrm dy^4+y^2\,\mathrm dy^3)H_{\mathbf m},
 \qquad
 F_{\mathbf m,R}=-2\pi H_{\mathbf m}\eta_R.
 \label{cusp-native-full-flux-suN-connection-curvature}
\end{equation}
We retain the invariant inner product and functional normalization
\begin{equation}
 \begin{gathered}
  \langle X,Y\rangle_{\mathfrak{su}(r)}
  =-\frac12\operatorname{Tr}(XY),
  \qquad
  \lVert H_{\mathbf m}\rVert^2=\frac12s_2(\mathbf m),\\
  \operatorname{YM}_R(A)=\frac1{2g_{\mathrm{YM}}^2}
  \int_{F_R}|F_A|^2\,\mathrm d\operatorname{vol}_R.
 \end{gathered}
 \label{cusp-native-full-flux-suN-normalizations}
\end{equation}

\begin{theorem}[Yang--Mills criticality and exact action]
\label{cusp-native-full-flux-suN-action-theorem}
The connection \(\nabla^{\mathbf m,R}\) is non-flat, reducible, and
Yang--Mills.  Its action is
\begin{align}
 \operatorname{YM}_R(\nabla^{\mathbf m,R})
 &=\frac{\pi^2s_2(\mathbf m)}{g_{\mathrm{YM}}^2D_R}
   \bigl(q_R^2+36L_R^2+72m_R^2\bigr)
 \label{cusp-native-full-flux-suN-action-coordinate}\\
 &=\frac{2\pi^2k_{\mathbf m}}{g_{\mathrm{YM}}^2}
   +\frac{\pi^2s_2(\mathbf m)S_R^2}
          {g_{\mathrm{YM}}^2D_R}.
 \label{cusp-native-full-flux-suN-action-excess}
\end{align}
On the canonical cusp tail,
\begin{equation}
 \operatorname{YM}_R(\nabla^{\mathbf m,R})
 =\frac{\pi^2s_2(\mathbf m)}{g_{\mathrm{YM}}^2}
 \left(37+\frac{35B}{R}+O(R^{-2})\right).
 \label{cusp-native-full-flux-suN-action-asymptotic}
\end{equation}
\end{theorem}

\begin{proof}
Both \(H_{\mathbf m}\) and \(\eta_R\) are parallel for the diagonal
connection and flat quotient metric, and their brackets vanish.  Hence
\(d_{\nabla^{\mathbf m,R}}^*F_{\mathbf m,R}=0\).  Non-flatness follows
from \(\mathbf m\ne0\) and \(\eta_R\ne0\).

In the physical orthonormal frame used in
\cref{cusp-native-full-flux-normal-form},
\[
 |\eta_R|^2
 =\frac{\mathcal B_{R,+}^2+\mathcal B_{R,-}^2}{4\pi^2}.
\]
Using \cref{cusp-native-full-flux-suN-normalizations},
\(\operatorname{vol}(F_R)=D_R\), and the exact frequency square sum from
the native module yields
\cref{cusp-native-full-flux-suN-action-coordinate}.  The identity
\[
 q_R^2+36L_R^2+72m_R^2
 =12(L_R q_R+6m_R^2)+(6L_R-q_R)^2
 =12D_R+S_R^2
\]
together with \(k_{\mathbf m}=6s_2(\mathbf m)\), proves
\cref{cusp-native-full-flux-suN-action-excess}.  Substituting the retained
expansions of \(L_R,q_R,m_R,D_R\), without suppressing \(H,B,m\), gives
\[
 \frac{q_R^2+36L_R^2+72m_R^2}{D_R}
 =37+\frac{35B}{R}+O(R^{-2}),
\]
which proves the last formula.
\end{proof}

\subsection{Every nonzero integral root charge}

The native module proves the charge-one ladder decomposition and the
charge-two specialization used by its \(SU(2)\) theorem.  The following
lemma records, and proves, the exact integer-charge extension needed here.

\begin{lemma}[Charge-\(d\) Landau decomposition]
\label{cusp-native-full-flux-suN-charge-d-Landau}
Let \(d\in\mathbb Z\setminus\{0\}\), put \(q=|d|\), and equip
\(\mathscr L_{\eta,R}^{\otimes d}\) with its tensor-power connection.  Its
scalar covariant Laplacian has joint levels
\begin{equation}
 q\bigl((2a+1)\mathcal B_{R,+}+(2b+1)\mathcal B_{R,-}\bigr),
 \qquad a,b\in\mathbb Z_{\ge0},
 \label{cusp-native-full-flux-suN-charge-d-levels}
\end{equation}
and every joint space has complex dimension \(6q^2\).  If different pairs
\((a,b)\) give the same numerical value, their orthogonal joint spaces are
to be summed.
\end{lemma}

\begin{proof}
The curvature is \(-2\pi i d\eta_R\).  In the polar construction of
\cref{cusp-native-full-flux-landau-theorem}, replace the auxiliary complex
structure by
\(J_{d,R}=\operatorname{sgn}(d)J_{\eta,R}\).  Then \(d\eta_R\) is positive
of type \((1,1)\), so the vacuum is the holomorphic-section space of the
positive bundle \(\mathscr L_{\eta,R}^{\otimes d}\) on that auxiliary
complex torus.  Its tangent and canonical bundles are trivial.  Kodaira
vanishing and Hirzebruch--Riemann--Roch give
\[
 \dim_{\mathbb C}H^0
 =\frac12\int_{[F_R]_{d\eta\text{-}\mathrm{symp}}}(d\eta_R)^2
 =\frac12q^2\int_{[F_R]_{\mathrm{mark}}}\eta_R^2
 =6q^2.
\]
Here changing \(d\) to \(-d\) does not change the four-dimensional
symplectic orientation.  Scaling the curvature by \(d\) scales both
positive oscillator frequencies by \(q\), while the two canonical
commutation relations and the lowering argument in
\cref{cusp-native-full-flux-landau-theorem} are unchanged after this
rescaling.  Thus every joint level is unitarily obtained from the vacuum and
has the asserted energy and dimension.  Orthogonal joint spaces remain a
direct sum when their numerical energies coincide.
\end{proof}

\subsection{Closed gauge-fixed and Coulomb operators}

Write \(A=\nabla^{\mathbf m,R}\) in this subsection.  Let
\begin{equation}
 \mathscr H_{\mathbf m,R}^1
 =L^2\Omega^1(F_R;\operatorname{ad}\mathscr E_{\mathbf m,R}),
 \qquad
 \mathscr H_{\mathbf m,R}^0
 =L^2\Omega^0(F_R;\operatorname{ad}\mathscr E_{\mathbf m,R})
 \label{cusp-native-full-flux-suN-real-Hilbert-spaces}
\end{equation}
as real Hilbert spaces.  Define the maximal distributional codifferential
by
\begin{align}
 \operatorname{Dom}D_{\mathbf m,R}^{\max}
 =\{a\in\mathscr H_{\mathbf m,R}^1:\;&
 \exists u\in\mathscr H_{\mathbf m,R}^0\text{ such that}\notag\\[-1mm]
 &\langle a,d_A\phi\rangle_{L^2}
 =\langle u,\phi\rangle_{L^2}
 \quad\text{for every }\notag\\[-1mm]
 &\phi\in C^\infty(F_R;\operatorname{ad}\mathscr E_{\mathbf m,R})\},
 \label{cusp-native-full-flux-suN-max-codiff-domain}\\
 D_{\mathbf m,R}^{\max}a={}&u,
 \qquad
 \mathscr S_{\mathbf m,R}:=\ker D_{\mathbf m,R}^{\max}.
 \label{cusp-native-full-flux-suN-Coulomb-slice}
\end{align}
Density of smooth test scalars makes \(u\) unique.  If
\(a_j\to a\) and \(D_{\mathbf m,R}^{\max}a_j\to u\) in \(L^2\), passage to
the limit in the defining identity gives
\(D_{\mathbf m,R}^{\max}a=u\).  Thus \(D_{\mathbf m,R}^{\max}\) is closed
and \(\mathscr S_{\mathbf m,R}\) is a closed subspace.

Let \(\Delta_A^{(0)}\) be the nonnegative self-adjoint scalar operator of
the closed form
\begin{equation}
 \mathfrak q_{\mathbf m,R}^{(0)}[\phi,\psi]
 =\langle d_A\phi,d_A\psi\rangle_{L^2},
 \qquad
 \operatorname{Dom}\mathfrak q_{\mathbf m,R}^{(0)}=H^1\Omega^0.
 \label{cusp-native-full-flux-suN-scalar-form}
\end{equation}
It has domain \(H^2\Omega^0\) and equals \(d_A^*d_A\).  In a constant
orthonormal frame set
\begin{equation}
 (V_{\mathbf m,R}a)_\nu
 =-2\sum_{\mu=1}^4[F_{\mathbf m,R,\nu\mu},a_\mu].
 \label{cusp-native-full-flux-suN-spin-endomorphism}
\end{equation}
Invariance of the compact Lie-algebra inner product and
\(F_{\nu\mu}=-F_{\mu\nu}\) show directly that \(V_{\mathbf m,R}\) is a
bounded symmetric endomorphism.  Define
\begin{equation}
 \mathfrak j_{\mathbf m,R}^{\mathrm{gf}}[a,b]
 =\langle\nabla_A a,\nabla_A b\rangle_{L^2}
  +\langle V_{\mathbf m,R}a,b\rangle_{L^2},
 \qquad
 \operatorname{Dom}\mathfrak j_{\mathbf m,R}^{\mathrm{gf}}
 =H^1\Omega^1.
 \label{cusp-native-full-flux-suN-Jacobi-form}
\end{equation}

\begin{proposition}[Closed realization and exact Coulomb reduction]
\label{cusp-native-full-flux-suN-Coulomb-reduction}
The form in \cref{cusp-native-full-flux-suN-Jacobi-form} is densely defined,
closed, symmetric, and bounded below.  Its self-adjoint operator has compact
resolvent, domain \(H^2\Omega^1\), and expression
\begin{equation}
 (\mathcal J_{\mathbf m,R}^{\mathrm{gf}}a)_\nu
 =-\sum_{\mu=1}^4\nabla_{A,e_\mu}\nabla_{A,e_\mu}a_\nu
  -2\sum_{\mu=1}^4[F_{\mathbf m,R,\nu\mu},a_\mu].
 \label{cusp-native-full-flux-suN-Jacobi-expression}
\end{equation}
For every \(z\in\mathbb C\setminus\mathbb R\) and
\(f\in\operatorname{Dom}D_{\mathbf m,R}^{\max}\),
\begin{equation}
 D_{\mathbf m,R}^{\max}
 (\mathcal J_{\mathbf m,R}^{\mathrm{gf}}-z)^{-1}f
 =(\Delta_A^{(0)}-z)^{-1}D_{\mathbf m,R}^{\max}f.
 \label{cusp-native-full-flux-suN-resolvent-intertwining}
\end{equation}
Consequently \(\mathscr S_{\mathbf m,R}\) reduces
\(\mathcal J_{\mathbf m,R}^{\mathrm{gf}}\), and
\begin{equation}
 \mathcal J_{\mathbf m,R}
 :=\mathcal J_{\mathbf m,R}^{\mathrm{gf}}
   \big|_{\mathscr S_{\mathbf m,R}},
 \qquad
 \operatorname{Dom}\mathcal J_{\mathbf m,R}
 =H^2\Omega^1\cap\mathscr S_{\mathbf m,R}
 \label{cusp-native-full-flux-suN-Coulomb-operator}
\end{equation}
is self-adjoint, with compact resolvent and closed form
\(\mathfrak j_{\mathbf m,R}\) equal to the restriction of
\(\mathfrak j_{\mathbf m,R}^{\mathrm{gf}}\) to
\(H^1\Omega^1\cap\mathscr S_{\mathbf m,R}\).  The physical Yang--Mills
Hessian on this slice is
\begin{equation}
 \operatorname{YM}_R''\big|_A
 =g_{\mathrm{YM}}^{-2}\mathfrak j_{\mathbf m,R}.
 \label{cusp-native-full-flux-suN-physical-Hessian}
\end{equation}
Thus an eigenvalue \(\nu\) of \(\mathcal J_{\mathbf m,R}\) is the physical
Hessian eigenvalue \(\nu/g_{\mathrm{YM}}^2\), while its sign and Morse-index
contribution are unchanged.
\end{proposition}

\begin{proof}
The first term of
\cref{cusp-native-full-flux-suN-Jacobi-form}, after addition of the
\(L^2\)-norm, is an equivalent \(H^1\)-norm.  The second term is a bounded
symmetric form.  The bounded perturbation theorem for closed forms and the
representation theorem therefore give a lower-bounded self-adjoint
operator.  Its principal symbol is \(|\xi|^2\) times the identity, so
elliptic regularity gives the \(H^2\) domain.  Rellich compactness on the
compact torus gives compact resolvent.

The two-parameter curvature expansion and the flat Weitzenb\"ock identity,
with the inner product and functional normalization in
\cref{cusp-native-full-flux-suN-normalizations}, give
\[
 \operatorname{YM}_R''\big|_A(a,b)
 +g_{\mathrm{YM}}^{-2}\langle d_A^*a,d_A^*b\rangle
 =g_{\mathrm{YM}}^{-2}
   \mathfrak j_{\mathbf m,R}^{\mathrm{gf}}[a,b].
\]
This proves the Hessian normalization after restriction to Coulomb gauge.

It remains to prove that the restriction is an operator rather than a formal
constraint.  For a smooth adjoint scalar \(\phi\), write
\(\nabla_\mu=\nabla_{A,e_\mu}\).  Commuting all three derivatives and using
the Yang--Mills equation gives
\begin{align*}
 &(\mathcal J_{\mathbf m,R}^{\mathrm{gf}}d_A\phi
       -d_A\Delta_A^{(0)}\phi)_\nu\\
 &\quad=-\nabla_\mu\nabla_\mu\nabla_\nu\phi
       +\nabla_\nu\nabla_\mu\nabla_\mu\phi
       -2[F_{\nu\mu},\nabla_\mu\phi]\\
 &\quad=-[\nabla_\mu F_{\mu\nu},\phi]
       =[(d_A^*F_A)_\nu,\phi]=0.
\end{align*}
Formal adjunction yields, on smooth one-forms,
\begin{equation}
 \mathcal J_{\mathbf m,R}^{\mathrm{gf}}d_A
 =d_A\Delta_A^{(0)},
 \qquad
 d_A^*\mathcal J_{\mathbf m,R}^{\mathrm{gf}}
 =\Delta_A^{(0)}d_A^*.
 \label{cusp-native-full-flux-suN-formal-intertwining}
\end{equation}
For
\(u=(\mathcal J_{\mathbf m,R}^{\mathrm{gf}}-z)^{-1}f\), the operator domain
gives \(u\in H^2\), hence \(D_{\mathbf m,R}^{\max}u\in H^1\).  Applying the
second identity distributionally to
\((\mathcal J_{\mathbf m,R}^{\mathrm{gf}}-z)u=f\) gives
\[
 (\Delta_A^{(0)}-z)D_{\mathbf m,R}^{\max}u
 =D_{\mathbf m,R}^{\max}f.
\]
Scalar elliptic regularity and the scalar resolvent give
\cref{cusp-native-full-flux-suN-resolvent-intertwining}.  The two conjugate
resolvents preserve \(\mathscr S_{\mathbf m,R}\); adjunction then makes its
orthogonal complement invariant as well.  The slice is reducing, so the
self-adjoint part has the operator and form domains displayed above.
\end{proof}

\subsection{Rootwise spectrum, coincidences, index, and nullity}

For \(i\ne j\), the matrix unit \(E_{ij}\) obeys
\[
 [H_{\mathbf m},E_{ij}]=i(m_i-m_j)E_{ij},
\]
so its complex root line is
\begin{equation}
 \operatorname{Hom}(\mathscr L_{\eta,R}^{m_j},
                     \mathscr L_{\eta,R}^{m_i})
 =\mathscr L_{\eta,R}^{m_i-m_j},
 \qquad q_{ij}:=|m_i-m_j|.
 \label{cusp-native-full-flux-suN-root-charge}
\end{equation}
The factor \(q_{ij}\) is the complete charge normalization; for the existing
\(SU(2)\) weights \((1,-1)\), it gives \(q_{12}=2\).

Let the distinct numerical values among \(m_1,\ldots,m_r\) have
multiplicities \(r_1,\ldots,r_\ell\).  The neutral centralizer is
\begin{equation}
 \mathfrak z_{\mathbf m}
 =\{X\in\mathfrak{su}(r):[H_{\mathbf m},X]=0\}
 \cong\mathfrak s\!\left(\bigoplus_{a=1}^{\ell}\mathfrak u(r_a)\right),
 \qquad
 \dim_{\mathbb R}\mathfrak z_{\mathbf m}
 =\sum_{a=1}^{\ell}r_a^2-1.
 \label{cusp-native-full-flux-suN-centralizer}
\end{equation}

\begin{theorem}[Complete negative spectrum and numerical multiplicities]
\label{cusp-native-full-flux-suN-Morse-theorem}
Put
\begin{equation}
 N_-(R):=\#\{n\in\mathbb Z_{\ge0}:2n+1<\rho_R^*\}
 =\left\lceil\frac{\rho_R^*-1}{2}\right\rceil.
 \label{cusp-native-full-flux-suN-N-minus}
\end{equation}
For each positive root \(i<j\) with \(q_{ij}>0\), the complete negative
spectrum of the normalized Coulomb operator is
\begin{equation}
 \nu_{ij,n}(R)
 =q_{ij}\bigl((2n+1)\mathcal B_{R,\min}
                  -\mathcal B_{R,\max}\bigr),
 \qquad 0\le n<N_-(R).
 \label{cusp-native-full-flux-suN-negative-bands}
\end{equation}
The contribution of that positive-root pair to each displayed band has real
dimension \(12q_{ij}^2\).  For an actual numerical eigenvalue
\(\lambda<0\), its full real multiplicity is therefore
\begin{equation}
 \operatorname{mult}_{\mathbb R}(\lambda)
 =\sum_{\substack{i<j,\ q_{ij}>0,\ 0\le n<N_-(R)\\
                   \nu_{ij,n}(R)=\lambda}}
   12q_{ij}^2.
 \label{cusp-native-full-flux-suN-coincident-multiplicity}
\end{equation}
In particular, coincident root charges or accidental coincidences between
different \(q_{ij},n\) add rather than being counted as one band.  The real
Morse index is
\begin{equation}
 \boxed{
 \operatorname{ind}_{\mathbb R}\mathcal J_{\mathbf m,R}
 =12N_-(R)\sum_{i<j}(m_i-m_j)^2
 =12rN_-(R)s_2(\mathbf m).}
 \label{cusp-native-full-flux-suN-index}
\end{equation}
Every eigenvalue in
\cref{cusp-native-full-flux-suN-negative-bands} is divided by
\(g_{\mathrm{YM}}^2\) for the physical Hessian.
\end{theorem}

\begin{proof}
On the \(E_{ij}\) line, the scalar oscillator is the charge-\(q_{ij}\)
operator of \cref{cusp-native-full-flux-suN-charge-d-Landau}.  In a
skew-normal orthonormal coframe, the four circular one-form polarizations
diagonalize \cref{cusp-native-full-flux-suN-spin-endomorphism}; the two
polarizations in a plane of charge-\(q_{ij}\) frequency
\(q_{ij}\mathcal B\) shift the scalar energy by
\(\pm2q_{ij}\mathcal B\).  The unfavorable polarizations are positive.
For the favorable polarization in the high-frequency plane, high-plane
oscillator number zero leaves
\[
 q_{ij}\bigl((2n+1)\mathcal B_{R,\min}
             -\mathcal B_{R,\max}\bigr).
\]
Increasing its high-plane oscillator number makes it positive.  The
favorable low-frequency polarization has bottom
\(q_{ij}(\mathcal B_{R,\max}-\mathcal B_{R,\min})\ge0\), and increasing
either number cannot lower it.  This proves that the list is exhaustive.

Every joint scalar space on one complex root line has dimension
\(6q_{ij}^2\).  The opposite root is its anti-linear conjugate.  The pair is
the complexification of a real root subspace, so its real dimension is
\(12q_{ij}^2\).  Root subspaces and joint oscillator spaces are orthogonal;
hence equality of numerical eigenvalues produces precisely the direct sum
in \cref{cusp-native-full-flux-suN-coincident-multiplicity}.

The strict inequality in
\cref{cusp-native-full-flux-suN-N-minus} counts the integers correctly even
when \(\rho_R^*\) is an odd integer.  Summing dimensions gives the first
index formula.  Finally,
\[
 \sum_{i<j}(m_i-m_j)^2
 =r\sum_i m_i^2-\left(\sum_i m_i\right)^2
 =rs_2(\mathbf m),
\]
which gives the second formula.

It remains to justify Coulomb gauge.  If
\(\mathcal J_{\mathbf m,R}^{\mathrm{gf}}a=\lambda a\) with
\(\lambda<0\), elliptic regularity and
\cref{cusp-native-full-flux-suN-formal-intertwining} give
\[
 \Delta_A^{(0)}d_A^*a=\lambda d_A^*a.
\]
The scalar operator is nonnegative, so \(d_A^*a=0\).  Thus the negative
gauge-fixed and Coulomb eigenspaces coincide, completing the proof.
\end{proof}

\begin{theorem}[Exact neutral and wall nullities]
\label{cusp-native-full-flux-suN-nullity-theorem}
Away from an odd wall, the neutral Coulomb kernel has real dimension
\begin{equation}
 \operatorname{null}_{\mathrm{neutral}}\mathcal J_{\mathbf m,R}
 =4\left(\sum_{a=1}^{\ell}r_a^2-1\right).
 \label{cusp-native-full-flux-suN-neutral-nullity}
\end{equation}
At a wall \(\rho_R^*=2n+1>1\), the nonzero-root part adds
\begin{equation}
 12\sum_{i<j}q_{ij}^2=12r s_2(\mathbf m)
 \label{cusp-native-full-flux-suN-odd-wall-nullity}
\end{equation}
real zero modes.  At the equality wall \(S_R=0\), equivalently
\(\rho_R^*=1\), both favorable plane polarizations vanish and the nonzero
root contribution doubles to
\begin{equation}
 24\sum_{i<j}q_{ij}^2=24r s_2(\mathbf m).
 \label{cusp-native-full-flux-suN-equality-wall-nullity}
\end{equation}
The total nullity at a wall is the sum of the neutral term and the stated
root term.
\end{theorem}

\begin{proof}
On \(\mathfrak z_{\mathbf m}\), both the induced connection commutator and
the spin endomorphism vanish.  The operator is therefore the ordinary
coefficientwise Laplacian on one-forms.  Its kernel on a flat four-torus is
the space of constant one-forms, all of which are Coulomb.  Tensoring the
four-dimensional real harmonic-one-form space with
\(\mathfrak z_{\mathbf m}\) proves
\cref{cusp-native-full-flux-suN-neutral-nullity}.

If \(\rho_R^*=2n+1>1\), the high-plane favorable polarization at high
oscillator number zero and low oscillator number \(n\) is the unique zero
branch for each nonzero root.  Its positive-root-pair real dimension is
\(12q_{ij}^2\).  If \(S_R=0\), the two frequencies agree and the lowest
favorable polarization in each of the two planes is zero, so this dimension
doubles.  These charged zero modes are Coulomb: the intertwining identity
shows that their divergences lie in the kernel of the charged scalar
Laplacian, but that scalar Laplacian has strictly positive ground energy by
\cref{cusp-native-full-flux-suN-charge-d-levels}.  Summing over positive
roots and using
\(\sum_{i<j}q_{ij}^2=rs_2(\mathbf m)\) proves the result.
\end{proof}

For \(n\ge1\), the \(n\)-th odd wall is equivalently
\begin{equation}
 \rho_R^*=2n+1
 \iff n\Xi_R=(n+1)|S_R|
 \iff 24n^2D_R=(2n+1)S_R^2.
 \label{cusp-native-full-flux-suN-wall-equations}
\end{equation}
The final polynomial equation must be accompanied by the unsigned ordering
in the middle expression if one needs to distinguish the two sides
\(S_R>0\) and \(S_R<0\).  The \(n=0\) wall is exactly \(S_R=0\).

On the sufficiently deep cusp tail,
\(S_R>0\), \(5<\rho_R^*<7\), and \(N_-(R)=3\).  Hence
\begin{equation}
 \operatorname{ind}_{\mathbb R}\mathcal J_{\mathbf m,R}
 =36r s_2(\mathbf m)=6r k_{\mathbf m},
 \label{cusp-native-full-flux-suN-deep-index}
\end{equation}
and the three negative branches have the retained asymptotics
\begin{equation}
 \nu_{ij,n}(R)
 =\frac{q_{ij}\pi}{R}(4n-10)+O(R^{-2}),
 \qquad n=0,1,2.
 \label{cusp-native-full-flux-suN-negative-asymptotics}
\end{equation}
For example, the genuinely \(SU(3)\) data are
\begin{center}
\renewcommand{\arraystretch}{1.2}
\begin{tabular}{@{}lcccc@{}}
\toprule
weights & nonzero \(q_{ij}\) & \(k_{\mathbf m}\) & deep index
& neutral nullity\\
\midrule
\((1,0,-1)\) & \(1,1,2\) & \(12\) & \(216\) & \(8\)\\
\((1,1,-2)\) & \(3,3\) & \(36\) & \(648\) & \(16\)\\
\((2,1,-3)\) & \(1,4,5\) & \(84\) & \(1512\) & \(8\)\\
\bottomrule
\end{tabular}
\end{center}
The middle row exhibits both required edge phenomena: one neutral root from
the repeated weight and coincident numerical bands from the two charge-three
positive roots.  The specialization \(r=2\), \(\mathbf m=(1,-1)\) is
exactly the already proved native \(SU(2)\) theorem and is not reproduced
here.

\subsection{Stability boundary, possible wall bifurcation, and exact scope}

\begin{proposition}[Where the split background is, and is not, unstable]
\label{cusp-native-full-flux-suN-stability-boundary}
If \(S_R\ne0\), then \(\rho_R^*>1\), \(N_-(R)\ge1\), and every nonzero
positive root contributes a negative Coulomb direction.  Since a nonzero
trace-zero weight vector has at least one nonzero root,
\(\nabla^{\mathbf m,R}\) is then not a local minimum modulo gauge.  This
applies in particular throughout the sufficiently deep cusp tail.

If \(S_R=0\), the negative root spectrum is empty.  In that case
\cref{cusp-native-full-flux-suN-action-excess} attains the topological value
\(2\pi^2k_{\mathbf m}/g_{\mathrm{YM}}^2\), while
\cref{cusp-native-full-flux-suN-equality-wall-nullity} supplies root zero
modes.  No assertion of strict stability follows because the Hessian is
degenerate.
\end{proposition}

\begin{proof}
The first paragraph follows immediately from
\cref{cusp-native-full-flux-suN-ordered-frequencies,%
cusp-native-full-flux-suN-Morse-theorem}.  For the second,
\(S_R=0\) makes the two frequencies equal and hence \(\rho_R^*=1\); the
strict negative-band inequality has no solution.  The action and nullity
statements are exactly
\cref{cusp-native-full-flux-suN-action-excess,%
cusp-native-full-flux-suN-nullity-theorem}.
\end{proof}

A negative eigenspace is a space of descent directions, not a
Lyapunov--Schmidt bifurcation kernel.  A local bifurcation analysis must
instead be posed at one of the zero-mode parameters in
\cref{cusp-native-full-flux-suN-nullity-theorem}.  It must first quotient or
slice the ever-present centralizer-valued harmonic kernel, then identify the
remaining kernel and cokernel of the gauge-fixed nonlinear Yang--Mills map,
and finally compute the parameter transversality and nonlinear obstruction
coefficients.  The present linear calculation supplies the exact candidate
kernel at each wall; it does not assert existence of a nearby irreducible
branch.

The theorem proved here is strictly lower-rung and fibrewise.  It constructs
explicit reducible split-Cartan \(SU(r)\) Yang--Mills critical points on the
already existing regular flat four-torus fibres and determines their
topology, action, normalized Jacobi spectrum, physical Hessian scaling,
Morse index, and nullity.  It constructs no new base, spacetime, cone,
\(S^{12}\), \(S^{24}\), transfer matrix, or interacting quantum Hamiltonian.
The radial multiplication operator already treated in the canonical native
module is not repeated here.

There is also an exact obstruction to interpreting the nonzero-flux bundle
as the restriction of a bundle on a proposed higher carrier.  If
\(i:F_R\to Y\) is any map and \(H^4(Y;\mathbb Z)=0\), then every complex
bundle \(\mathscr V\to Y\) has \(c_2(\mathscr V)=0\), so naturality gives
\[
 c_2(i^*\mathscr V)=i^*c_2(\mathscr V)=0.
\]
This contradicts
\cref{cusp-native-full-flux-suN-Chern-classes,%
cusp-native-full-flux-suN-Chern-numbers} for
\(\mathscr E_{\mathbf m,R}\).  Therefore this nonzero-\(c_2\) bundle cannot
extend over any carrier with vanishing fourth integral cohomology, including
a space diffeomorphic to \(S^6\), \(S^{12}\), or \(S^{24}\).  Changing the
gauge group changes the root system on the fibre; it does not by itself
produce the proposed higher rung or any conclusion about a physical mass
gap.




\section{The exceptional first-coordinate elliptic quotient and Landau sector}

The four period generators in $\C^2$ are
\begin{equation}\label{eq:generators}
 v_\gamma=(6\mu,\beta),\quad
 v_u=(\tau,\mu),\quad
 v_w=(1,0),\quad
 v_\delta=(0,1).
\end{equation}
Let $E_\tau=\C/(\Z\tau+\Z)$.

\begin{proposition}[Exact descent criterion]\label{prop:descent}
The first-coordinate projection $\operatorname{pr}_1:\C^2\to\C$ descends to a surjective holomorphic homomorphism
\[
 \overline{\operatorname{pr}}_1:F_z\longrightarrow E_\tau
\]
if and only if
\begin{equation}\label{eq:descent-condition}
 6\mu=a\tau+b
 \qquad\text{for some }(a,b)\in\Z^2.
\end{equation}
The pair $(a,b)$ is necessarily unique.
On this locus, put
\begin{equation}\label{eq:tau-f}
 \tau_f=a\mu-\beta.
\end{equation}
Then $\Im\tau_f=-D>0$ and the connected kernel is the elliptic curve
$E_{\tau_f}=\C/(\Z+\Z\tau_f)$, up to the harmless reversal of one lattice generator.
\end{proposition}

\begin{proof}
The first coordinates of \eqref{eq:generators} are $6\mu,\tau,1,0$.  Descent to $E_\tau$ is therefore equivalent to $6\mu\in\Z\tau+\Z$, which is \eqref{eq:descent-condition}; the images of $v_u,v_w$ prove surjectivity.  Uniqueness follows from $\Im\tau>0$.  The unimodular replacement
\[
 v_\gamma\longmapsto v_\gamma-av_u-bv_w=(0,\beta-a\mu)
\]
shows that the kernel lattice is generated by $1$ and $\beta-a\mu=-\tau_f$.  Taking imaginary parts in \eqref{eq:descent-condition} gives $6m=aT$, and hence
\[
 D=p-\frac{6m^2}{T}=p-am=\Im(\beta-a\mu)=-\Im\tau_f.
\]
Since $D<0$, $\tau_f\in\mathfrak h$.
\end{proof}

Define the marked exceptional set
\begin{equation}\label{eq:descent-locus}
 \mathcal D=\bigcup_{(a,b)\in\Z^2}
 \mathcal D_{a,b},
 \qquad
 \mathcal D_{a,b}=\{z:6\mu(z)-a\tau(z)-b=0\}.
\end{equation}

\begin{theorem}[No cusp accumulation and finiteness on the orbifold base]\label{thm:descent-finite}
Assume the transformation laws and extension properties in Section~1.
The generator laws act on the integral labels by
\begin{equation}\label{eq:descent-label-monodromy}
 g_0:(a,b)\longmapsto(a,a+b),
 \qquad
 g_2:(a,b)\longmapsto(-b,6+a).
\end{equation}
Consequently $\mathcal D$ is invariant under the triangle group and has a
well-defined image in the completed orbifold base.

\begin{enumerate}[label=(\roman*),leftmargin=2.3em]
\item There is a punctured neighborhood of the distinguished cusp containing no point of $\mathcal D$.
\item The image of $\mathcal D$ in the completed triangle-orbifold base is finite.
\end{enumerate}
\end{theorem}

\begin{proof}
If $z\in\mathcal D_{a,b}$, then the $g_0$ laws give
\[
 6\mu(g_0z)=a(\tau(g_0z)+1)+b
 =a\tau(g_0z)+(a+b),
\]
while the $g_2$ laws give
\[
 6\mu(g_2z)=6+a+\frac b{\tau(z)}
 =-b\tau(g_2z)+(6+a).
\]
Both label maps are bijections, and $g_0,g_2$ generate the triangle group.
This proves \eqref{eq:descent-label-monodromy} and invariance of $\mathcal D$.

Suppose first that $z_j\in\mathcal D_{a_j,b_j}$ approaches the cusp in a fixed cusp component.  Taking imaginary parts gives
\begin{equation}\label{eq:a-cusp}
 a_j=\frac{6\Im\mu(z_j)}{\Im\tau(z_j)}\longrightarrow0.
\end{equation}
Because $a_j\in\Z$, one has $a_j=0$ eventually.  Then
$6\mu(z_j)=b_j$.  The holomorphic cusp extension of $\mu$ makes it bounded, so the integers $b_j$ are bounded.  After passage to a subsequence, $b_j=b$ is constant.  The bounded holomorphic function $6\mu-b$ extends across the cusp and has zeros accumulating at the added interior point, hence vanishes identically.

No identity $6\mu=a\tau+b$ with fixed integers $a,b$ is possible.  Indeed, applying $g_0$ and using \eqref{eq:generator-laws} forces $a=0$.  The identity would then make $\mu=b/6$ constant, whereas the $g_2$ law would require
\[
 \frac b6=1+\frac{b}{6\tau}
\]
for the nonconstant function $\tau$, an impossibility.  This contradiction proves (i).

For (ii), truncate a fundamental polygon by the cusp neighborhood supplied by (i).  Its orbifold closure is compact, and the stated elliptic regularity makes $\tau$ and $\mu$ bounded there, while $\Im\tau$ has a positive lower bound.  At a point of $\mathcal D_{a,b}$,
\[
 a=\frac{6\Im\mu}{\Im\tau},
 \qquad b=6\mu-a\tau.
\]
Thus only finitely many integer pairs $(a,b)$ can occur on the compact truncation.  For each such pair, the holomorphic function $6\mu-a\tau-b$ is not identically zero by the preceding transformation-law argument, so it has only finitely many zeros on the compact orbifold chart system.  The union is finite.  The period monodromy merely permutes the marked pairs, so this is a statement on the quotient base.
\end{proof}

On a descent fibre, the degree-six holomorphic line bundle isolated in the supplied workbench has a unitary constant-curvature connection.  For the Hamiltonian $H_{\phi,\psi}=\tfrac12\nabla^*\nabla$ its spectrum is
\begin{equation}\label{eq:landau-spectrum}
 \Spec(H_{\phi,\psi})=
 \left\{
 \frac{12\pi}{\Im\tau_f}\left(j+\frac12\right):
 j\in\Z_{\ge0}
 \right\},
\end{equation}
with multiplicity six at each level.  The ground-shifted operator has levels
$12\pi j/\Im\tau_f$ and a six-dimensional kernel.

\begin{corollary}[The first-coordinate Landau sector does not run into the cusp]\label{cor:no-landau-cusp}
There is no sequence of actual descent fibres in the proposed $S^6$ family approaching the cusp.  In particular, the limit $\Im\tau_f\to\infty$ cannot be inferred from the $S^6$ cusp and cannot be used to close the Landau-level spacing in \eqref{eq:landau-spectrum} inside this family.
\end{corollary}

\begin{remark}[Correction of an ambient-family argument]
The pointwise matrices
\[
 (\tau_T,\mu_T,\beta_T)=(i/T,0,-iT)
\]
 do satisfy a nonzero period determinant and make certain flat or magnetic gaps small.  They are, however, an arbitrary path in the ambient coefficient space.  They are not shown to belong to the image of the asserted period map, and this specific first-coordinate magnetic construction also requires the independent arithmetic condition \eqref{eq:descent-condition}.  Theorem~\ref{thm:cusp-gap} supplies the correct scalar gap-closing statement on the actual cusp; Corollary~\ref{cor:no-landau-cusp} shows that this exceptional first-coordinate Landau sector does not realize it.  The independent full-torus construction in Theorem~\ref{thm:full-torus-magnetic-spectrum} does give a magnetic gap-closing sequence for every fixed nondegenerate integral flux on the marked cusp fibres, without asserting a monodromy-invariant global flux class or classifying holomorphic elliptic quotients.  The native full-flux theorem separately treats the explicit invariant class $u\wedge w+6\gamma\wedge\delta$ and does not use this exceptional elliptic quotient.

The descent theorem classifies only the specified first-coordinate map to
$\C/(\Z\tau+\Z)$.  It does not prove that $F_z$ has no other elliptic
quotient or elliptic subtorus, that its Néron--Severi group vanishes, or that
it admits no smooth integral magnetic line bundle; those are different typed
questions.
\end{remark}

\section{The compact total space is a different spectral problem}

\begin{proposition}[Compact total-space scalar gap]\label{prop:total-space}
Assume that the proposed completion is a smooth compact manifold $X\cong S^6$.  For every smooth Riemannian metric $g_X$ on $X$, the scalar Laplace--Beltrami operator has compact resolvent,
\[
 0=\lambda_0(X,g_X)<\lambda_1(X,g_X)\le\lambda_2(X,g_X)\le\cdots,
 \qquad \lambda_j\to\infty,
\]
and $\lambda_1(X,g_X)>0$ because $X$ is connected.
\end{proposition}

\begin{proof}
This is the standard elliptic spectral theorem on a connected closed Riemannian manifold \citep{Chavel_1984}.
\end{proof}

The proposition does not contradict Theorem~\ref{thm:cusp-gap}.  The latter concerns a singularly degenerating family of auxiliary vertical metrics on regular torus fibres.  No argument in the period construction shows that these Euclidean fibre metrics are restrictions of one smooth metric on the completed total space.  The same warning applies to the quotient metrics in Theorem~\ref{thm:global-real-quotient}: a pointed $\R^4$ limit of their regular-fibre targets is not a metric or operator on the compact total space.  A complex structure alone determines neither a preferred Riemannian metric nor a Hamiltonian.

\section{Metric indeterminacy and the physical boundary}

The preceding results attach gaplessness only to explicitly displayed scalar,
magnetic, classical Hessian, and finite Gaussian operators.  The following
example shows why the complex-geometric data alone cannot select any one of
those metric spectra.

\begin{proposition}[Fixed complex torus, fixed volume, arbitrarily small gap]
\label{prop:fixed-complex-metric-counterexample}
Let
\[
 A=(\C/(\Z+i\Z))^2,
 \qquad z_j=x_j+iy_j.
\]
For every $R\ge1$,
\begin{equation}\label{eq:fixed-complex-metric}
 g_R=R^2(\dd x_1^2+\dd y_1^2)
     +R^{-2}(\dd x_2^2+\dd y_2^2)
\end{equation}
is a flat Kähler metric of volume one on this one fixed complex torus, and
\begin{equation}\label{eq:fixed-complex-small-gap}
 \lambda_1(\Delta_{g_R})=\frac{4\pi^2}{R^2}\longrightarrow0.
\end{equation}
Therefore the complex structure, Hodge decomposition, integral cohomology,
and coherent direct-image data, even together with total volume, do not
determine a positive lower bound for the scalar Laplacian gap.
\end{proposition}

\begin{proof}
The two-form
\[
 \omega_R=R^2\dd x_1\wedge\dd y_1
          +R^{-2}\dd x_2\wedge\dd y_2
\]
is closed, positive, and invariant under the fixed complex structure, so
$g_R$ is Kähler.  Its real metric matrix is
$\operatorname{diag}(R^2,R^2,R^{-2},R^{-2})$, of determinant one, and the
coordinate fundamental domain has volume one.  The Fourier character
\[
 e_{a,b,c,d}=e^{2\pi i(ax_1+by_1+cx_2+dy_2)}
\]
has eigenvalue
\[
 4\pi^2\bigl(R^{-2}(a^2+b^2)+R^2(c^2+d^2)\bigr).
\]
For $R\ge1$ its minimum over nonzero integer quadruples is
$4\pi^2R^{-2}$, attained exactly at
$(\pm1,0,0,0)$ and $(0,\pm1,0,0)$.  This proves the claim; none of the
listed complex or cohomological data varies with $R$.
\end{proof}

The mathematical mass gap in the Clay formulation concerns the Hamiltonian
of a constructed four-dimensional quantum Yang--Mills theory, not an internal
scalar Laplacian or a classical Hessian \citep{JaffeWitten_2006}.  Euclidean
data reach such a Hamiltonian only after a positive-metric reconstruction
satisfying, in particular, the Osterwalder--Schrader requirements
\citep{OsterwalderSchrader_1973,OsterwalderSchrader_1975}.  The exact
operator-level implication that a transfer argument would need is the next
proposition.

\begin{proposition}[Typed Rayleigh-sequence transfer]
\label{prop:typed-physical-transfer}
Let $H\ge0$ be self-adjoint on a Hilbert space $\mathcal H_{\mathrm{phys}}$.
Put
\begin{equation}\label{eq:physical-excitation-space}
 \mathcal V_0=\ker H,\qquad P_0=\mathbf1_{\{0\}}(H),\qquad
 \mathcal H_{\mathrm{exc}}=(I-P_0)\mathcal H_{\mathrm{phys}},
\end{equation}
and let
\begin{equation}\label{eq:physical-energy-form}
 \mathfrak h[\psi]=\|H^{1/2}\psi\|^2,
 \qquad \mathcal Q(H)=\operatorname{Dom}(H^{1/2}).
\end{equation}
For each $n$, let $\mathcal K_n$ be a Hilbert space,
$\mathfrak q_n$ a nonnegative quadratic form on
$D(\mathfrak q_n)\subseteq\mathcal K_n$, and
$0\ne u_n\in D(\mathfrak q_n)$.  Suppose that a linear map
\[
 J_n:\operatorname{span}\{u_n\}\longrightarrow\mathcal Q(H)
\]
is given.  Set $\psi_n=(I-P_0)J_nu_n$.  If there are $c>0$, $C_n\ge0$,
and $\varepsilon_n\ge0$ such that
\begin{align}
 \|\psi_n\|&\ge c\|u_n\|,
 \label{eq:bridge-norm}\\
 \mathfrak h[\psi_n]
 &\le C_n\mathfrak q_n[u_n]+\varepsilon_n\|u_n\|^2,
 \label{eq:bridge-energy}\\
 C_n\frac{\mathfrak q_n[u_n]}{\|u_n\|^2}+\varepsilon_n
 &\longrightarrow0,
 \label{eq:bridge-smallness}
\end{align}
then the reducing restriction
\[
 H_{\mathrm{exc}}
 :=H\big|_{\operatorname{Dom}(H)\cap\mathcal H_{\mathrm{exc}}}
\]
is self-adjoint and
\begin{equation}\label{eq:bridge-conclusion}
 \inf\Spec(H_{\mathrm{exc}})=0,
 \qquad 0\in\sigma_{\mathrm{ess}}(H_{\mathrm{exc}}).
\end{equation}
Thus $H$ has no positive gap above its full zero-energy sector.
\end{proposition}

\begin{proof}
The spectral projection $P_0$ commutes with $H$ and $H^{1/2}$, so
$\mathcal H_{\mathrm{phys}}=\ker H\oplus\mathcal H_{\mathrm{exc}}$ is a
reducing decomposition.  Hence $H_{\mathrm{exc}}$ is self-adjoint,
nonnegative, has form domain
$\mathcal Q(H)\cap\mathcal H_{\mathrm{exc}}$, and has zero kernel.  Normalize
$w_n=\psi_n/\|\psi_n\|$.  Then
\[
 \mathfrak h[w_n]
 \le c^{-2}\left(
 C_n\frac{\mathfrak q_n[u_n]}{\|u_n\|^2}+\varepsilon_n\right)
 \longrightarrow0.
\]
The variational characterization of the bottom of a nonnegative
self-adjoint operator therefore gives
$\inf\Spec(H_{\mathrm{exc}})=0$.  Its spectrum is closed, so zero belongs to
it; its kernel is zero, so zero is not an eigenvalue.  A self-adjoint operator
has no residual spectrum.  Hence zero is a continuous spectral point and
belongs to the essential spectrum under the standard convention that removes
isolated eigenvalues of finite multiplicity.
\end{proof}

\begin{remark}[The missing physical morphism]
The native unscaled limit retains two compact directions, and the cusp
parameter $T$ has not been identified with a spacetime box size, lattice
spacing, correlation length, Euclidean time, or renormalization scale.  No
reflection-positive measure, transfer matrix, Osterwalder--Schrader
reconstruction, Gauss-law physical quotient, local observable algebra,
thermodynamic limit, or continuum limit in fixed physical units has been
constructed.  An eigenvalue of a four-dimensional elliptic Hessian has units
of inverse length squared; without a time reconstruction and a scale
conversion it is not an eigenvalue of a physical Hamiltonian.

The cusp calculations supply many possible auxiliary soft states: the scalar
characters in Theorem~\ref{thm:cusp-gap}, the continuum transverse modes in
Theorem~\ref{thm:continuum-YM-gap}, the finite transverse modes in
Theorem~\ref{thm:finite-YM-gap}, the native full-flux Landau bands and their
radial multiplication operator, and the finite Gaussian oscillator modes in
Theorem~\ref{cusp-link-gaussian-theorem}.  The reducible full-flux critical
connections additionally have negative Jacobi bands, so their Hessians are
not positive candidate vacuum Hamiltonians.  None of these constructions
supplies one already constructed, gauge-invariant, positive-metric physical
Hilbert space for four-dimensional quantum Yang--Mills theory, its Hamiltonian
$H$, or maps $J_n$ satisfying
\eqref{eq:bridge-norm}--\eqref{eq:bridge-smallness} with constants controlled
in physical units after gauge reduction, reconstruction, renormalization,
the infinite-volume limit, and the continuum limit.  Orthogonality to one
chosen vacuum vector would not suffice unless vacuum uniqueness were also
proved: the projection $I-P_0$ must remove the full zero-energy sector.
Consequently Proposition~\ref{prop:typed-physical-transfer} is an exact
transfer theorem, not an assertion that the required transfer data exist.
No nonlinear or quantum mass-gap conclusion follows here.
\end{remark}

\section{Conclusions}

Conditional on the manuscript's period, monodromy, and toric cusp formulae
in the precise scopes stated above, the following hierarchy is exact.

\begin{enumerate}[leftmargin=2.3em]
\item Every fixed regular flat fibre has a positive scalar spectral gap.
\item The family has no uniform flat-fibre gap at the actual cusp.  Its first
eigenvalue is
$4\pi^2T^{-2}(1+2\min\{b_{\mathrm I,0},0\}T^{-1}+O(T^{-2}))$, and after unit-volume
normalization it is
$4\pi^2T^{-1}(1+(-b_{\mathrm I,0}/2+2\min\{b_{\mathrm I,0},0\})T^{-1}+O(T^{-2}))$.
\item More than one eigenvalue becomes soft: if $0<E<4\pi^2$, the exact
fixed-energy count with Fourier multiplicity is
$N_T(E)=E\mathcal D_T/(4\pi)+O_E(T)$, where
$\mathcal D_T=|\det\mathsf P_z|=T^2-b_{\mathrm I,0}T+O(1)$.  Thus the low-energy
population has a two-dimensional leading density.  The unit-volume count is
$E\mathcal D_T^{1/2}/(4\pi)+O_E(\mathcal D_T^{1/4}+1)$, and at the
collapsing $T^{-2}$ scale the finite limiting count is the lattice-point count
in a two-dimensional Euclidean disc, with the stated threshold refinement.
\item The native pointed smooth limit of the marked four-tori is
$\T^2\times\R^2$, not $\R^4$.  The degree-$M(T)^2$ marked local fibre covers
have pointed $\R^4$ limit precisely when $M(T)\to\infty$ in the unscaled
metric, and precisely when $M(T)/\sqrt T\to\infty$ for the pullback of the
base-fibre unit-volume metric.  Normalizing every covering torus to its own
volume one never gives a pointed smooth $\R^4$ limit.
\item After $t_c=v^M$, the cover sublattice is preserved by the pulled-back
cusp monodromy.  Its exact Fourier spectrum separates the deck-invariant
pullback sector from the new momenta, and those momentum lattices become dense
in $\R^4$ when $T,M\to\infty$.  This punctured statement is distinct from
the completed local toric theorem below.
\item Up to sign, the unique primitive full-monodromy-invariant integral
character gives a real, non-holomorphic circle quotient of the regular group
family.  In this paper's convention its quotient circumference is exactly
$-TD/\sqrt{T^2+m^2}$ and is
$T-b_{\mathrm I,0}+O(T^{-1})$ on the positive radial marked lift.  The
character is global; for the chosen radial fibre metrics its
Fourier mode gap-closes, but it need not be the exact shortest mode at every
parameter.
\item Four synchronized copies, without a base change or fibre cover, give an
all-expanding real four-torus quotient with pointed $\R^4$ radial limit.
Four copies are minimal only within the stated fixed primitive
full-monodromy-invariant kernel class with its induced quotient metric.  The
different two-copy expanding quotient of the native fibre is cusp-local, and
no finite base cover globalizes its particular kernel.  Neither quotient is
the native fibre or its degree-$M^2$ cover.
\item The literal invariant-character map does not extend through the
asserted toric cusp with its fixed circle target.  Conditional on
$X\cong_{\mathrm{diff}}S^6$, its fourfold product has no extension to the
fixed target $(\R/\Z)^4$ from $X^4$ with the stated regular-fibre restriction,
and no such fixed-target extension factors through
$X^{\wedge4}\cong S^{24}$.  This is not an obstruction to a construction
that changes the codomain.
\item The exact local Kummer construction after $t=v^N$ requires the source
action
$\Psi^{(N)}_{\bar\lambda}(\xi,v)=
(\exp(2\pi iC(v^N)\bar\lambda/N)v^{B_0\bar\lambda}\xi,v)$.
It gives a finite flat degree-$N^3$ map to the original completed cusp
family.  The factor to the base-changed target has generic degree $N^2$ and,
for $N>1$, is quasi-etale but nonflat precisely over the double and triple
central-fibre strata.  These are local cusp statements, not a full-base
cover.
\item Separately, the canonical workbench constructs
$q_2^{\wedge2}$ from the sphere source $S^{12}$ and $q_1^{\wedge4}$ from
the sphere source $S^{24}$ by changing to compact cone targets.  This paper
does not reproduce their proofs.  Their radial target slices have intrinsic
pointed $\R^4$ limits but collapse in every fixed compatible ambient target
metric.
\item The associated mean-zero scalar vertical direct integral is injective
and gapless, with zero in its essential continuous spectrum.  This statement
identifies the threshold but does not assert the stronger full-half-line
spectral classification proved below for the native magnetic multiplication
operator.
\item The classical Yang--Mills Hessian at the product connection, after
Coulomb reduction and removal of harmonic one-forms, has exact gap
\eqref{eq:continuum-YM-gap-exact} and closes as
$4\pi^2/(g_{\mathrm{YM}}^2T^2)$.  The exact finite cubical cochain Hessian
closes uniformly as $T^{-2}$ for every $N\ge2$; resolving the varying
continuum geometry is the separate condition $N/T\to\infty$.
\item Keeping the covering degree $M_{\mathrm{cov}}^2$ separate from the
mesh count $N_{\mathrm{mesh}}^4$ gives an honest transported cubical complex,
typed continuum holonomy, and an explicit periodic compact-link
$\mathrm{SU}(2)$ flux.  Its logarithmic action is exactly
$2\pi^2M_{\mathrm{cov}}^2/(g_{\mathrm{YM}}^2D_R)$, and its displayed
$(1,2)$-sector Wilson functional converges to that value as
$N_{\mathrm{mesh}}\to\infty$.
\item On the reduced finite linear Coulomb space, the Gaussian Hamiltonian
has a unique ground state and exact gap
$\sqrt{\lambda_{\min}(g_{\mathrm{YM}}^{-2}d^*d)}$.  The diagonal sequence
$R_j=j^2$, $M_{\mathrm{cov},j}=j$,
$N_{\mathrm{mesh},j}=j^3$ simultaneously has pointed locally isometric
$\R^4$ cover limit, vanishing maximal physical edge length, vanishing
Gaussian gap, and vanishing displayed logarithmic action.
\item For every fixed nondegenerate integral alternating form $\mathsf E$ on
the marked four-torus, the full-torus constant-curvature magnetic Laplacian
has frequency product
$b_1b_2=4\pi^2|\operatorname{Pf}(\mathsf E)|/\mathcal D_T$ and level spacing
at most
$4\pi\sqrt{|\operatorname{Pf}(\mathsf E)|/\mathcal D_T}=O(T^{-1})$.
This is a fibrewise marked-cusp result, not an assertion of monodromy
invariance or a physical mass gap.
\item The distinguished class
$\eta=u\wedge w+6\gamma\wedge\delta$ is invariant under the asserted
monodromy generators, has elementary divisors $(1,6)$ and Pfaffian $6$, and
is represented by the explicit cocycle and constant-curvature connection
given below.  Its exact charge-one frequencies have product
$24\pi^2/D_R$ and asymptotics
$\mathcal B_{R,+}=12\pi/R+O(R^{-2})$ and
$\mathcal B_{R,-}=2\pi/R+O(R^{-2})$.  Every joint Landau level has complex
multiplicity six, the ground energy is $14\pi/R+O(R^{-2})$, and the first
distinct spacing is $4\pi/R+O(R^{-2})$.
\item After the fixed marking and volume-correcting unitary trivialize the
varying fibre Hilbert spaces, the full-flux magnetic family has common
$H^2$ operator domain and locally norm-continuous resolvents.  Its exact
maximal radial multiplication operators, in both fixed- and varying-fibre
forms over $(R_1,\infty)$ with Lebesgue measure, are self-adjoint and satisfy
\[
 \sigma=\sigma_{\mathrm{ess}}=\sigma_{\mathrm c}=[0,\infty),
 \qquad \sigma_{\mathrm p}=\varnothing,
 \qquad \ker=\{0\}.
\]
The parameter $R$ is not differentiated.
\item The nontrivial split bundle
$\mathscr L_{\eta,R}\oplus\mathscr L_{\eta,R}^{-1}$ carries a reducible
$\mathrm{SU}(2)$ Yang--Mills critical connection.  Its gauge-fixed Jacobi
form is closed, its operator domain is $H^2\Omega^1$, the maximal Coulomb
slice is reducing, and on a sufficiently deep cusp tail its complete
negative spectrum comprises three bands of real multiplicity $48$, giving
real Morse index $144$.  Thus this Hessian is not a positive fluctuation
operator.
\item More generally, every nonzero trace-zero integral weight vector
$\mathbf m=(m_1,\ldots,m_r)$ gives a split-Cartan $\mathrm{SU}(r)$ critical
connection with $k_{\mathbf m}=6\sum_jm_j^2$ and exact rootwise Jacobi
spectrum.  On the deep tail its real Morse index is
$36r\sum_jm_j^2=6rk_{\mathbf m}$; the neutral and odd-wall nullities are
given exactly in Theorem~\ref{cusp-native-full-flux-suN-nullity-theorem}.
These nonzero-$c_2$ bundles cannot be restrictions of bundles on a carrier
with vanishing $H^4(-;\mathbb Z)$.
\item The specific degree-six Landau construction pulled back from the
first-coordinate elliptic quotient is confined to a finite exceptional locus
and cannot approach the cusp.  This does not contradict the preceding
full-torus magnetic theorem, whose line bundle is defined directly from an
integral alternating form and is not pulled back from that elliptic quotient.
\item None of the preceding statements determines the scalar spectrum of a
smooth metric on the compact six-dimensional total space.  Even on one fixed
complex two-torus, unit-volume flat Kähler metrics can have arbitrarily small
scalar gap.
\item The finite Gaussian oscillator is not the compact-link Haar Hilbert
space, the exponential map from cochains is only a local chart, and the
magnetic Landau operators and reducible Jacobi operators remain fibrewise.
Transfer of any auxiliary soft sequence to a physical Hamiltonian would
require the explicit norm-noncollapse and energy-comparison morphisms in
Proposition~\ref{prop:typed-physical-transfer}.  No such physical Hilbert
space, Hamiltonian, or morphisms are constructed here, so no nonlinear or
quantum Yang--Mills mass-gap conclusion follows.
\end{enumerate}

The main structural lesson is that compactness and a positive gap at every
fibre do not survive integration over a noncompact degenerating modulus
without a uniform coercive estimate.  Conversely, a gapless decomposable
vertical operator or classical Hessian cannot be identified with a
field-theoretic Hamiltonian without the explicit operator-level bridge.
Likewise, the native $\T^2\times\R^2$ limit, finite-cover $\R^4$ unspooling,
fixed-target quotient limits, and cone-target constructions are four
different operations.  The fixed-target obstruction concerns only the
displayed character map to $(\R/\Z)^4$ with its stated fibre restriction;
the cone constructions change the codomain and distinguish intrinsic
expanding slice metrics from collapse in a fixed compact ambient metric.

\appendix
\section{Reproducibility formulas}

For a direct algebraic reproducibility check, expansion along the last two
columns of \eqref{eq:Z-P} gives
\[
 \det\mathsf P_z=T\left(p-\frac{6m^2}{T}\right)
\]
and multiplication of \eqref{eq:Z-P}$^{\mathsf T}$ by the vector in
\eqref{eq:complete-dual-vector} returns $(r,s,n_1,n_2)^{\mathsf T}$ exactly.
On the radial cusp, substitution of $p=B_R-T$ gives
\[
 -TD=T(T-B_R)+6m^2,
\]
while squaring the first two columns of the inverse transpose gives
\eqref{eq:unit-mode-exact-lengths} and their subtraction gives
\eqref{eq:unit-mode-comparison}.  These identities are exact.  The proof of
Theorem~\ref{thm:cusp-gap} then uses the uniform singular-value estimate
\eqref{eq:singular-value-bound}; no finite lattice search or numerical
sampling is part of the proof.

\bibliographystyle{apalike}
\bibliography{references}
\end{document}
